% \pragmaonce

% adapted from https://www.overleaf.com/learn/latex/Commands
\providecommand{\dissertationelse}[2]{%
% adapted from https://tex.stackexchange.com/a/33577
\ifdefined\DISSERTATION
#1
\else
#2
\fi
}


% One row per genome site.  The rows live in a generated data file rather than
% inline here: data/annotated-genome-16005-stint100-root1.tex is written by
% script/make_annotated_genome_table.py straight from the committed annotated
% genome JSON, alongside an RFC 4180 CSV of the same records for analysis.
% Instruction semantics contain TeX-special characters (%, &, ~, ^, _, |), so
% the generator escapes them on the way into the rows file; the CSV keeps them
% verbatim.
\newcommand{\genomerow}[9]{%
  #1 & #2 & #3 & \texttt{#4} & \texttt{#5} & \texttt{#6}
    & \texttt{#7} & \texttt{#8} & #9 \\
}

\newcommand{\genomehead}{%
\textbf{Site} & \textbf{Mod.} & \textbf{Off.} & \textbf{Operation}
  & \textbf{Semantics} & \textbf{a,b,c} & \textbf{State Target}
  & \textbf{Const.} & \textbf{Tag Moniker} \\
\midrule
}

\begin{landscape}
\begin{scriptsize}
% nine columns of default padding would push the table past the text block
\setlength{\tabcolsep}{2pt}
\begin{longtable}{%
  r@{\hspace{0.6em}}r@{\hspace{0.6em}}r
  >{\raggedright\arraybackslash}p{4.2cm}
  >{\raggedright\arraybackslash}p{6.15cm}
  c
  >{\raggedright\arraybackslash}p{4.6cm}
  r
  >{\raggedright\arraybackslash}p{2.8cm}
}
\caption{%
\textbf{Annotated genome of the round-100 co-evolved ``background'' strain isolate.}
\footnotesize
One row per instruction for the 2,357-instruction genome screened in Figure \ref{fig:ecogwas}, tabulated from the annotated JSON that DISHTINY emits for a specimen.
\emph{Site} indexes the linear instruction sequence, matching the knockout site labels in Figure \labelcref{fig:ecogwas:sites,fig:ecogwas-nobb-supp}.
\emph{Mod.} numbers the tag-addressable module a site falls in and \emph{Off.} gives its position within that module; modules are delimited by \texttt{Global Anchor} instructions, which register the tag by which a module can be triggered, so an anchor always sits at offset 0 of its own module.
\emph{Semantics}, \emph{State Target}, \emph{Const.} and \emph{Tag Moniker} reproduce the human-readable descriptors that the virtual CPU generates for each operation (Section \ref{sec:virtual_cpu}), and \emph{a,b,c} are the instruction's three raw register arguments.
A blank \emph{State Target} marks an operation that neither reads nor writes cell state, and a blank \emph{Const.} one that carries no immediate value.
Source data are committed alongside the manuscript: \texttt{data/genome-16005-stint100-root1.json} is the genome dump verbatim, and \texttt{script/make\_annotated\_genome\_table.py} flattens it to \texttt{data/annotated-genome-16005-stint100-root1.csv} --- which additionally records each site's enclosing module size, its full 64-bit tag, the jump-table index of control-flow operations, and the draw probability of \texttt{Generate Random Bool} --- and to the row file this table reads.
} \label{tab:annotated_genome} \\
\toprule
\genomehead
\endfirsthead
\caption*{\tablename{} \thetable{} (cont'd)} \\
\toprule
\genomehead
\endhead
\bottomrule
\endlastfoot
% Generated by script/make_annotated_genome_table.py from data/genome-16005-stint100-root1.json.
% Do not edit by hand.  One \genomerow per genome site; see
% table/annotated_genome.tex for the macro that typesets them.
\genomerow{0}{0}{0}{Global Anchor}{register a global jump-to destination, maybe terminate}{5,0,4}{}{}
\genomerow{1}{0}{1}{Terminate If}{if a, terminate core}{6,4,2}{}{}
\genomerow{2}{0}{2}{Terminal}{a = value}{2,2,3}{}{-0.953273}
\genomerow{3}{1}{0}{Global Anchor}{register a global jump-to destination, maybe terminate}{4,0,4}{}{}
\genomerow{4}{1}{1}{Set Global Regulator}{set global regulator value to a}{3,1,6}{}{}
\genomerow{5}{1}{2}{Subtract}{a = b - c}{2,4,7}{}{}
\genomerow{6}{1}{3}{Get Local Regulator}{a = local regulator value}{4,6,1}{}{}
\genomerow{7}{2}{0}{Global Anchor}{register a global jump-to destination, maybe terminate}{1,0,4}{}{}
\genomerow{8}{2}{1}{Terminate If}{if a, terminate core}{6,4,2}{}{}
\genomerow{9}{2}{2}{Draw Random Value}{a = weighted draw from random number generator}{2,2,3}{}{}
\genomerow{10}{3}{0}{Global Anchor}{register a global jump-to destination, maybe terminate}{4,0,4}{}{}
\genomerow{11}{3}{1}{Set Global Regulator}{set global regulator value to a}{3,1,6}{}{}
\genomerow{12}{3}{2}{Greater Than}{a = ( b > c )}{5,2,7}{}{}
\genomerow{13}{3}{3}{Generate Random Bool}{if p, a = 1; else a = 0}{3,2,6}{}{}
\genomerow{14}{3}{4}{Bitwise Shift}{a = b << c}{3,7,7}{}{}
\genomerow{15}{3}{5}{Terminal}{a = value}{7,4,4}{}{1.96058}
\genomerow{16}{3}{6}{Read Own State}{a = state}{0,6,5}{IsChildGroupOf[1]}{}
\genomerow{17}{3}{7}{Terminal}{a = value}{7,2,2}{}{0.650246}
\genomerow{18}{3}{8}{Local Anchor}{register a local jump-to destination}{6,5,4}{}{}
\genomerow{19}{4}{0}{Global Anchor}{register a global jump-to destination, maybe terminate}{5,0,4}{}{}
\genomerow{20}{4}{1}{Terminate If}{if a, terminate core}{6,4,2}{}{}
\genomerow{21}{4}{2}{Terminal}{a = value}{2,2,3}{}{-0.953273}
\genomerow{22}{4}{3}{Terminal}{a = value}{0,5,7}{}{-1.27139}
\genomerow{23}{4}{4}{Nop-0}{perform no operation}{3,5,4}{}{}
\genomerow{24}{4}{5}{Terminal}{a = value}{1,3,6}{}{0.283679}
\genomerow{25}{4}{6}{Write Own State}{if b, state = a}{6,7,5}{ResourceReceiveResistance[0]}{}
\genomerow{26}{4}{7}{Terminal}{a = value}{2,0,3}{}{-0.994727}
\genomerow{27}{4}{8}{Write Own State}{if b, state = a}{6,5,2}{TransientNopState[2]}{}
\genomerow{28}{5}{0}{Global Anchor}{register a global jump-to destination, maybe terminate}{0,2,4}{}{}
\genomerow{29}{5}{1}{Draw Random Value}{a = weighted draw from random number generator}{5,5,4}{}{}
\genomerow{30}{5}{2}{Nop-2}{perform no operation}{5,0,6}{}{}
\genomerow{31}{5}{3}{Send Intra-Cell Message}{if a, send message another cardinal within cell}{3,3,6}{}{}
\genomerow{32}{5}{4}{Logical Or}{a = b \textbar{}\textbar{} c}{0,3,3}{}{}
\genomerow{33}{5}{5}{Local Anchor}{register a local jump-to destination}{5,4,3}{}{}
\genomerow{34}{5}{6}{Divide}{a = b / c}{3,1,0}{}{}
\genomerow{35}{6}{0}{Global Anchor}{register a global jump-to destination, maybe terminate}{1,2,5}{}{}
\genomerow{36}{7}{0}{Global Anchor}{register a global jump-to destination, maybe terminate}{7,2,5}{}{}
\genomerow{37}{7}{1}{Decay Secondary Global Regulator}{decay global regulator value by register a}{5,2,7}{}{}
\genomerow{38}{7}{2}{Local Anchor}{register a local jump-to destination}{5,6,4}{}{}
\genomerow{39}{7}{3}{Read Neighbor State}{a = state}{7,4,4}{NumKnownQuorumBits[1]}{}
\genomerow{40}{7}{4}{Local Anchor}{register a local jump-to destination}{2,5,7}{}{}
\genomerow{41}{7}{5}{Get Global Regulator}{a = global regulator value}{0,5,2}{}{}
\genomerow{42}{7}{6}{Count Ones}{a = popcnt( b )}{4,0,3}{}{}
\genomerow{43}{7}{7}{Terminal}{a = value}{7,7,6}{}{872.396}
\genomerow{44}{7}{8}{Terminal}{a = value}{3,3,5}{}{0.412523}
\genomerow{45}{7}{9}{Count Ones}{a = popcnt( b )}{6,0,5}{}{}
\genomerow{46}{7}{10}{Local Anchor}{register a local jump-to destination}{4,2,0}{}{}
\genomerow{47}{7}{11}{Draw Random Value}{a = weighted draw from random number generator}{5,1,6}{}{}
\genomerow{48}{7}{12}{Less Than}{a = ( b < c )}{6,7,6}{}{}
\genomerow{49}{7}{13}{Terminal}{a = value}{3,2,7}{}{0.484778}
\genomerow{50}{7}{14}{Logical And}{a = b \&\& c}{4,2,0}{}{}
\genomerow{51}{7}{15}{Terminal}{a = value}{4,5,6}{}{-0.743729}
\genomerow{52}{7}{16}{Read Neighbor State}{a = state}{6,1,5}{ResourceSendRequest[0]}{}
\genomerow{53}{7}{17}{Read Own State}{a = state}{1,4,0}{MostRecentCauseOfDeath[0]}{}
\genomerow{54}{7}{18}{Terminal}{a = value}{6,0,2}{}{-0.205839}
\genomerow{55}{7}{19}{Adjust Secondary Global Regulator}{adjust global regulator value by a}{5,4,3}{}{}
\genomerow{56}{7}{20}{Terminal}{a = value}{4,4,7}{}{0.157256}
\genomerow{57}{7}{21}{Terminal}{a = value}{4,4,1}{}{-0.937871}
\genomerow{58}{7}{22}{Local Anchor}{register a local jump-to destination}{7,4,7}{}{}
\genomerow{59}{7}{23}{Not Equal}{a = (b != c)}{5,3,1}{}{}
\genomerow{60}{7}{24}{Divide}{a = b / c}{6,1,6}{}{}
\genomerow{61}{7}{25}{Terminal}{a = value}{4,3,0}{}{4.27122}
\genomerow{62}{7}{26}{Adjust Secondary Global Regulator}{adjust global regulator value by a}{0,7,5}{}{}
\genomerow{63}{7}{27}{Write Own State}{if b, state = a}{5,4,3}{RepLevRequest[1]}{}
\genomerow{64}{7}{28}{Not}{a = !a}{4,4,3}{}{}
\genomerow{65}{7}{29}{Send Intra-Cell Message}{if a, send message another cardinal within cell}{0,4,6}{}{}
\genomerow{66}{7}{30}{Global Jump If}{if a, goto tag match global anchor; if b, clear registers}{2,7,7}{}{}
\genomerow{67}{7}{31}{Terminal}{a = value}{1,5,0}{}{-13.5294}
\genomerow{68}{7}{32}{Decay Secondary Global Regulator}{decay global regulator value by register a}{3,1,6}{}{}
\genomerow{69}{7}{33}{Bitwise And}{a = b \& c}{4,4,3}{}{}
\genomerow{70}{7}{34}{Decay Global Regulator}{decay global regulator value by register a}{3,4,1}{}{}
\genomerow{71}{7}{35}{Terminal}{a = value}{4,0,7}{}{-0.184987}
\genomerow{72}{7}{36}{Terminal}{a = value}{4,3,0}{}{82.5566}
\genomerow{73}{7}{37}{Decay Global Regulator}{decay global regulator value by register a}{3,2,0}{}{}
\genomerow{74}{7}{38}{Write Own State}{if b, state = a}{1,7,3}{RepLevRequest[0]}{}
\genomerow{75}{7}{39}{Write Own State}{if b, state = a}{6,0,1}{RepLevRequest[1]}{}
\genomerow{76}{7}{40}{Write Own State}{if b, state = a}{0,4,6}{ResourceSendRequest[0]}{}
\genomerow{77}{7}{41}{Read Neighbor State}{a = state}{3,0,3}{KinGroupWillExpire[1]}{}
\genomerow{78}{7}{42}{Terminal}{a = value}{6,3,6}{}{-1.35619}
\genomerow{79}{7}{43}{Add}{a = b + c}{3,2,4}{}{}
\genomerow{80}{8}{0}{Global Anchor}{register a global jump-to destination, maybe terminate}{2,3,1}{}{}
\genomerow{81}{8}{1}{Modulo}{a = b \% c}{5,5,2}{}{}
\genomerow{82}{8}{2}{Write Own State}{if b, state = a}{1,5,1}{HeirRequest[0]}{}
\genomerow{83}{8}{3}{Terminal}{a = value}{4,2,0}{}{-0.242841}
\genomerow{84}{9}{0}{Global Anchor}{register a global jump-to destination, maybe terminate}{4,6,0}{}{}
\genomerow{85}{9}{1}{Modulo}{a = b \% c}{5,5,2}{}{}
\genomerow{86}{9}{2}{Write Own State}{if b, state = a}{1,5,1}{HeirRequest[0]}{}
\genomerow{87}{9}{3}{Bitwise Not}{a = \textasciitilde{}b}{4,2,0}{}{}
\genomerow{88}{10}{0}{Global Anchor}{register a global jump-to destination, maybe terminate}{4,6,0}{}{}
\genomerow{89}{10}{1}{Terminal}{a = value}{7,6,6}{}{-2.04925}
\genomerow{90}{10}{2}{Decay Local Regulator}{decay local regulator value by register a}{0,4,1}{}{}
\genomerow{91}{10}{3}{Add}{a = b + c}{3,6,4}{}{}
\genomerow{92}{10}{4}{Read Neighbor State}{a = state}{2,3,1}{TransientNopState[2]}{}
\genomerow{93}{10}{5}{Read Own State}{a = state}{3,3,1}{TransientNopState[0]}{}
\genomerow{94}{10}{6}{Read Neighbor State}{a = state}{6,6,7}{RepLevRequest[1]}{}
\genomerow{95}{10}{7}{Local Anchor}{register a local jump-to destination}{1,2,5}{}{}
\genomerow{96}{10}{8}{Send Inter-Cell Message}{if a, send message to neighbor cell}{6,7,3}{}{}
\genomerow{97}{10}{9}{Local Anchor}{register a local jump-to destination}{5,3,1}{}{}
\genomerow{98}{10}{10}{Local Anchor}{register a local jump-to destination}{2,2,4}{}{}
\genomerow{99}{10}{11}{Bitwise Not}{a = \textasciitilde{}b}{7,1,1}{}{}
\genomerow{100}{10}{12}{Adjust Global Regulator}{adjust global regulator value by a}{4,5,5}{}{}
\genomerow{101}{10}{13}{Set Global Regulator}{set global regulator value to a}{7,0,0}{}{}
\genomerow{102}{10}{14}{RandomFill}{fill a with uniformly-distributed random bits}{7,0,7}{}{}
\genomerow{103}{11}{0}{Global Anchor}{register a global jump-to destination, maybe terminate}{5,5,3}{}{}
\genomerow{104}{12}{0}{Global Anchor}{register a global jump-to destination, maybe terminate}{4,2,0}{}{}
\genomerow{105}{12}{1}{Add}{a = b + c}{4,6,2}{}{}
\genomerow{106}{12}{2}{Read Neighbor State}{a = state}{6,7,3}{TransientNopState[0]}{}
\genomerow{107}{13}{0}{Global Anchor}{register a global jump-to destination, maybe terminate}{4,2,0}{}{}
\genomerow{108}{13}{1}{Add}{a = b + c}{4,6,2}{}{}
\genomerow{109}{13}{2}{Read Neighbor State}{a = state}{6,7,3}{TransientNopState[0]}{}
\genomerow{110}{13}{3}{Bitwise And}{a = b \& c}{3,1,4}{}{}
\genomerow{111}{13}{4}{Read Neighbor State}{a = state}{4,6,6}{Epoch[0]}{}
\genomerow{112}{13}{5}{Not Equal}{a = (b != c)}{3,4,2}{}{}
\genomerow{113}{13}{6}{Logical Or}{a = b \textbar{}\textbar{} c}{6,1,4}{}{}
\genomerow{114}{13}{7}{Bitwise And}{a = b \& c}{3,1,4}{}{}
\genomerow{115}{13}{8}{Read Neighbor State}{a = state}{4,6,6}{Epoch[0]}{}
\genomerow{116}{13}{9}{Not Equal}{a = (b != c)}{3,4,2}{}{}
\genomerow{117}{13}{10}{Logical Or}{a = b \textbar{}\textbar{} c}{6,1,4}{}{}
\genomerow{118}{13}{11}{Terminal}{a = value}{1,2,7}{}{-4.43737}
\genomerow{119}{13}{12}{Local Anchor}{register a local jump-to destination}{4,2,2}{}{}
\genomerow{120}{13}{13}{Terminal}{a = value}{5,4,6}{}{392.349}
\genomerow{121}{13}{14}{Local Anchor}{register a local jump-to destination}{1,3,4}{}{}
\genomerow{122}{13}{15}{Increment}{a = a + 1}{4,6,1}{}{}
\genomerow{123}{13}{16}{Not Equal}{a = (b != c)}{0,1,0}{}{}
\genomerow{124}{14}{0}{Global Anchor}{register a global jump-to destination, maybe terminate}{6,7,0}{}{}
\genomerow{125}{15}{0}{Global Anchor}{register a global jump-to destination, maybe terminate}{1,3,1}{}{}
\genomerow{126}{15}{1}{Read Own State}{a = state}{5,6,2}{NeighborApoptosis[0]}{}
\genomerow{127}{15}{2}{Read Own State}{a = state}{1,3,4}{TransientNopState[3]}{}
\genomerow{128}{15}{3}{Terminal}{a = value}{4,6,1}{}{-0.898717}
\genomerow{129}{16}{0}{Global Anchor}{register a global jump-to destination, maybe terminate}{0,1,0}{}{}
\genomerow{130}{16}{1}{Local Jump If}{if a, goto tag match local anchor}{6,7,0}{}{}
\genomerow{131}{17}{0}{Global Anchor}{register a global jump-to destination, maybe terminate}{1,3,1}{}{}
\genomerow{132}{17}{1}{Bitwise Not}{a = \textasciitilde{}b}{7,5,3}{}{}
\genomerow{133}{17}{2}{Bitwise Shift}{a = b << c}{2,6,0}{}{}
\genomerow{134}{17}{3}{Send Inter-Cell Message}{if a, send message to neighbor cell}{1,3,3}{}{}
\genomerow{135}{18}{0}{Global Anchor}{register a global jump-to destination, maybe terminate}{4,4,1}{}{}
\genomerow{136}{18}{1}{Terminal}{a = value}{4,2,1}{}{0.0558399}
\genomerow{137}{18}{2}{Get Local Regulator}{a = local regulator value}{5,5,1}{}{}
\genomerow{138}{18}{3}{Subtract}{a = b - c}{5,4,2}{}{}
\genomerow{139}{18}{4}{Add}{a = b + c}{7,6,5}{}{}
\genomerow{140}{18}{5}{Terminal}{a = value}{6,7,0}{}{-50.511}
\genomerow{141}{19}{0}{Global Anchor}{register a global jump-to destination, maybe terminate}{1,3,1}{}{}
\genomerow{142}{19}{1}{Read Neighbor State}{a = state}{7,2,0}{ParentFragmented[0]}{}
\genomerow{143}{20}{0}{Global Anchor}{register a global jump-to destination, maybe terminate}{2,4,1}{}{}
\genomerow{144}{20}{1}{Logical Or}{a = b \textbar{}\textbar{} c}{1,2,6}{}{}
\genomerow{145}{20}{2}{Read Own State}{a = state}{5,0,5}{PhylogeneticRootView[0]}{}
\genomerow{146}{20}{3}{Terminal}{a = value}{1,3,4}{}{24.6398}
\genomerow{147}{20}{4}{Bitwise Not}{a = \textasciitilde{}b}{7,5,3}{}{}
\genomerow{148}{20}{5}{Bitwise Shift}{a = b << c}{2,6,0}{}{}
\genomerow{149}{20}{6}{Add To Own State}{state += a}{6,0,6}{TransientNopState[3]}{}
\genomerow{150}{20}{7}{Terminal}{a = value}{4,7,3}{}{-0.377056}
\genomerow{151}{21}{0}{Global Anchor}{register a global jump-to destination, maybe terminate}{0,1,0}{}{}
\genomerow{152}{22}{0}{Global Anchor}{register a global jump-to destination, maybe terminate}{6,7,0}{}{}
\genomerow{153}{23}{0}{Global Anchor}{register a global jump-to destination, maybe terminate}{5,6,1}{}{}
\genomerow{154}{23}{1}{Bitwise Not}{a = \textasciitilde{}b}{5,5,3}{}{}
\genomerow{155}{23}{2}{Terminal}{a = value}{1,2,5}{}{-19.6883}
\genomerow{156}{23}{3}{Greater Than}{a = ( b > c )}{4,1,4}{}{}
\genomerow{157}{23}{4}{Add To Own State}{state += a}{6,0,6}{ResourceSendRequest[0]}{}
\genomerow{158}{23}{5}{Terminal}{a = value}{4,7,3}{}{-0.377048}
\genomerow{159}{24}{0}{Global Anchor}{register a global jump-to destination, maybe terminate}{0,1,0}{}{}
\genomerow{160}{24}{1}{Terminal}{a = value}{4,7,3}{}{-0.377048}
\genomerow{161}{25}{0}{Global Anchor}{register a global jump-to destination, maybe terminate}{0,1,0}{}{}
\genomerow{162}{26}{0}{Global Anchor}{register a global jump-to destination, maybe terminate}{6,7,0}{}{}
\genomerow{163}{27}{0}{Global Anchor}{register a global jump-to destination, maybe terminate}{5,6,1}{}{}
\genomerow{164}{27}{1}{Bitwise Not}{a = \textasciitilde{}b}{1,5,3}{}{}
\genomerow{165}{27}{2}{Bitwise Shift}{a = b << c}{2,6,0}{}{}
\genomerow{166}{27}{3}{Send Inter-Cell Message}{if a, send message to neighbor cell}{0,1,0}{}{}
\genomerow{167}{27}{4}{Local Anchor}{register a local jump-to destination}{7,7,5}{}{}
\genomerow{168}{27}{5}{Multiply Own State}{state *= a}{5,6,3}{TransientNopState[2]}{}
\genomerow{169}{27}{6}{Terminal}{a = value}{5,6,3}{}{-0.332087}
\genomerow{170}{28}{0}{Global Anchor}{register a global jump-to destination, maybe terminate}{0,5,0}{}{}
\genomerow{171}{29}{0}{Global Anchor}{register a global jump-to destination, maybe terminate}{6,7,0}{}{}
\genomerow{172}{30}{0}{Global Anchor}{register a global jump-to destination, maybe terminate}{1,7,1}{}{}
\genomerow{173}{30}{1}{Bitwise Not}{a = \textasciitilde{}b}{5,5,3}{}{}
\genomerow{174}{30}{2}{Local Anchor}{register a local jump-to destination}{2,6,0}{}{}
\genomerow{175}{30}{3}{Write Own State}{if b, state = a}{0,1,1}{NopState[0]}{}
\genomerow{176}{30}{4}{Nop-2}{perform no operation}{4,6,6}{}{}
\genomerow{177}{30}{5}{Terminal}{a = value}{7,2,2}{}{1.08432}
\genomerow{178}{30}{6}{Count Ones}{a = popcnt( b )}{6,2,0}{}{}
\genomerow{179}{30}{7}{Local Anchor}{register a local jump-to destination}{3,4,6}{}{}
\genomerow{180}{30}{8}{Send Intra-Cell Message}{if a, send message another cardinal within cell}{2,2,4}{}{}
\genomerow{181}{31}{0}{Global Anchor}{register a global jump-to destination, maybe terminate}{2,5,2}{}{}
\genomerow{182}{31}{1}{Local Anchor}{register a local jump-to destination}{7,5,3}{}{}
\genomerow{183}{31}{2}{Get Local Regulator}{a = local regulator value}{3,7,5}{}{}
\genomerow{184}{31}{3}{Terminal}{a = value}{3,4,4}{}{-1.6443}
\genomerow{185}{32}{0}{Global Anchor}{register a global jump-to destination, maybe terminate}{5,4,0}{}{}
\genomerow{186}{32}{1}{Less Than}{a = ( b < c )}{1,5,7}{}{}
\genomerow{187}{32}{2}{Local Anchor}{register a local jump-to destination}{3,5,3}{}{}
\genomerow{188}{32}{3}{Local Jump If Not}{if !a, goto tag match local anchor}{2,6,1}{}{}
\genomerow{189}{32}{4}{Decay Local Regulator}{decay local regulator value by register a}{3,1,5}{}{}
\genomerow{190}{32}{5}{Local Anchor}{register a local jump-to destination}{6,7,7}{}{}
\genomerow{191}{32}{6}{Decay Local Regulator}{decay local regulator value by register a}{5,1,4}{}{}
\genomerow{192}{32}{7}{Bitwise Xor}{a = b \textasciicircum{} c}{4,0,1}{}{}
\genomerow{193}{32}{8}{Local Jump If Not}{if !a, goto tag match local anchor}{2,6,1}{}{}
\genomerow{194}{32}{9}{Local Anchor}{register a local jump-to destination}{6,7,7}{}{}
\genomerow{195}{32}{10}{Multiply Own State}{state *= a}{1,1,4}{RepLevRequest[1]}{}
\genomerow{196}{32}{11}{Divide}{a = b / c}{4,0,1}{}{}
\genomerow{197}{32}{12}{Local Anchor}{register a local jump-to destination}{1,4,6}{}{}
\genomerow{198}{32}{13}{Nop-0}{perform no operation}{4,2,4}{}{}
\genomerow{199}{32}{14}{Local Jump If}{if a, goto tag match local anchor}{3,7,6}{}{}
\genomerow{200}{32}{15}{Send Intra-Cell Message}{if a, send message another cardinal within cell}{4,5,5}{}{}
\genomerow{201}{33}{0}{Global Anchor}{register a global jump-to destination, maybe terminate}{3,3,6}{}{}
\genomerow{202}{33}{1}{Read Own State}{a = state}{4,5,4}{TransientNopState[0]}{}
\genomerow{203}{33}{2}{Terminal}{a = value}{4,3,4}{}{0.922447}
\genomerow{204}{33}{3}{Multiply Own State}{state *= a}{5,4,1}{NopState[3]}{}
\genomerow{205}{33}{4}{Read Own State}{a = state}{3,1,4}{IsNewborn[0]}{}
\genomerow{206}{33}{5}{Get Global Regulator}{a = global regulator value}{2,1,6}{}{}
\genomerow{207}{33}{6}{Add To Own State}{state += a}{5,2,5}{RepLevRequest[1]}{}
\genomerow{208}{33}{7}{Read Own State}{a = state}{7,0,6}{SpawnRequest[0]}{}
\genomerow{209}{33}{8}{Terminal}{a = value}{3,4,5}{}{-4.9108}
\genomerow{210}{33}{9}{Bitwise Not}{a = \textasciitilde{}b}{6,5,4}{}{}
\genomerow{211}{33}{10}{Local Anchor}{register a local jump-to destination}{6,4,1}{}{}
\genomerow{212}{33}{11}{Terminal}{a = value}{7,5,2}{}{0.286325}
\genomerow{213}{33}{12}{Send Intra-Cell Message}{if a, send message another cardinal within cell}{6,2,1}{}{}
\genomerow{214}{33}{13}{Multiply Own State}{state *= a}{0,0,0}{NopState[0]}{}
\genomerow{215}{33}{14}{Adjust Global Regulator}{adjust global regulator value by a}{5,7,0}{}{}
\genomerow{216}{33}{15}{Add To Own State}{state += a}{7,7,1}{ResourceReserveRequest[0]}{}
\genomerow{217}{33}{16}{Modulo}{a = b \% c}{1,5,7}{}{}
\genomerow{218}{33}{17}{Terminal}{a = value}{1,6,7}{}{0.560025}
\genomerow{219}{33}{18}{Read Own State}{a = state}{4,5,0}{ResourceSendRequest[0]}{}
\genomerow{220}{33}{19}{Adjust Global Regulator}{adjust global regulator value by a}{4,4,4}{}{}
\genomerow{221}{33}{20}{Read Own State}{a = state}{7,5,4}{KinGroupWillExpire[1]}{}
\genomerow{222}{33}{21}{Send Inter-Cell Message}{if a, send message to neighbor cell}{6,4,1}{}{}
\genomerow{223}{33}{22}{Set Local Regulator}{set local regulator value to a}{7,5,5}{}{}
\genomerow{224}{33}{23}{Broadcast Intra-Cell Message}{if a, send message to all other cardinals within cell}{4,0,1}{}{}
\genomerow{225}{33}{24}{Set Local Regulator}{set local regulator value to a}{5,5,5}{}{}
\genomerow{226}{33}{25}{Read Own State}{a = state}{4,5,0}{ResourceSendRequest[0]}{}
\genomerow{227}{33}{26}{Adjust Global Regulator}{adjust global regulator value by a}{4,4,4}{}{}
\genomerow{228}{33}{27}{Read Own State}{a = state}{7,5,4}{KinGroupWillExpire[1]}{}
\genomerow{229}{33}{28}{Send Inter-Cell Message}{if a, send message to neighbor cell}{6,4,1}{}{}
\genomerow{230}{33}{29}{Set Local Regulator}{set local regulator value to a}{7,5,5}{}{}
\genomerow{231}{33}{30}{Broadcast Intra-Cell Message}{if a, send message to all other cardinals within cell}{4,0,1}{}{}
\genomerow{232}{33}{31}{Set Local Regulator}{set local regulator value to a}{5,1,5}{}{}
\genomerow{233}{33}{32}{Local Anchor}{register a local jump-to destination}{4,0,1}{}{}
\genomerow{234}{33}{33}{Write Own State}{if b, state = a}{0,5,1}{RepLevRequest[0]}{}
\genomerow{235}{33}{34}{Logical And}{a = b \&\& c}{4,5,1}{}{}
\genomerow{236}{33}{35}{Local Anchor}{register a local jump-to destination}{4,2,7}{}{}
\genomerow{237}{33}{36}{Read Neighbor State}{a = state}{2,3,1}{KinGroupWillExpire[0]}{}
\genomerow{238}{34}{0}{Global Anchor}{register a global jump-to destination, maybe terminate}{7,5,1}{}{}
\genomerow{239}{35}{0}{Global Anchor}{register a global jump-to destination, maybe terminate}{3,7,2}{}{}
\genomerow{240}{35}{1}{Local Anchor}{register a local jump-to destination}{6,3,0}{}{}
\genomerow{241}{35}{2}{Terminal}{a = value}{4,6,0}{}{-28.3538}
\genomerow{242}{35}{3}{Read Own State}{a = state}{1,0,4}{KinGroupIDAncestorView[0]}{}
\genomerow{243}{35}{4}{Read Own State}{a = state}{5,2,7}{SpawnRequest[0]}{}
\genomerow{244}{35}{5}{Send Inter-Cell Message}{if a, send message to neighbor cell}{3,6,4}{}{}
\genomerow{245}{35}{6}{Count Ones}{a = popcnt( b )}{5,3,6}{}{}
\genomerow{246}{36}{0}{Global Anchor}{register a global jump-to destination, maybe terminate}{3,5,3}{}{}
\genomerow{247}{36}{1}{Terminal}{a = value}{5,3,2}{}{-1.65855}
\genomerow{248}{36}{2}{Local Anchor}{register a local jump-to destination}{4,1,1}{}{}
\genomerow{249}{36}{3}{Read Own State}{a = state}{1,0,4}{KinGroupIDAncestorView[0]}{}
\genomerow{250}{36}{4}{Read Own State}{a = state}{5,3,7}{SpawnRequest[0]}{}
\genomerow{251}{36}{5}{Send Inter-Cell Message}{if a, send message to neighbor cell}{3,6,4}{}{}
\genomerow{252}{36}{6}{Logical And}{a = b \&\& c}{6,5,2}{}{}
\genomerow{253}{36}{7}{Bitwise Not}{a = \textasciitilde{}b}{0,1,1}{}{}
\genomerow{254}{36}{8}{Send Intra-Cell Message}{if a, send message another cardinal within cell}{7,2,3}{}{}
\genomerow{255}{36}{9}{Read Neighbor State}{a = state}{1,2,1}{NopState[0]}{}
\genomerow{256}{37}{0}{Global Anchor}{register a global jump-to destination, maybe terminate}{7,3,2}{}{}
\genomerow{257}{37}{1}{Terminal}{a = value}{2,6,0}{}{2.77907}
\genomerow{258}{37}{2}{RandomFill}{fill a with uniformly-distributed random bits}{2,6,4}{}{}
\genomerow{259}{37}{3}{Count Ones}{a = popcnt( b )}{5,3,6}{}{}
\genomerow{260}{38}{0}{Global Anchor}{register a global jump-to destination, maybe terminate}{3,5,3}{}{}
\genomerow{261}{38}{1}{Terminal}{a = value}{5,3,2}{}{-1.65856}
\genomerow{262}{38}{2}{Local Anchor}{register a local jump-to destination}{4,1,5}{}{}
\genomerow{263}{38}{3}{Set Secondary Global Regulator}{set global regulator value to a}{2,0,0}{}{}
\genomerow{264}{38}{4}{Modulo}{a = b \% c}{3,1,5}{}{}
\genomerow{265}{38}{5}{Terminal}{a = value}{1,4,2}{}{9.29557}
\genomerow{266}{38}{6}{Local Anchor}{register a local jump-to destination}{3,5,3}{}{}
\genomerow{267}{38}{7}{Local Anchor}{register a local jump-to destination}{4,0,4}{}{}
\genomerow{268}{38}{8}{Generate Random Bool}{if p, a = 1; else a = 0}{3,3,1}{}{}
\genomerow{269}{38}{9}{Multiply Own State}{state *= a}{4,0,0}{HeirRequest[0]}{}
\genomerow{270}{38}{10}{Read Own State}{a = state}{3,0,1}{IsChildCellOf[0]}{}
\genomerow{271}{38}{11}{Negate}{a = -a}{7,1,2}{}{}
\genomerow{272}{38}{12}{Decay Secondary Global Regulator}{decay global regulator value by register a}{3,1,1}{}{}
\genomerow{273}{38}{13}{Terminal}{a = value}{4,6,2}{}{0.787094}
\genomerow{274}{38}{14}{Get Global Regulator}{a = global regulator value}{6,1,6}{}{}
\genomerow{275}{38}{15}{Local Jump If}{if a, goto tag match local anchor}{6,2,6}{}{}
\genomerow{276}{38}{16}{Read Neighbor State}{a = state}{7,5,2}{KinGroupIDAncestorView[1]}{}
\genomerow{277}{38}{17}{Local Jump If}{if a, goto tag match local anchor}{7,5,2}{}{}
\genomerow{278}{38}{18}{Nop-0}{perform no operation}{2,4,7}{}{}
\genomerow{279}{38}{19}{Logical Or}{a = b \textbar{}\textbar{} c}{2,0,3}{}{}
\genomerow{280}{38}{20}{Read Own State}{a = state}{3,6,4}{IsNewborn[0]}{}
\genomerow{281}{38}{21}{Local Anchor}{register a local jump-to destination}{2,3,0}{}{}
\genomerow{282}{38}{22}{Bitwise And}{a = b \& c}{2,7,7}{}{}
\genomerow{283}{38}{23}{Local Anchor}{register a local jump-to destination}{6,2,3}{}{}
\genomerow{284}{38}{24}{Terminal}{a = value}{0,3,1}{}{-16292.7}
\genomerow{285}{38}{25}{Read Own State}{a = state}{5,3,3}{NeighborKinGroupWillExpire[0]}{}
\genomerow{286}{38}{26}{Local Anchor}{register a local jump-to destination}{4,6,0}{}{}
\genomerow{287}{39}{0}{Global Anchor}{register a global jump-to destination, maybe terminate}{4,4,5}{}{}
\genomerow{288}{39}{1}{Send Inter-Cell Message}{if a, send message to neighbor cell}{3,6,4}{}{}
\genomerow{289}{39}{2}{Terminal}{a = value}{2,3,0}{}{-2.35291}
\genomerow{290}{39}{3}{Bitwise And}{a = b \& c}{2,5,7}{}{}
\genomerow{291}{40}{0}{Global Anchor}{register a global jump-to destination, maybe terminate}{4,2,4}{}{}
\genomerow{292}{40}{1}{Send Inter-Cell Message}{if a, send message to neighbor cell}{3,6,4}{}{}
\genomerow{293}{40}{2}{Local Anchor}{register a local jump-to destination}{2,3,0}{}{}
\genomerow{294}{40}{3}{Bitwise And}{a = b \& c}{2,5,7}{}{}
\genomerow{295}{40}{4}{Nop-2}{perform no operation}{4,2,4}{}{}
\genomerow{296}{40}{5}{Terminal}{a = value}{2,5,2}{}{-0.831769}
\genomerow{297}{40}{6}{Bitwise Shift}{a = b << c}{4,3,6}{}{}
\genomerow{298}{40}{7}{Negate}{a = -a}{2,2,7}{}{}
\genomerow{299}{40}{8}{Multiply}{a = b * c}{3,3,5}{}{}
\genomerow{300}{40}{9}{Local Anchor}{register a local jump-to destination}{4,2,6}{}{}
\genomerow{301}{40}{10}{Local Anchor}{register a local jump-to destination}{6,2,3}{}{}
\genomerow{302}{40}{11}{Send Intra-Cell Message}{if a, send message another cardinal within cell}{0,3,3}{}{}
\genomerow{303}{40}{12}{Logical And}{a = b \&\& c}{5,3,3}{}{}
\genomerow{304}{40}{13}{Local Anchor}{register a local jump-to destination}{5,6,0}{}{}
\genomerow{305}{41}{0}{Global Anchor}{register a global jump-to destination, maybe terminate}{0,4,5}{}{}
\genomerow{306}{41}{1}{Terminal}{a = value}{4,1,1}{}{-10.7109}
\genomerow{307}{41}{2}{Terminal}{a = value}{1,0,0}{}{0.173226}
\genomerow{308}{41}{3}{Write Own State}{if b, state = a}{7,1,6}{NopState[2]}{}
\genomerow{309}{41}{4}{Draw Random Value}{a = weighted draw from random number generator}{2,0,7}{}{}
\genomerow{310}{41}{5}{Terminal}{a = value}{6,7,5}{}{-0.356372}
\genomerow{311}{41}{6}{Bitwise And}{a = b \& c}{6,0,6}{}{}
\genomerow{312}{41}{7}{Nop-2}{perform no operation}{7,4,3}{}{}
\genomerow{313}{41}{8}{Terminal}{a = value}{6,6,5}{}{-0.352512}
\genomerow{314}{41}{9}{Terminal}{a = value}{4,6,5}{}{0.313987}
\genomerow{315}{42}{0}{Global Anchor}{register a global jump-to destination, maybe terminate}{4,7,0}{}{}
\genomerow{316}{42}{1}{Bitwise Xor}{a = b \textasciicircum{} c}{2,4,2}{}{}
\genomerow{317}{42}{2}{Terminal}{a = value}{2,3,0}{}{89.2276}
\genomerow{318}{42}{3}{Multiply Own State}{state *= a}{1,2,4}{ResourceReserveRequest[0]}{}
\genomerow{319}{43}{0}{Global Anchor}{register a global jump-to destination, maybe terminate}{6,7,2}{}{}
\genomerow{320}{43}{1}{Read Neighbor State}{a = state}{7,6,0}{NumKnownQuorumBits[0]}{}
\genomerow{321}{43}{2}{Bitwise Or}{a = b \textbar{} c}{2,3,2}{}{}
\genomerow{322}{43}{3}{Adjust Local Regulator}{adjust local regulator value by a}{1,5,4}{}{}
\genomerow{323}{43}{4}{Increment}{a = a + 1}{5,1,1}{}{}
\genomerow{324}{43}{5}{Send Intra-Cell Message}{if a, send message another cardinal within cell}{5,5,3}{}{}
\genomerow{325}{43}{6}{Write Own State}{if b, state = a}{5,5,1}{SpawnRequest[0]}{}
\genomerow{326}{43}{7}{Nop-1}{perform no operation}{2,7,3}{}{}
\genomerow{327}{43}{8}{Divide}{a = b / c}{2,3,0}{}{}
\genomerow{328}{43}{9}{Send Inter-Cell Message}{if a, send message to neighbor cell}{7,1,0}{}{}
\genomerow{329}{43}{10}{Send Intra-Cell Message}{if a, send message another cardinal within cell}{7,5,5}{}{}
\genomerow{330}{43}{11}{Increment}{a = a + 1}{5,7,5}{}{}
\genomerow{331}{43}{12}{Write Own State}{if b, state = a}{7,5,2}{SpawnArrest[0]}{}
\genomerow{332}{43}{13}{Equal}{a = ( b == c )}{4,0,2}{}{}
\genomerow{333}{43}{14}{Equal}{a = ( b == c )}{3,7,2}{}{}
\genomerow{334}{44}{0}{Global Anchor}{register a global jump-to destination, maybe terminate}{4,5,5}{}{}
\genomerow{335}{44}{1}{Send Intra-Cell Message}{if a, send message another cardinal within cell}{7,5,5}{}{}
\genomerow{336}{44}{2}{Increment}{a = a + 1}{5,7,5}{}{}
\genomerow{337}{44}{3}{Read Neighbor State}{a = state}{7,5,2}{IsChildGroupOf[1]}{}
\genomerow{338}{44}{4}{Equal}{a = ( b == c )}{4,0,2}{}{}
\genomerow{339}{44}{5}{Logical Or}{a = b \textbar{}\textbar{} c}{2,7,2}{}{}
\genomerow{340}{44}{6}{Bitwise Or}{a = b \textbar{} c}{0,7,0}{}{}
\genomerow{341}{44}{7}{Not Equal}{a = (b != c)}{1,3,0}{}{}
\genomerow{342}{44}{8}{Logical Or}{a = b \textbar{}\textbar{} c}{2,7,0}{}{}
\genomerow{343}{44}{9}{Fork If}{if a, submit fork request (processed when core terminates)}{2,4,0}{}{}
\genomerow{344}{45}{0}{Global Anchor}{register a global jump-to destination, maybe terminate}{1,3,2}{}{}
\genomerow{345}{45}{1}{Read Own State}{a = state}{4,0,1}{NopState[0]}{}
\genomerow{346}{45}{2}{Write Own State}{if b, state = a}{5,7,3}{TransientNopState[1]}{}
\genomerow{347}{45}{3}{Terminal}{a = value}{1,1,0}{}{156.557}
\genomerow{348}{45}{4}{Local Anchor}{register a local jump-to destination}{1,3,2}{}{}
\genomerow{349}{45}{5}{Write Own State}{if b, state = a}{5,5,1}{ResourceReserveRequest[0]}{}
\genomerow{350}{45}{6}{Generate Random Bool}{if p, a = 1; else a = 0}{1,1,0}{}{}
\genomerow{351}{45}{7}{Read Own State}{a = state}{5,7,1}{ApoptosisRequest[0]}{}
\genomerow{352}{45}{8}{Local Anchor}{register a local jump-to destination}{1,1,7}{}{}
\genomerow{353}{45}{9}{Local Jump If}{if a, goto tag match local anchor}{3,1,7}{}{}
\genomerow{354}{45}{10}{Read Neighbor State}{a = state}{1,4,2}{NumKnownQuorumBits[1]}{}
\genomerow{355}{45}{11}{Get Global Regulator}{a = global regulator value}{3,0,7}{}{}
\genomerow{356}{45}{12}{Draw Random Value}{a = weighted draw from random number generator}{7,7,5}{}{}
\genomerow{357}{45}{13}{Terminal}{a = value}{7,0,7}{}{0.343176}
\genomerow{358}{45}{14}{Local Anchor}{register a local jump-to destination}{0,7,6}{}{}
\genomerow{359}{45}{15}{Nop-0}{perform no operation}{0,4,7}{}{}
\genomerow{360}{45}{16}{Terminal}{a = value}{6,0,5}{}{8.64109}
\genomerow{361}{45}{17}{Nop-0}{perform no operation}{0,4,7}{}{}
\genomerow{362}{45}{18}{Bitwise Shift}{a = b << c}{5,2,7}{}{}
\genomerow{363}{45}{19}{Local Anchor}{register a local jump-to destination}{3,2,1}{}{}
\genomerow{364}{45}{20}{Read Own State}{a = state}{7,7,1}{IsNewborn[0]}{}
\genomerow{365}{45}{21}{Send Inter-Cell Message}{if a, send message to neighbor cell}{5,4,1}{}{}
\genomerow{366}{45}{22}{Set Global Regulator}{set global regulator value to a}{0,7,4}{}{}
\genomerow{367}{45}{23}{Bitwise Not}{a = \textasciitilde{}b}{7,4,7}{}{}
\genomerow{368}{45}{24}{Nop-1}{perform no operation}{3,7,3}{}{}
\genomerow{369}{45}{25}{Terminal}{a = value}{0,2,1}{}{-13.8505}
\genomerow{370}{45}{26}{Terminal}{a = value}{6,0,5}{}{8.6403}
\genomerow{371}{45}{27}{Set Secondary Global Regulator}{set global regulator value to a}{7,2,6}{}{}
\genomerow{372}{45}{28}{Local Anchor}{register a local jump-to destination}{4,2,3}{}{}
\genomerow{373}{45}{29}{Less Than}{a = ( b < c )}{5,6,7}{}{}
\genomerow{374}{46}{0}{Global Anchor}{register a global jump-to destination, maybe terminate}{4,5,0}{}{}
\genomerow{375}{46}{1}{Decay Global Regulator}{decay global regulator value by register a}{7,0,6}{}{}
\genomerow{376}{46}{2}{Local Anchor}{register a local jump-to destination}{4,7,6}{}{}
\genomerow{377}{46}{3}{Nop-0}{perform no operation}{0,0,7}{}{}
\genomerow{378}{46}{4}{Terminal}{a = value}{7,0,5}{}{-491791}
\genomerow{379}{46}{5}{Decay Secondary Global Regulator}{decay global regulator value by register a}{2,1,2}{}{}
\genomerow{380}{46}{6}{Bitwise And}{a = b \& c}{6,4,3}{}{}
\genomerow{381}{46}{7}{Nop-2}{perform no operation}{3,7,5}{}{}
\genomerow{382}{47}{0}{Global Anchor}{register a global jump-to destination, maybe terminate}{5,1,3}{}{}
\genomerow{383}{47}{1}{Bitwise Shift}{a = b << c}{2,3,7}{}{}
\genomerow{384}{47}{2}{Nop-1}{perform no operation}{0,5,3}{}{}
\genomerow{385}{47}{3}{Get Secondary Global Regulator}{a = global regulator value}{0,3,7}{}{}
\genomerow{386}{47}{4}{Bitwise Shift}{a = b << c}{2,3,7}{}{}
\genomerow{387}{47}{5}{Nop-1}{perform no operation}{0,5,3}{}{}
\genomerow{388}{47}{6}{Generate Random Bool}{if p, a = 1; else a = 0}{0,5,0}{}{}
\genomerow{389}{47}{7}{Global Jump If Not}{if !a, goto global anchor; if !b, clear registers}{0,5,2}{}{}
\genomerow{390}{47}{8}{Bitwise Shift}{a = b << c}{0,7,7}{}{}
\genomerow{391}{47}{9}{Bitwise Shift}{a = b << c}{6,3,7}{}{}
\genomerow{392}{47}{10}{Bitwise Shift}{a = b << c}{6,3,7}{}{}
\genomerow{393}{47}{11}{Set Local Regulator}{set local regulator value to a}{5,0,0}{}{}
\genomerow{394}{47}{12}{Bitwise Shift}{a = b << c}{2,3,7}{}{}
\genomerow{395}{47}{13}{Nop-1}{perform no operation}{0,4,3}{}{}
\genomerow{396}{47}{14}{Bitwise Shift}{a = b << c}{0,3,7}{}{}
\genomerow{397}{47}{15}{Get Secondary Global Regulator}{a = global regulator value}{2,3,7}{}{}
\genomerow{398}{47}{16}{Nop-1}{perform no operation}{0,5,3}{}{}
\genomerow{399}{47}{17}{Bitwise Shift}{a = b << c}{0,3,7}{}{}
\genomerow{400}{47}{18}{Nop-1}{perform no operation}{0,5,7}{}{}
\genomerow{401}{47}{19}{Draw Random Value}{a = weighted draw from random number generator}{0,5,1}{}{}
\genomerow{402}{47}{20}{Divide}{a = b / c}{0,5,3}{}{}
\genomerow{403}{47}{21}{Logical Or}{a = b \textbar{}\textbar{} c}{0,3,7}{}{}
\genomerow{404}{47}{22}{Read Own State}{a = state}{4,1,2}{IsChildCellOf[0]}{}
\genomerow{405}{47}{23}{Divide}{a = b / c}{0,5,3}{}{}
\genomerow{406}{47}{24}{Local Anchor}{register a local jump-to destination}{5,7,3}{}{}
\genomerow{407}{47}{25}{Subtract}{a = b - c}{1,6,1}{}{}
\genomerow{408}{47}{26}{Read Own State}{a = state}{1,0,4}{IsChildCellOf[0]}{}
\genomerow{409}{47}{27}{Read Neighbor State}{a = state}{6,5,7}{NeighborFragmented[0]}{}
\genomerow{410}{47}{28}{Write Own State}{if b, state = a}{4,7,7}{SpawnRequest[0]}{}
\genomerow{411}{47}{29}{Increment}{a = a + 1}{0,5,3}{}{}
\genomerow{412}{47}{30}{Send Intra-Cell Message}{if a, send message another cardinal within cell}{0,5,7}{}{}
\genomerow{413}{47}{31}{Write Own State}{if b, state = a}{4,7,1}{SpawnArrest[0]}{}
\genomerow{414}{47}{32}{Divide}{a = b / c}{3,7,3}{}{}
\genomerow{415}{47}{33}{Logical Or}{a = b \textbar{}\textbar{} c}{0,7,3}{}{}
\genomerow{416}{47}{34}{Local Anchor}{register a local jump-to destination}{0,5,3}{}{}
\genomerow{417}{47}{35}{Get Local Regulator}{a = local regulator value}{5,5,1}{}{}
\genomerow{418}{47}{36}{Generate Random Bool}{if p, a = 1; else a = 0}{6,0,2}{}{}
\genomerow{419}{47}{37}{Not}{a = !a}{3,5,7}{}{}
\genomerow{420}{47}{38}{Read Own State}{a = state}{6,5,4}{PhylogeneticRootView[0]}{}
\genomerow{421}{47}{39}{Set Local Regulator}{set local regulator value to a}{2,4,4}{}{}
\genomerow{422}{47}{40}{Adjust Local Regulator}{adjust local regulator value by a}{5,7,7}{}{}
\genomerow{423}{47}{41}{Bitwise Shift}{a = b << c}{1,6,6}{}{}
\genomerow{424}{48}{0}{Global Anchor}{register a global jump-to destination, maybe terminate}{0,0,2}{}{}
\genomerow{425}{48}{1}{Divide}{a = b / c}{5,5,4}{}{}
\genomerow{426}{48}{2}{Read Neighbor State}{a = state}{2,5,2}{KinGroupIDAncestorView[0]}{}
\genomerow{427}{48}{3}{Read Neighbor State}{a = state}{3,6,2}{CellAge[0]}{}
\genomerow{428}{48}{4}{Local Anchor}{register a local jump-to destination}{5,2,2}{}{}
\genomerow{429}{48}{5}{Read Own State}{a = state}{6,7,6}{NeighborOptimumQuorumExceeded[1]}{}
\genomerow{430}{48}{6}{Set Local Regulator}{set local regulator value to a}{2,5,4}{}{}
\genomerow{431}{48}{7}{Terminal}{a = value}{3,2,3}{}{0.224082}
\genomerow{432}{48}{8}{Adjust Global Regulator}{adjust global regulator value by a}{7,6,2}{}{}
\genomerow{433}{48}{9}{Bitwise And}{a = b \& c}{1,1,4}{}{}
\genomerow{434}{48}{10}{RandomFill}{fill a with uniformly-distributed random bits}{5,4,0}{}{}
\genomerow{435}{48}{11}{Get Local Regulator}{a = local regulator value}{7,2,7}{}{}
\genomerow{436}{48}{12}{Local Anchor}{register a local jump-to destination}{7,0,3}{}{}
\genomerow{437}{48}{13}{Logical And}{a = b \&\& c}{0,0,2}{}{}
\genomerow{438}{48}{14}{Bitwise Shift}{a = b << c}{4,4,5}{}{}
\genomerow{439}{48}{15}{Terminal}{a = value}{3,6,6}{}{0.864839}
\genomerow{440}{48}{16}{Local Anchor}{register a local jump-to destination}{1,3,3}{}{}
\genomerow{441}{48}{17}{Logical Or}{a = b \textbar{}\textbar{} c}{0,4,7}{}{}
\genomerow{442}{48}{18}{Nop-2}{perform no operation}{5,6,7}{}{}
\genomerow{443}{49}{0}{Global Anchor}{register a global jump-to destination, maybe terminate}{6,2,5}{}{}
\genomerow{444}{49}{1}{Decay Local Regulator}{decay local regulator value by register a}{0,6,1}{}{}
\genomerow{445}{49}{2}{Terminal}{a = value}{4,4,0}{}{-0.303057}
\genomerow{446}{49}{3}{Local Anchor}{register a local jump-to destination}{3,7,6}{}{}
\genomerow{447}{49}{4}{Bitwise Shift}{a = b << c}{7,4,0}{}{}
\genomerow{448}{50}{0}{Global Anchor}{register a global jump-to destination, maybe terminate}{0,5,6}{}{}
\genomerow{449}{50}{1}{Bitwise Shift}{a = b << c}{7,0,0}{}{}
\genomerow{450}{50}{2}{Read Neighbor State}{a = state}{6,1,7}{Epoch[0]}{}
\genomerow{451}{50}{3}{Read Own State}{a = state}{6,0,2}{IsParentGroupOf[1]}{}
\genomerow{452}{50}{4}{Local Anchor}{register a local jump-to destination}{6,6,7}{}{}
\genomerow{453}{50}{5}{Local Anchor}{register a local jump-to destination}{4,5,5}{}{}
\genomerow{454}{50}{6}{Write Own State}{if b, state = a}{1,5,0}{TransientNopState[3]}{}
\genomerow{455}{50}{7}{Local Anchor}{register a local jump-to destination}{6,6,7}{}{}
\genomerow{456}{50}{8}{Local Anchor}{register a local jump-to destination}{4,5,5}{}{}
\genomerow{457}{50}{9}{Write Own State}{if b, state = a}{1,5,0}{TransientNopState[1]}{}
\genomerow{458}{51}{0}{Global Anchor}{register a global jump-to destination, maybe terminate}{2,0,2}{}{}
\genomerow{459}{51}{1}{Local Anchor}{register a local jump-to destination}{6,4,3}{}{}
\genomerow{460}{51}{2}{Terminal}{a = value}{0,2,5}{}{-0.0348913}
\genomerow{461}{52}{0}{Global Anchor}{register a global jump-to destination, maybe terminate}{5,1,1}{}{}
\genomerow{462}{52}{1}{Read Own State}{a = state}{4,6,0}{RicherThanNeighbor[0]}{}
\genomerow{463}{53}{0}{Global Anchor}{register a global jump-to destination, maybe terminate}{1,1,1}{}{}
\genomerow{464}{53}{1}{Set Global Regulator}{set global regulator value to a}{7,2,1}{}{}
\genomerow{465}{53}{2}{Less Than}{a = ( b < c )}{6,3,0}{}{}
\genomerow{466}{53}{3}{Local Anchor}{register a local jump-to destination}{1,4,6}{}{}
\genomerow{467}{53}{4}{Local Jump If}{if a, goto tag match local anchor}{4,2,2}{}{}
\genomerow{468}{53}{5}{Local Anchor}{register a local jump-to destination}{0,4,6}{}{}
\genomerow{469}{53}{6}{Terminal}{a = value}{1,0,2}{}{-10.5958}
\genomerow{470}{53}{7}{Set Local Regulator}{set local regulator value to a}{5,4,1}{}{}
\genomerow{471}{54}{0}{Global Anchor}{register a global jump-to destination, maybe terminate}{1,1,2}{}{}
\genomerow{472}{54}{1}{Less Than}{a = ( b < c )}{7,2,1}{}{}
\genomerow{473}{54}{2}{Get Local Regulator}{a = local regulator value}{1,0,6}{}{}
\genomerow{474}{54}{3}{Add To Own State}{state += a}{5,4,1}{TransientNopState[3]}{}
\genomerow{475}{54}{4}{Less Than}{a = ( b < c )}{5,0,1}{}{}
\genomerow{476}{54}{5}{Less Than}{a = ( b < c )}{3,0,1}{}{}
\genomerow{477}{54}{6}{Logical And}{a = b \&\& c}{3,5,7}{}{}
\genomerow{478}{54}{7}{Bitwise Not}{a = \textasciitilde{}b}{2,5,4}{}{}
\genomerow{479}{54}{8}{Terminal}{a = value}{7,6,4}{}{-0.736573}
\genomerow{480}{55}{0}{Global Anchor}{register a global jump-to destination, maybe terminate}{2,6,5}{}{}
\genomerow{481}{55}{1}{Add To Own State}{state += a}{3,6,0}{RepLevRequest[1]}{}
\genomerow{482}{56}{0}{Global Anchor}{register a global jump-to destination, maybe terminate}{0,7,3}{}{}
\genomerow{483}{56}{1}{Decay Secondary Global Regulator}{decay global regulator value by register a}{0,1,6}{}{}
\genomerow{484}{56}{2}{Fork If}{if a, submit fork request (processed when core terminates)}{4,2,6}{}{}
\genomerow{485}{57}{0}{Global Anchor}{register a global jump-to destination, maybe terminate}{2,6,5}{}{}
\genomerow{486}{57}{1}{Read Own State}{a = state}{3,4,6}{KinGroupIDAncestorView[1]}{}
\genomerow{487}{57}{2}{Get Local Regulator}{a = local regulator value}{4,7,5}{}{}
\genomerow{488}{57}{3}{Multiply}{a = b * c}{3,4,4}{}{}
\genomerow{489}{57}{4}{Nop-0}{perform no operation}{4,1,4}{}{}
\genomerow{490}{57}{5}{Not Equal}{a = (b != c)}{6,7,4}{}{}
\genomerow{491}{57}{6}{Not Equal}{a = (b != c)}{6,7,4}{}{}
\genomerow{492}{57}{7}{Get Global Regulator}{a = global regulator value}{5,5,5}{}{}
\genomerow{493}{58}{0}{Global Anchor}{register a global jump-to destination, maybe terminate}{0,7,3}{}{}
\genomerow{494}{58}{1}{Get Secondary Global Regulator}{a = global regulator value}{0,1,6}{}{}
\genomerow{495}{58}{2}{Read Own State}{a = state}{3,6,6}{KinGroupIDAncestorView[1]}{}
\genomerow{496}{58}{3}{Get Local Regulator}{a = local regulator value}{4,7,5}{}{}
\genomerow{497}{58}{4}{Multiply}{a = b * c}{3,5,4}{}{}
\genomerow{498}{58}{5}{Nop-0}{perform no operation}{4,1,4}{}{}
\genomerow{499}{58}{6}{Not Equal}{a = (b != c)}{6,7,4}{}{}
\genomerow{500}{58}{7}{Get Global Regulator}{a = global regulator value}{5,5,5}{}{}
\genomerow{501}{58}{8}{Read Neighbor State}{a = state}{1,5,7}{KinGroupIDAncestorView[0]}{}
\genomerow{502}{58}{9}{Read Own State}{a = state}{7,4,0}{NeighborIsNewborn[0]}{}
\genomerow{503}{58}{10}{Terminal}{a = value}{6,5,4}{}{0.896754}
\genomerow{504}{58}{11}{Read Own State}{a = state}{7,5,3}{NeighborKinGroupWillExpire[0]}{}
\genomerow{505}{58}{12}{Bitwise Not}{a = \textasciitilde{}b}{2,2,3}{}{}
\genomerow{506}{58}{13}{Local Anchor}{register a local jump-to destination}{2,1,1}{}{}
\genomerow{507}{58}{14}{Not Equal}{a = (b != c)}{6,2,4}{}{}
\genomerow{508}{58}{15}{Local Anchor}{register a local jump-to destination}{2,1,1}{}{}
\genomerow{509}{58}{16}{Not Equal}{a = (b != c)}{6,2,4}{}{}
\genomerow{510}{58}{17}{Send Inter-Cell Message}{if a, send message to neighbor cell}{7,4,7}{}{}
\genomerow{511}{58}{18}{Read Own State}{a = state}{3,2,2}{TransientNopState[2]}{}
\genomerow{512}{58}{19}{Read Neighbor State}{a = state}{6,2,4}{StockpileFecund[0]}{}
\genomerow{513}{58}{20}{Less Than}{a = ( b < c )}{2,7,5}{}{}
\genomerow{514}{58}{21}{Draw Random Value}{a = weighted draw from random number generator}{3,7,1}{}{}
\genomerow{515}{59}{0}{Global Anchor}{register a global jump-to destination, maybe terminate}{5,1,6}{}{}
\genomerow{516}{60}{0}{Global Anchor}{register a global jump-to destination, maybe terminate}{2,1,5}{}{}
\genomerow{517}{61}{0}{Global Anchor}{register a global jump-to destination, maybe terminate}{6,7,1}{}{}
\genomerow{518}{61}{1}{Decrement}{a = a - 1}{6,4,0}{}{}
\genomerow{519}{61}{2}{Greater Than}{a = ( b > c )}{0,3,2}{}{}
\genomerow{520}{61}{3}{Send Inter-Cell Message}{if a, send message to neighbor cell}{7,7,7}{}{}
\genomerow{521}{61}{4}{Read Own State}{a = state}{6,2,6}{KinGroupIDView[0]}{}
\genomerow{522}{61}{5}{Read Own State}{a = state}{7,4,0}{KinGroupIDAncestorView[1]}{}
\genomerow{523}{61}{6}{Adjust Global Regulator}{adjust global regulator value by a}{6,2,5}{}{}
\genomerow{524}{61}{7}{Add To Own State}{state += a}{1,0,4}{RepLevRequest[1]}{}
\genomerow{525}{62}{0}{Global Anchor}{register a global jump-to destination, maybe terminate}{2,5,3}{}{}
\genomerow{526}{62}{1}{Read Neighbor State}{a = state}{6,2,4}{StockpileFecund[0]}{}
\genomerow{527}{62}{2}{Send Inter-Cell Message}{if a, send message to neighbor cell}{7,7,7}{}{}
\genomerow{528}{62}{3}{Read Own State}{a = state}{6,2,6}{Epoch[0]}{}
\genomerow{529}{62}{4}{Read Own State}{a = state}{7,4,0}{KinGroupIDAncestorView[1]}{}
\genomerow{530}{62}{5}{Local Anchor}{register a local jump-to destination}{6,2,5}{}{}
\genomerow{531}{62}{6}{Add To Own State}{state += a}{3,0,4}{TransientNopState[3]}{}
\genomerow{532}{62}{7}{Logical And}{a = b \&\& c}{2,5,3}{}{}
\genomerow{533}{62}{8}{Generate Random Bool}{if p, a = 1; else a = 0}{4,7,4}{}{}
\genomerow{534}{62}{9}{Read Own State}{a = state}{0,7,7}{ApoptosisRequest[0]}{}
\genomerow{535}{62}{10}{Generate Random Bool}{if p, a = 1; else a = 0}{2,0,1}{}{}
\genomerow{536}{62}{11}{Negate}{a = -a}{7,4,7}{}{}
\genomerow{537}{62}{12}{Write Own State}{if b, state = a}{1,4,7}{ResourceReserveRequest[0]}{}
\genomerow{538}{62}{13}{Adjust Global Regulator}{adjust global regulator value by a}{6,6,1}{}{}
\genomerow{539}{63}{0}{Global Anchor}{register a global jump-to destination, maybe terminate}{0,7,7}{}{}
\genomerow{540}{63}{1}{Generate Random Bool}{if p, a = 1; else a = 0}{2,0,1}{}{}
\genomerow{541}{63}{2}{Negate}{a = -a}{7,4,5}{}{}
\genomerow{542}{63}{3}{Logical And}{a = b \&\& c}{0,7,7}{}{}
\genomerow{543}{63}{4}{Generate Random Bool}{if p, a = 1; else a = 0}{2,0,1}{}{}
\genomerow{544}{63}{5}{Negate}{a = -a}{7,4,5}{}{}
\genomerow{545}{63}{6}{Write Own State}{if b, state = a}{1,4,7}{SpawnRequest[0]}{}
\genomerow{546}{64}{0}{Global Anchor}{register a global jump-to destination, maybe terminate}{6,6,1}{}{}
\genomerow{547}{64}{1}{Read Own State}{a = state}{0,7,7}{ApoptosisRequest[0]}{}
\genomerow{548}{64}{2}{Read Own State}{a = state}{2,0,1}{KinGroupMatch[0]}{}
\genomerow{549}{64}{3}{Negate}{a = -a}{7,4,5}{}{}
\genomerow{550}{64}{4}{Write Own State}{if b, state = a}{1,5,7}{ResourceReserveRequest[0]}{}
\genomerow{551}{64}{5}{Set Global Regulator}{set global regulator value to a}{4,6,7}{}{}
\genomerow{552}{64}{6}{Add To Own State}{state += a}{6,4,2}{HeirRequest[0]}{}
\genomerow{553}{64}{7}{Get Global Regulator}{a = global regulator value}{1,5,0}{}{}
\genomerow{554}{64}{8}{Local Anchor}{register a local jump-to destination}{5,3,1}{}{}
\genomerow{555}{65}{0}{Global Anchor}{register a global jump-to destination, maybe terminate}{5,5,2}{}{}
\genomerow{556}{66}{0}{Global Anchor}{register a global jump-to destination, maybe terminate}{1,6,0}{}{}
\genomerow{557}{66}{1}{Local Anchor}{register a local jump-to destination}{1,3,1}{}{}
\genomerow{558}{67}{0}{Global Anchor}{register a global jump-to destination, maybe terminate}{2,7,1}{}{}
\genomerow{559}{67}{1}{Terminal}{a = value}{2,2,4}{}{3.10007}
\genomerow{560}{67}{2}{Local Jump If Not}{if !a, goto tag match local anchor}{0,0,5}{}{}
\genomerow{561}{67}{3}{Local Anchor}{register a local jump-to destination}{1,5,1}{}{}
\genomerow{562}{67}{4}{Bitwise Xor}{a = b \textasciicircum{} c}{5,4,4}{}{}
\genomerow{563}{68}{0}{Global Anchor}{register a global jump-to destination, maybe terminate}{5,7,6}{}{}
\genomerow{564}{68}{1}{Less Than}{a = ( b < c )}{3,1,6}{}{}
\genomerow{565}{68}{2}{Equal}{a = ( b == c )}{0,0,6}{}{}
\genomerow{566}{68}{3}{Local Anchor}{register a local jump-to destination}{6,4,6}{}{}
\genomerow{567}{68}{4}{Local Anchor}{register a local jump-to destination}{0,1,2}{}{}
\genomerow{568}{68}{5}{Read Own State}{a = state}{6,2,6}{Epoch[0]}{}
\genomerow{569}{68}{6}{Multiply Own State}{state *= a}{5,1,7}{SpawnArrest[0]}{}
\genomerow{570}{69}{0}{Global Anchor}{register a global jump-to destination, maybe terminate}{1,1,3}{}{}
\genomerow{571}{69}{1}{Terminal}{a = value}{3,5,6}{}{-89.6354}
\genomerow{572}{69}{2}{Read Own State}{a = state}{2,7,2}{NopState[3]}{}
\genomerow{573}{69}{3}{Local Anchor}{register a local jump-to destination}{2,6,0}{}{}
\genomerow{574}{69}{4}{Less Than}{a = ( b < c )}{4,4,5}{}{}
\genomerow{575}{69}{5}{Send Intra-Cell Message}{if a, send message another cardinal within cell}{4,3,0}{}{}
\genomerow{576}{69}{6}{Subtract}{a = b - c}{0,6,7}{}{}
\genomerow{577}{69}{7}{Send Intra-Cell Message}{if a, send message another cardinal within cell}{6,3,0}{}{}
\genomerow{578}{69}{8}{Subtract}{a = b - c}{0,6,7}{}{}
\genomerow{579}{69}{9}{Read Own State}{a = state}{2,7,2}{ResourceSendRequest[0]}{}
\genomerow{580}{69}{10}{Send Intra-Cell Message}{if a, send message another cardinal within cell}{4,3,0}{}{}
\genomerow{581}{69}{11}{Multiply}{a = b * c}{0,6,7}{}{}
\genomerow{582}{69}{12}{Read Own State}{a = state}{2,7,2}{NumKnownQuorumBits[0]}{}
\genomerow{583}{70}{0}{Global Anchor}{register a global jump-to destination, maybe terminate}{4,4,5}{}{}
\genomerow{584}{71}{0}{Global Anchor}{register a global jump-to destination, maybe terminate}{5,3,0}{}{}
\genomerow{585}{71}{1}{Read Own State}{a = state}{6,7,2}{SpawnRequest[0]}{}
\genomerow{586}{71}{2}{Decay Secondary Global Regulator}{decay global regulator value by register a}{4,6,4}{}{}
\genomerow{587}{71}{3}{Decay Secondary Global Regulator}{decay global regulator value by register a}{2,6,0}{}{}
\genomerow{588}{71}{4}{Bitwise Not}{a = \textasciitilde{}b}{2,3,0}{}{}
\genomerow{589}{72}{0}{Global Anchor}{register a global jump-to destination, maybe terminate}{2,6,0}{}{}
\genomerow{590}{72}{1}{Modulo}{a = b \% c}{2,3,0}{}{}
\genomerow{591}{72}{2}{Send Intra-Cell Message}{if a, send message another cardinal within cell}{1,7,0}{}{}
\genomerow{592}{72}{3}{Local Anchor}{register a local jump-to destination}{0,5,2}{}{}
\genomerow{593}{73}{0}{Global Anchor}{register a global jump-to destination, maybe terminate}{2,6,0}{}{}
\genomerow{594}{73}{1}{Modulo}{a = b \% c}{2,3,0}{}{}
\genomerow{595}{73}{2}{Send Intra-Cell Message}{if a, send message another cardinal within cell}{1,7,0}{}{}
\genomerow{596}{73}{3}{Local Anchor}{register a local jump-to destination}{0,5,2}{}{}
\genomerow{597}{73}{4}{Decay Local Regulator}{decay local regulator value by register a}{7,6,0}{}{}
\genomerow{598}{73}{5}{Terminal}{a = value}{4,7,0}{}{5.98985}
\genomerow{599}{74}{0}{Global Anchor}{register a global jump-to destination, maybe terminate}{5,1,6}{}{}
\genomerow{600}{74}{1}{Get Local Regulator}{a = local regulator value}{0,0,5}{}{}
\genomerow{601}{74}{2}{Set Secondary Global Regulator}{set global regulator value to a}{0,5,6}{}{}
\genomerow{602}{74}{3}{Get Local Regulator}{a = local regulator value}{7,2,6}{}{}
\genomerow{603}{74}{4}{Increment}{a = a + 1}{1,7,0}{}{}
\genomerow{604}{74}{5}{Bitwise Xor}{a = b \textasciicircum{} c}{0,5,2}{}{}
\genomerow{605}{74}{6}{Terminal}{a = value}{4,7,0}{}{4.33849}
\genomerow{606}{75}{0}{Global Anchor}{register a global jump-to destination, maybe terminate}{5,1,6}{}{}
\genomerow{607}{75}{1}{Get Local Regulator}{a = local regulator value}{0,0,5}{}{}
\genomerow{608}{75}{2}{Set Secondary Global Regulator}{set global regulator value to a}{0,5,6}{}{}
\genomerow{609}{75}{3}{Get Local Regulator}{a = local regulator value}{7,3,6}{}{}
\genomerow{610}{75}{4}{Increment}{a = a + 1}{1,7,0}{}{}
\genomerow{611}{75}{5}{Bitwise Or}{a = b \textbar{} c}{0,5,2}{}{}
\genomerow{612}{76}{0}{Global Anchor}{register a global jump-to destination, maybe terminate}{5,6,0}{}{}
\genomerow{613}{77}{0}{Global Anchor}{register a global jump-to destination, maybe terminate}{7,7,4}{}{}
\genomerow{614}{77}{1}{Send Inter-Cell Message}{if a, send message to neighbor cell}{3,0,0}{}{}
\genomerow{615}{77}{2}{Decay Local Regulator}{decay local regulator value by register a}{4,5,7}{}{}
\genomerow{616}{77}{3}{Terminal}{a = value}{3,7,3}{}{-0.23512}
\genomerow{617}{77}{4}{Send Intra-Cell Message}{if a, send message another cardinal within cell}{5,2,4}{}{}
\genomerow{618}{77}{5}{Draw Random Value}{a = weighted draw from random number generator}{3,2,3}{}{}
\genomerow{619}{77}{6}{Read Own State}{a = state}{7,6,0}{IsChildGroupOf[0]}{}
\genomerow{620}{77}{7}{Terminal}{a = value}{5,7,0}{}{5.98985}
\genomerow{621}{78}{0}{Global Anchor}{register a global jump-to destination, maybe terminate}{5,1,6}{}{}
\genomerow{622}{78}{1}{Get Local Regulator}{a = local regulator value}{0,0,5}{}{}
\genomerow{623}{79}{0}{Global Anchor}{register a global jump-to destination, maybe terminate}{6,3,6}{}{}
\genomerow{624}{79}{1}{Logical Or}{a = b \textbar{}\textbar{} c}{3,0,6}{}{}
\genomerow{625}{79}{2}{Terminal}{a = value}{3,3,2}{}{-2.75462}
\genomerow{626}{79}{3}{Set Secondary Global Regulator}{set global regulator value to a}{0,5,6}{}{}
\genomerow{627}{79}{4}{Terminal}{a = value}{7,0,7}{}{-0.588531}
\genomerow{628}{79}{5}{Terminal}{a = value}{7,4,0}{}{-0.168869}
\genomerow{629}{79}{6}{Terminal}{a = value}{1,3,6}{}{-0.0379}
\genomerow{630}{79}{7}{Local Jump If}{if a, goto tag match local anchor}{2,7,6}{}{}
\genomerow{631}{79}{8}{Subtract}{a = b - c}{2,4,7}{}{}
\genomerow{632}{79}{9}{Get Local Regulator}{a = local regulator value}{4,4,1}{}{}
\genomerow{633}{79}{10}{Get Local Regulator}{a = local regulator value}{2,3,4}{}{}
\genomerow{634}{79}{11}{Get Global Regulator}{a = global regulator value}{4,1,1}{}{}
\genomerow{635}{80}{0}{Global Anchor}{register a global jump-to destination, maybe terminate}{2,5,2}{}{}
\genomerow{636}{81}{0}{Global Anchor}{register a global jump-to destination, maybe terminate}{4,7,2}{}{}
\genomerow{637}{82}{0}{Global Anchor}{register a global jump-to destination, maybe terminate}{2,5,2}{}{}
\genomerow{638}{82}{1}{Decay Local Regulator}{decay local regulator value by register a}{3,6,3}{}{}
\genomerow{639}{82}{2}{Adjust Global Regulator}{adjust global regulator value by a}{6,1,4}{}{}
\genomerow{640}{83}{0}{Global Anchor}{register a global jump-to destination, maybe terminate}{6,6,2}{}{}
\genomerow{641}{83}{1}{Read Own State}{a = state}{7,2,6}{IsNewborn[0]}{}
\genomerow{642}{84}{0}{Global Anchor}{register a global jump-to destination, maybe terminate}{5,3,7}{}{}
\genomerow{643}{84}{1}{Read Neighbor State}{a = state}{0,3,5}{HeirRequest[0]}{}
\genomerow{644}{84}{2}{Greater Than}{a = ( b > c )}{1,4,6}{}{}
\genomerow{645}{84}{3}{Send Inter-Cell Message}{if a, send message to neighbor cell}{0,1,0}{}{}
\genomerow{646}{84}{4}{Bitwise Or}{a = b \textbar{} c}{4,5,1}{}{}
\genomerow{647}{84}{5}{Local Anchor}{register a local jump-to destination}{1,3,5}{}{}
\genomerow{648}{84}{6}{Less Than}{a = ( b < c )}{6,1,3}{}{}
\genomerow{649}{84}{7}{Write Own State}{if b, state = a}{1,1,0}{RepLevRequest[1]}{}
\genomerow{650}{85}{0}{Global Anchor}{register a global jump-to destination, maybe terminate}{4,6,4}{}{}
\genomerow{651}{85}{1}{Terminal}{a = value}{6,6,7}{}{-0.0702862}
\genomerow{652}{85}{2}{Terminal}{a = value}{0,4,5}{}{-0.028735}
\genomerow{653}{86}{0}{Global Anchor}{register a global jump-to destination, maybe terminate}{3,6,7}{}{}
\genomerow{654}{86}{1}{Add To Own State}{state += a}{3,4,5}{TransientNopState[3]}{}
\genomerow{655}{87}{0}{Global Anchor}{register a global jump-to destination, maybe terminate}{5,0,6}{}{}
\genomerow{656}{87}{1}{Add To Own State}{state += a}{0,4,6}{NopState[0]}{}
\genomerow{657}{87}{2}{Read Neighbor State}{a = state}{5,5,1}{NopState[2]}{}
\genomerow{658}{87}{3}{Terminal}{a = value}{3,0,1}{}{0.722374}
\genomerow{659}{87}{4}{Bitwise Not}{a = \textasciitilde{}b}{3,0,1}{}{}
\genomerow{660}{88}{0}{Global Anchor}{register a global jump-to destination, maybe terminate}{5,5,0}{}{}
\genomerow{661}{88}{1}{Less Than}{a = ( b < c )}{6,6,7}{}{}
\genomerow{662}{88}{2}{Less Than}{a = ( b < c )}{3,4,2}{}{}
\genomerow{663}{88}{3}{Not Equal}{a = (b != c)}{6,6,3}{}{}
\genomerow{664}{89}{0}{Global Anchor}{register a global jump-to destination, maybe terminate}{0,0,6}{}{}
\genomerow{665}{89}{1}{Count Ones}{a = popcnt( b )}{1,4,2}{}{}
\genomerow{666}{89}{2}{Terminate If}{if a, terminate core}{5,2,6}{}{}
\genomerow{667}{89}{3}{Write Own State}{if b, state = a}{5,2,7}{TransientNopState[2]}{}
\genomerow{668}{89}{4}{Multiply}{a = b * c}{1,1,3}{}{}
\genomerow{669}{89}{5}{Generate Random Bool}{if p, a = 1; else a = 0}{7,1,0}{}{}
\genomerow{670}{89}{6}{Read Own State}{a = state}{0,6,6}{IsParentGroupOf[1]}{}
\genomerow{671}{89}{7}{Local Anchor}{register a local jump-to destination}{6,6,0}{}{}
\genomerow{672}{89}{8}{Add To Own State}{state += a}{2,2,6}{TransientNopState[2]}{}
\genomerow{673}{89}{9}{Multiply Own State}{state *= a}{1,6,5}{SpawnArrest[0]}{}
\genomerow{674}{89}{10}{Get Local Regulator}{a = local regulator value}{6,5,3}{}{}
\genomerow{675}{89}{11}{Local Anchor}{register a local jump-to destination}{2,0,5}{}{}
\genomerow{676}{90}{0}{Global Anchor}{register a global jump-to destination, maybe terminate}{0,7,2}{}{}
\genomerow{677}{90}{1}{Draw Random Value}{a = weighted draw from random number generator}{7,4,7}{}{}
\genomerow{678}{90}{2}{Generate Random Bool}{if p, a = 1; else a = 0}{0,4,5}{}{}
\genomerow{679}{90}{3}{Get Global Regulator}{a = global regulator value}{4,4,5}{}{}
\genomerow{680}{90}{4}{Negate}{a = -a}{0,7,2}{}{}
\genomerow{681}{90}{5}{Local Jump If Not}{if !a, goto tag match local anchor}{3,0,1}{}{}
\genomerow{682}{90}{6}{Not}{a = !a}{6,6,6}{}{}
\genomerow{683}{90}{7}{Get Local Regulator}{a = local regulator value}{0,6,7}{}{}
\genomerow{684}{90}{8}{Negate}{a = -a}{0,7,2}{}{}
\genomerow{685}{90}{9}{Local Jump If Not}{if !a, goto tag match local anchor}{3,0,1}{}{}
\genomerow{686}{90}{10}{Draw Random Value}{a = weighted draw from random number generator}{2,6,6}{}{}
\genomerow{687}{90}{11}{Local Anchor}{register a local jump-to destination}{3,5,1}{}{}
\genomerow{688}{90}{12}{Terminal}{a = value}{1,1,4}{}{27.1161}
\genomerow{689}{90}{13}{Nop-0}{perform no operation}{5,6,2}{}{}
\genomerow{690}{90}{14}{Terminal}{a = value}{1,3,3}{}{-1.13946}
\genomerow{691}{90}{15}{Read Neighbor State}{a = state}{5,1,2}{TransientNopState[2]}{}
\genomerow{692}{91}{0}{Global Anchor}{register a global jump-to destination, maybe terminate}{4,5,2}{}{}
\genomerow{693}{91}{1}{Terminal}{a = value}{4,5,7}{}{-4.38582}
\genomerow{694}{91}{2}{Terminal}{a = value}{5,0,2}{}{6.11917}
\genomerow{695}{91}{3}{Read Own State}{a = state}{5,6,3}{KinGroupWillExpire[1]}{}
\genomerow{696}{91}{4}{Local Anchor}{register a local jump-to destination}{4,4,3}{}{}
\genomerow{697}{91}{5}{Terminal}{a = value}{3,1,5}{}{67.2374}
\genomerow{698}{91}{6}{Send Intra-Cell Message}{if a, send message another cardinal within cell}{0,0,3}{}{}
\genomerow{699}{91}{7}{Terminal}{a = value}{4,4,2}{}{19.6949}
\genomerow{700}{91}{8}{Terminal}{a = value}{6,1,3}{}{-0.495556}
\genomerow{701}{91}{9}{Draw Random Value}{a = weighted draw from random number generator}{2,6,6}{}{}
\genomerow{702}{91}{10}{Terminal}{a = value}{0,6,7}{}{-0.924392}
\genomerow{703}{91}{11}{Less Than}{a = ( b < c )}{2,0,5}{}{}
\genomerow{704}{91}{12}{Local Jump If}{if a, goto tag match local anchor}{7,6,2}{}{}
\genomerow{705}{91}{13}{Local Jump If Not}{if !a, goto tag match local anchor}{3,0,1}{}{}
\genomerow{706}{91}{14}{Generate Random Bool}{if p, a = 1; else a = 0}{7,3,7}{}{}
\genomerow{707}{91}{15}{Local Anchor}{register a local jump-to destination}{4,4,6}{}{}
\genomerow{708}{91}{16}{Terminal}{a = value}{6,5,2}{}{-0.0810032}
\genomerow{709}{91}{17}{Nop-0}{perform no operation}{5,6,2}{}{}
\genomerow{710}{91}{18}{Less Than}{a = ( b < c )}{0,0,5}{}{}
\genomerow{711}{91}{19}{Local Jump If}{if a, goto tag match local anchor}{7,6,2}{}{}
\genomerow{712}{91}{20}{Logical Or}{a = b \textbar{}\textbar{} c}{3,5,2}{}{}
\genomerow{713}{91}{21}{Bitwise Xor}{a = b \textasciicircum{} c}{7,3,1}{}{}
\genomerow{714}{91}{22}{Equal}{a = ( b == c )}{3,1,6}{}{}
\genomerow{715}{91}{23}{RandomFill}{fill a with uniformly-distributed random bits}{0,1,6}{}{}
\genomerow{716}{91}{24}{Nop-2}{perform no operation}{5,4,0}{}{}
\genomerow{717}{91}{25}{Local Anchor}{register a local jump-to destination}{6,5,1}{}{}
\genomerow{718}{91}{26}{Terminal}{a = value}{6,5,1}{}{8.62912}
\genomerow{719}{91}{27}{Send Intra-Cell Message}{if a, send message another cardinal within cell}{5,2,7}{}{}
\genomerow{720}{91}{28}{Terminal}{a = value}{3,1,2}{}{291358}
\genomerow{721}{91}{29}{Local Anchor}{register a local jump-to destination}{6,5,1}{}{}
\genomerow{722}{92}{0}{Global Anchor}{register a global jump-to destination, maybe terminate}{5,5,4}{}{}
\genomerow{723}{92}{1}{Bitwise Not}{a = \textasciitilde{}b}{1,5,5}{}{}
\genomerow{724}{92}{2}{Global Jump If}{if a, goto tag match global anchor; if b, clear registers}{0,5,7}{}{}
\genomerow{725}{92}{3}{Terminal}{a = value}{6,6,4}{}{-24.9342}
\genomerow{726}{92}{4}{Generate Random Bool}{if p, a = 1; else a = 0}{7,4,2}{}{}
\genomerow{727}{92}{5}{Subtract}{a = b - c}{4,2,5}{}{}
\genomerow{728}{92}{6}{Decay Secondary Global Regulator}{decay global regulator value by register a}{2,2,4}{}{}
\genomerow{729}{92}{7}{Logical Or}{a = b \textbar{}\textbar{} c}{2,1,7}{}{}
\genomerow{730}{92}{8}{Nop-1}{perform no operation}{0,2,6}{}{}
\genomerow{731}{93}{0}{Global Anchor}{register a global jump-to destination, maybe terminate}{4,3,0}{}{}
\genomerow{732}{93}{1}{Read Own State}{a = state}{3,4,5}{IsAlive[0]}{}
\genomerow{733}{93}{2}{Local Anchor}{register a local jump-to destination}{6,5,2}{}{}
\genomerow{734}{93}{3}{Global Jump If Not}{if !a, goto global anchor; if !b, clear registers}{4,2,7}{}{}
\genomerow{735}{93}{4}{Read Own State}{a = state}{1,3,0}{NeighborFragmented[0]}{}
\genomerow{736}{93}{5}{Broadcast Intra-Cell Message}{if a, send message to all other cardinals within cell}{1,6,5}{}{}
\genomerow{737}{93}{6}{Get Global Regulator}{a = global regulator value}{7,4,3}{}{}
\genomerow{738}{93}{7}{Read Neighbor State}{a = state}{4,2,2}{KinGroupIDAncestorView[0]}{}
\genomerow{739}{93}{8}{Read Own State}{a = state}{0,3,4}{ReceivedResourceFrom[0]}{}
\genomerow{740}{93}{9}{Count Ones}{a = popcnt( b )}{3,0,2}{}{}
\genomerow{741}{93}{10}{Read Own State}{a = state}{1,3,0}{MostRecentCauseOfDeath[0]}{}
\genomerow{742}{93}{11}{Local Anchor}{register a local jump-to destination}{3,6,6}{}{}
\genomerow{743}{93}{12}{Decay Local Regulator}{decay local regulator value by register a}{7,2,3}{}{}
\genomerow{744}{93}{13}{Multiply Own State}{state *= a}{0,4,6}{NopState[1]}{}
\genomerow{745}{93}{14}{Logical And}{a = b \&\& c}{5,5,2}{}{}
\genomerow{746}{93}{15}{Send Inter-Cell Message}{if a, send message to neighbor cell}{7,0,7}{}{}
\genomerow{747}{94}{0}{Global Anchor}{register a global jump-to destination, maybe terminate}{3,4,2}{}{}
\genomerow{748}{94}{1}{Send Intra-Cell Message}{if a, send message another cardinal within cell}{3,0,0}{}{}
\genomerow{749}{94}{2}{Write Own State}{if b, state = a}{6,4,4}{ApoptosisRequest[0]}{}
\genomerow{750}{94}{3}{Multiply Own State}{state *= a}{5,6,6}{RepLevRequest[1]}{}
\genomerow{751}{94}{4}{Terminal}{a = value}{4,6,1}{}{-13.4992}
\genomerow{752}{94}{5}{Modulo}{a = b \% c}{7,2,1}{}{}
\genomerow{753}{94}{6}{Read Neighbor State}{a = state}{5,4,0}{KinGroupWillExpire[1]}{}
\genomerow{754}{94}{7}{Terminal}{a = value}{6,1,4}{}{-0.0529842}
\genomerow{755}{94}{8}{Read Neighbor State}{a = state}{5,4,0}{SpawnArrest[0]}{}
\genomerow{756}{94}{9}{Negate}{a = -a}{4,1,5}{}{}
\genomerow{757}{94}{10}{Terminal}{a = value}{4,6,1}{}{-1.28429}
\genomerow{758}{94}{11}{Read Neighbor State}{a = state}{5,4,0}{SpawnArrest[0]}{}
\genomerow{759}{94}{12}{Negate}{a = -a}{4,1,5}{}{}
\genomerow{760}{94}{13}{Terminal}{a = value}{4,6,1}{}{-1.28429}
\genomerow{761}{94}{14}{Modulo}{a = b \% c}{7,2,1}{}{}
\genomerow{762}{94}{15}{Modulo}{a = b \% c}{7,2,1}{}{}
\genomerow{763}{94}{16}{Send Intra-Cell Message}{if a, send message another cardinal within cell}{5,4,0}{}{}
\genomerow{764}{94}{17}{Local Anchor}{register a local jump-to destination}{5,0,5}{}{}
\genomerow{765}{94}{18}{Write Own State}{if b, state = a}{0,5,5}{SpawnArrest[0]}{}
\genomerow{766}{94}{19}{Local Anchor}{register a local jump-to destination}{1,5,2}{}{}
\genomerow{767}{94}{20}{Terminal}{a = value}{6,1,4}{}{-0.0540199}
\genomerow{768}{94}{21}{Read Neighbor State}{a = state}{5,4,0}{RicherThanNeighbor[0]}{}
\genomerow{769}{94}{22}{Negate}{a = -a}{5,1,5}{}{}
\genomerow{770}{94}{23}{Greater Than}{a = ( b > c )}{0,2,2}{}{}
\genomerow{771}{94}{24}{Set Local Regulator}{set local regulator value to a}{4,3,4}{}{}
\genomerow{772}{94}{25}{Decay Local Regulator}{decay local regulator value by register a}{7,6,3}{}{}
\genomerow{773}{94}{26}{Write Own State}{if b, state = a}{7,6,5}{RepLevRequest[0]}{}
\genomerow{774}{94}{27}{Increment}{a = a + 1}{0,6,7}{}{}
\genomerow{775}{94}{28}{Logical And}{a = b \&\& c}{7,4,1}{}{}
\genomerow{776}{94}{29}{Terminal}{a = value}{6,1,4}{}{0.949885}
\genomerow{777}{94}{30}{Local Anchor}{register a local jump-to destination}{5,3,4}{}{}
\genomerow{778}{94}{31}{Terminal}{a = value}{4,3,6}{}{-1.41653}
\genomerow{779}{94}{32}{Terminal}{a = value}{0,4,2}{}{197.593}
\genomerow{780}{94}{33}{Nop-2}{perform no operation}{7,5,7}{}{}
\genomerow{781}{94}{34}{Adjust Secondary Global Regulator}{adjust global regulator value by a}{5,5,6}{}{}
\genomerow{782}{94}{35}{Local Anchor}{register a local jump-to destination}{1,1,0}{}{}
\genomerow{783}{94}{36}{Terminal}{a = value}{0,4,2}{}{197.593}
\genomerow{784}{94}{37}{Nop-0}{perform no operation}{7,5,7}{}{}
\genomerow{785}{94}{38}{Adjust Secondary Global Regulator}{adjust global regulator value by a}{5,5,6}{}{}
\genomerow{786}{94}{39}{Local Anchor}{register a local jump-to destination}{1,1,0}{}{}
\genomerow{787}{94}{40}{Local Jump If Not}{if !a, goto tag match local anchor}{6,2,4}{}{}
\genomerow{788}{94}{41}{Generate Random Bool}{if p, a = 1; else a = 0}{4,2,5}{}{}
\genomerow{789}{94}{42}{Send Intra-Cell Message}{if a, send message another cardinal within cell}{5,0,0}{}{}
\genomerow{790}{94}{43}{Less Than}{a = ( b < c )}{2,2,2}{}{}
\genomerow{791}{94}{44}{Terminal}{a = value}{4,2,5}{}{-0.746062}
\genomerow{792}{94}{45}{Adjust Local Regulator}{adjust local regulator value by a}{1,1,3}{}{}
\genomerow{793}{94}{46}{Logical Or}{a = b \textbar{}\textbar{} c}{5,0,5}{}{}
\genomerow{794}{94}{47}{Get Secondary Global Regulator}{a = global regulator value}{5,3,2}{}{}
\genomerow{795}{94}{48}{Logical Or}{a = b \textbar{}\textbar{} c}{5,0,5}{}{}
\genomerow{796}{94}{49}{Get Secondary Global Regulator}{a = global regulator value}{5,3,2}{}{}
\genomerow{797}{94}{50}{Generate Random Bool}{if p, a = 1; else a = 0}{4,2,5}{}{}
\genomerow{798}{94}{51}{Get Secondary Global Regulator}{a = global regulator value}{5,3,2}{}{}
\genomerow{799}{94}{52}{Logical Or}{a = b \textbar{}\textbar{} c}{5,0,5}{}{}
\genomerow{800}{94}{53}{Get Secondary Global Regulator}{a = global regulator value}{5,3,2}{}{}
\genomerow{801}{94}{54}{Generate Random Bool}{if p, a = 1; else a = 0}{4,2,5}{}{}
\genomerow{802}{94}{55}{Send Intra-Cell Message}{if a, send message another cardinal within cell}{5,0,0}{}{}
\genomerow{803}{95}{0}{Global Anchor}{register a global jump-to destination, maybe terminate}{2,2,2}{}{}
\genomerow{804}{95}{1}{Terminal}{a = value}{4,2,5}{}{0.253938}
\genomerow{805}{95}{2}{Adjust Local Regulator}{adjust local regulator value by a}{1,1,3}{}{}
\genomerow{806}{95}{3}{Get Secondary Global Regulator}{a = global regulator value}{5,3,2}{}{}
\genomerow{807}{96}{0}{Global Anchor}{register a global jump-to destination, maybe terminate}{1,6,5}{}{}
\genomerow{808}{96}{1}{Read Own State}{a = state}{0,7,6}{NeighborKinGroupWillExpire[1]}{}
\genomerow{809}{96}{2}{Modulo}{a = b \% c}{7,1,3}{}{}
\genomerow{810}{96}{3}{Global Jump If Not}{if !a, goto global anchor; if !b, clear registers}{5,7,5}{}{}
\genomerow{811}{96}{4}{Negate}{a = -a}{7,7,0}{}{}
\genomerow{812}{96}{5}{Bitwise Xor}{a = b \textasciicircum{} c}{5,0,3}{}{}
\genomerow{813}{96}{6}{Divide}{a = b / c}{5,7,2}{}{}
\genomerow{814}{96}{7}{Global Jump If}{if a, goto tag match global anchor; if b, clear registers}{1,6,7}{}{}
\genomerow{815}{96}{8}{Local Jump If}{if a, goto tag match local anchor}{2,2,3}{}{}
\genomerow{816}{96}{9}{Terminal}{a = value}{4,2,5}{}{0.750031}
\genomerow{817}{96}{10}{Adjust Global Regulator}{adjust global regulator value by a}{4,7,6}{}{}
\genomerow{818}{96}{11}{Terminal}{a = value}{0,0,6}{}{-0.622412}
\genomerow{819}{96}{12}{Read Own State}{a = state}{1,6,3}{OptimumQuorumExceeded[1]}{}
\genomerow{820}{96}{13}{Terminal}{a = value}{2,0,3}{}{-0.431705}
\genomerow{821}{96}{14}{Read Own State}{a = state}{6,5,2}{NumKnownQuorumBits[0]}{}
\genomerow{822}{96}{15}{Local Anchor}{register a local jump-to destination}{4,3,3}{}{}
\genomerow{823}{96}{16}{Logical And}{a = b \&\& c}{7,6,3}{}{}
\genomerow{824}{96}{17}{Local Anchor}{register a local jump-to destination}{3,3,5}{}{}
\genomerow{825}{96}{18}{Divide}{a = b / c}{7,1,6}{}{}
\genomerow{826}{96}{19}{Bitwise Or}{a = b \textbar{} c}{2,5,7}{}{}
\genomerow{827}{96}{20}{Read Own State}{a = state}{1,5,2}{KinGroupMatch[0]}{}
\genomerow{828}{97}{0}{Global Anchor}{register a global jump-to destination, maybe terminate}{7,1,0}{}{}
\genomerow{829}{97}{1}{Terminal}{a = value}{5,0,1}{}{0.377067}
\genomerow{830}{97}{2}{Greater Than}{a = ( b > c )}{7,6,0}{}{}
\genomerow{831}{97}{3}{RandomFill}{fill a with uniformly-distributed random bits}{3,1,2}{}{}
\genomerow{832}{97}{4}{Terminal}{a = value}{5,5,7}{}{-1175.87}
\genomerow{833}{97}{5}{Read Own State}{a = state}{7,7,1}{SpawnRequest[0]}{}
\genomerow{834}{97}{6}{Bitwise And}{a = b \& c}{4,5,4}{}{}
\genomerow{835}{97}{7}{Terminal}{a = value}{0,0,2}{}{-0.34375}
\genomerow{836}{97}{8}{Equal}{a = ( b == c )}{0,1,5}{}{}
\genomerow{837}{97}{9}{Local Anchor}{register a local jump-to destination}{2,3,3}{}{}
\genomerow{838}{97}{10}{Local Jump If}{if a, goto tag match local anchor}{5,3,0}{}{}
\genomerow{839}{97}{11}{Local Anchor}{register a local jump-to destination}{7,3,0}{}{}
\genomerow{840}{97}{12}{Greater Than}{a = ( b > c )}{5,7,6}{}{}
\genomerow{841}{97}{13}{Local Anchor}{register a local jump-to destination}{7,4,6}{}{}
\genomerow{842}{97}{14}{Add To Own State}{state += a}{1,1,4}{ResourceSendRequest[0]}{}
\genomerow{843}{97}{15}{Local Jump If}{if a, goto tag match local anchor}{5,2,0}{}{}
\genomerow{844}{97}{16}{Write Own State}{if b, state = a}{3,0,3}{SpawnArrest[0]}{}
\genomerow{845}{97}{17}{Increment}{a = a + 1}{5,1,5}{}{}
\genomerow{846}{97}{18}{Terminal}{a = value}{5,2,7}{}{-0.604155}
\genomerow{847}{97}{19}{Local Anchor}{register a local jump-to destination}{4,3,2}{}{}
\genomerow{848}{97}{20}{Bitwise Xor}{a = b \textasciicircum{} c}{5,4,3}{}{}
\genomerow{849}{97}{21}{Write Own State}{if b, state = a}{6,7,0}{SpawnRequest[0]}{}
\genomerow{850}{97}{22}{Adjust Global Regulator}{adjust global regulator value by a}{5,3,0}{}{}
\genomerow{851}{97}{23}{Multiply Own State}{state *= a}{5,7,6}{TransientNopState[1]}{}
\genomerow{852}{97}{24}{Local Anchor}{register a local jump-to destination}{7,6,6}{}{}
\genomerow{853}{97}{25}{Add To Own State}{state += a}{1,1,4}{SpawnRequest[0]}{}
\genomerow{854}{97}{26}{Greater Than}{a = ( b > c )}{5,3,4}{}{}
\genomerow{855}{97}{27}{Write Own State}{if b, state = a}{3,0,3}{NopState[0]}{}
\genomerow{856}{97}{28}{Decay Local Regulator}{decay local regulator value by register a}{5,1,5}{}{}
\genomerow{857}{97}{29}{Adjust Global Regulator}{adjust global regulator value by a}{7,3,0}{}{}
\genomerow{858}{97}{30}{Multiply Own State}{state *= a}{5,7,6}{TransientNopState[3]}{}
\genomerow{859}{97}{31}{Fork If}{if a, submit fork request (processed when core terminates)}{1,7,7}{}{}
\genomerow{860}{97}{32}{Local Anchor}{register a local jump-to destination}{0,6,1}{}{}
\genomerow{861}{97}{33}{Global Jump If}{if a, goto tag match global anchor; if b, clear registers}{0,5,1}{}{}
\genomerow{862}{97}{34}{Terminal}{a = value}{5,1,3}{}{2.90619}
\genomerow{863}{97}{35}{Add To Own State}{state += a}{1,3,4}{TransientNopState[1]}{}
\genomerow{864}{97}{36}{Greater Than}{a = ( b > c )}{5,3,6}{}{}
\genomerow{865}{97}{37}{Bitwise Xor}{a = b \textasciicircum{} c}{7,6,4}{}{}
\genomerow{866}{97}{38}{Add To Own State}{state += a}{3,1,5}{NopState[0]}{}
\genomerow{867}{97}{39}{Terminal}{a = value}{5,1,3}{}{2.90619}
\genomerow{868}{97}{40}{Add To Own State}{state += a}{3,1,5}{NopState[0]}{}
\genomerow{869}{97}{41}{Read Own State}{a = state}{6,2,4}{NumKnownQuorumBits[0]}{}
\genomerow{870}{97}{42}{Global Jump If}{if a, goto tag match global anchor; if b, clear registers}{0,4,1}{}{}
\genomerow{871}{97}{43}{Bitwise Xor}{a = b \textasciicircum{} c}{7,4,4}{}{}
\genomerow{872}{97}{44}{Add To Own State}{state += a}{3,1,5}{NopState[0]}{}
\genomerow{873}{97}{45}{Read Own State}{a = state}{6,2,5}{PhylogeneticRootMatch[0]}{}
\genomerow{874}{97}{46}{Global Jump If}{if a, goto tag match global anchor; if b, clear registers}{0,4,1}{}{}
\genomerow{875}{97}{47}{Terminal}{a = value}{5,1,3}{}{2.90619}
\genomerow{876}{97}{48}{Terminal}{a = value}{5,1,3}{}{2.90619}
\genomerow{877}{97}{49}{Read Own State}{a = state}{6,2,5}{KinGroupIDAncestorView[1]}{}
\genomerow{878}{97}{50}{Global Jump If}{if a, goto tag match global anchor; if b, clear registers}{0,4,1}{}{}
\genomerow{879}{97}{51}{Terminal}{a = value}{5,1,3}{}{2.90619}
\genomerow{880}{97}{52}{Terminal}{a = value}{5,1,1}{}{2.90619}
\genomerow{881}{97}{53}{Read Own State}{a = state}{6,2,5}{CellAge[0]}{}
\genomerow{882}{97}{54}{Global Jump If}{if a, goto tag match global anchor; if b, clear registers}{0,4,1}{}{}
\genomerow{883}{97}{55}{Terminal}{a = value}{1,1,3}{}{2.88203}
\genomerow{884}{97}{56}{Terminal}{a = value}{5,1,3}{}{2.90468}
\genomerow{885}{97}{57}{Get Global Regulator}{a = global regulator value}{6,7,5}{}{}
\genomerow{886}{97}{58}{Terminal}{a = value}{6,1,1}{}{0.587469}
\genomerow{887}{97}{59}{Send Intra-Cell Message}{if a, send message another cardinal within cell}{0,7,1}{}{}
\genomerow{888}{97}{60}{Terminal}{a = value}{0,5,5}{}{0.592154}
\genomerow{889}{97}{61}{Decay Global Regulator}{decay global regulator value by register a}{2,3,5}{}{}
\genomerow{890}{97}{62}{Terminal}{a = value}{1,2,4}{}{2.16337}
\genomerow{891}{97}{63}{Read Own State}{a = state}{1,3,0}{RicherThanNeighbor[0]}{}
\genomerow{892}{97}{64}{Set Local Regulator}{set local regulator value to a}{0,3,2}{}{}
\genomerow{893}{97}{65}{Set Local Regulator}{set local regulator value to a}{4,2,4}{}{}
\genomerow{894}{97}{66}{Local Anchor}{register a local jump-to destination}{5,0,4}{}{}
\genomerow{895}{97}{67}{Local Anchor}{register a local jump-to destination}{4,5,2}{}{}
\genomerow{896}{98}{0}{Global Anchor}{register a global jump-to destination, maybe terminate}{2,0,6}{}{}
\genomerow{897}{99}{0}{Global Anchor}{register a global jump-to destination, maybe terminate}{5,6,2}{}{}
\genomerow{898}{99}{1}{Set Secondary Global Regulator}{set global regulator value to a}{5,0,4}{}{}
\genomerow{899}{99}{2}{Set Global Regulator}{set global regulator value to a}{3,6,2}{}{}
\genomerow{900}{99}{3}{Increment}{a = a + 1}{5,0,4}{}{}
\genomerow{901}{99}{4}{Set Secondary Global Regulator}{set global regulator value to a}{5,0,4}{}{}
\genomerow{902}{99}{5}{Set Global Regulator}{set global regulator value to a}{6,4,2}{}{}
\genomerow{903}{99}{6}{Increment}{a = a + 1}{5,0,4}{}{}
\genomerow{904}{100}{0}{Global Anchor}{register a global jump-to destination, maybe terminate}{0,4,3}{}{}
\genomerow{905}{101}{0}{Global Anchor}{register a global jump-to destination, maybe terminate}{3,1,0}{}{}
\genomerow{906}{101}{1}{Adjust Secondary Global Regulator}{adjust global regulator value by a}{5,4,7}{}{}
\genomerow{907}{101}{2}{Global Jump If}{if a, goto tag match global anchor; if b, clear registers}{0,7,3}{}{}
\genomerow{908}{101}{3}{Nop-2}{perform no operation}{5,3,6}{}{}
\genomerow{909}{101}{4}{Local Anchor}{register a local jump-to destination}{0,5,5}{}{}
\genomerow{910}{101}{5}{Write Own State}{if b, state = a}{3,4,6}{NopState[2]}{}
\genomerow{911}{101}{6}{Write Own State}{if b, state = a}{0,7,3}{RepLevRequest[0]}{}
\genomerow{912}{101}{7}{Local Anchor}{register a local jump-to destination}{3,1,1}{}{}
\genomerow{913}{101}{8}{Local Anchor}{register a local jump-to destination}{4,3,1}{}{}
\genomerow{914}{101}{9}{Less Than}{a = ( b < c )}{1,6,7}{}{}
\genomerow{915}{101}{10}{Read Own State}{a = state}{3,1,2}{NeighborKinGroupWillExpire[1]}{}
\genomerow{916}{101}{11}{Subtract}{a = b - c}{0,2,3}{}{}
\genomerow{917}{102}{0}{Global Anchor}{register a global jump-to destination, maybe terminate}{1,7,5}{}{}
\genomerow{918}{102}{1}{Set Secondary Global Regulator}{set global regulator value to a}{5,2,0}{}{}
\genomerow{919}{102}{2}{Terminal}{a = value}{7,0,4}{}{1.84153}
\genomerow{920}{102}{3}{Bitwise Shift}{a = b << c}{2,6,2}{}{}
\genomerow{921}{102}{4}{Increment}{a = a + 1}{5,0,4}{}{}
\genomerow{922}{103}{0}{Global Anchor}{register a global jump-to destination, maybe terminate}{0,4,3}{}{}
\genomerow{923}{103}{1}{Bitwise Shift}{a = b << c}{6,6,2}{}{}
\genomerow{924}{103}{2}{Increment}{a = a + 1}{5,0,4}{}{}
\genomerow{925}{104}{0}{Global Anchor}{register a global jump-to destination, maybe terminate}{0,0,3}{}{}
\genomerow{926}{104}{1}{Local Anchor}{register a local jump-to destination}{1,7,5}{}{}
\genomerow{927}{104}{2}{Subtract}{a = b - c}{1,3,3}{}{}
\genomerow{928}{104}{3}{Decay Secondary Global Regulator}{decay global regulator value by register a}{2,6,7}{}{}
\genomerow{929}{104}{4}{Terminate If}{if a, terminate core}{5,3,7}{}{}
\genomerow{930}{104}{5}{Read Own State}{a = state}{6,2,5}{Epoch[0]}{}
\genomerow{931}{105}{0}{Global Anchor}{register a global jump-to destination, maybe terminate}{2,6,4}{}{}
\genomerow{932}{105}{1}{Local Anchor}{register a local jump-to destination}{3,2,4}{}{}
\genomerow{933}{105}{2}{Read Neighbor State}{a = state}{7,1,6}{IsParentGroupOf[1]}{}
\genomerow{934}{106}{0}{Global Anchor}{register a global jump-to destination, maybe terminate}{4,6,5}{}{}
\genomerow{935}{106}{1}{Decrement}{a = a - 1}{1,2,0}{}{}
\genomerow{936}{106}{2}{Terminal}{a = value}{1,1,1}{}{-0.951405}
\genomerow{937}{106}{3}{Logical Or}{a = b \textbar{}\textbar{} c}{1,3,0}{}{}
\genomerow{938}{106}{4}{Bitwise Xor}{a = b \textasciicircum{} c}{5,2,2}{}{}
\genomerow{939}{106}{5}{Local Anchor}{register a local jump-to destination}{5,2,0}{}{}
\genomerow{940}{106}{6}{Decay Secondary Global Regulator}{decay global regulator value by register a}{2,4,7}{}{}
\genomerow{941}{106}{7}{Terminal}{a = value}{2,0,6}{}{0.265235}
\genomerow{942}{106}{8}{Terminal}{a = value}{2,2,7}{}{-42.6459}
\genomerow{943}{106}{9}{Terminal}{a = value}{2,5,5}{}{1.87245}
\genomerow{944}{106}{10}{Send Intra-Cell Message}{if a, send message another cardinal within cell}{3,5,1}{}{}
\genomerow{945}{106}{11}{Local Anchor}{register a local jump-to destination}{1,2,0}{}{}
\genomerow{946}{106}{12}{Greater Than}{a = ( b > c )}{3,3,5}{}{}
\genomerow{947}{106}{13}{Local Anchor}{register a local jump-to destination}{2,0,6}{}{}
\genomerow{948}{106}{14}{Terminal}{a = value}{2,5,5}{}{1.87245}
\genomerow{949}{107}{0}{Global Anchor}{register a global jump-to destination, maybe terminate}{3,5,1}{}{}
\genomerow{950}{107}{1}{Local Anchor}{register a local jump-to destination}{1,2,0}{}{}
\genomerow{951}{107}{2}{Not Equal}{a = (b != c)}{3,3,1}{}{}
\genomerow{952}{107}{3}{Local Anchor}{register a local jump-to destination}{2,0,6}{}{}
\genomerow{953}{107}{4}{Local Anchor}{register a local jump-to destination}{6,6,5}{}{}
\genomerow{954}{107}{5}{Generate Random Bool}{if p, a = 1; else a = 0}{7,3,2}{}{}
\genomerow{955}{107}{6}{Local Anchor}{register a local jump-to destination}{3,6,2}{}{}
\genomerow{956}{107}{7}{Logical And}{a = b \&\& c}{2,0,2}{}{}
\genomerow{957}{108}{0}{Global Anchor}{register a global jump-to destination, maybe terminate}{7,6,7}{}{}
\genomerow{958}{108}{1}{RandomFill}{fill a with uniformly-distributed random bits}{5,4,6}{}{}
\genomerow{959}{108}{2}{Nop-2}{perform no operation}{5,4,2}{}{}
\genomerow{960}{108}{3}{Logical And}{a = b \&\& c}{4,1,6}{}{}
\genomerow{961}{108}{4}{Terminal}{a = value}{0,3,7}{}{0.428078}
\genomerow{962}{109}{0}{Global Anchor}{register a global jump-to destination, maybe terminate}{1,0,3}{}{}
\genomerow{963}{109}{1}{Read Own State}{a = state}{0,4,2}{NeighborKinGroupWillExpire[0]}{}
\genomerow{964}{109}{2}{Multiply Own State}{state *= a}{6,0,7}{TransientNopState[2]}{}
\genomerow{965}{109}{3}{Write Own State}{if b, state = a}{5,1,0}{NopState[0]}{}
\genomerow{966}{110}{0}{Global Anchor}{register a global jump-to destination, maybe terminate}{6,1,3}{}{}
\genomerow{967}{110}{1}{Add To Own State}{state += a}{4,2,1}{SpawnRequest[0]}{}
\genomerow{968}{110}{2}{Send Inter-Cell Message}{if a, send message to neighbor cell}{5,1,4}{}{}
\genomerow{969}{111}{0}{Global Anchor}{register a global jump-to destination, maybe terminate}{6,1,3}{}{}
\genomerow{970}{111}{1}{Local Jump If}{if a, goto tag match local anchor}{5,7,6}{}{}
\genomerow{971}{112}{0}{Global Anchor}{register a global jump-to destination, maybe terminate}{0,1,7}{}{}
\genomerow{972}{112}{1}{Get Global Regulator}{a = global regulator value}{1,0,3}{}{}
\genomerow{973}{113}{0}{Global Anchor}{register a global jump-to destination, maybe terminate}{7,5,0}{}{}
\genomerow{974}{113}{1}{Get Local Regulator}{a = local regulator value}{3,7,6}{}{}
\genomerow{975}{113}{2}{Get Global Regulator}{a = global regulator value}{1,0,3}{}{}
\genomerow{976}{114}{0}{Global Anchor}{register a global jump-to destination, maybe terminate}{7,5,0}{}{}
\genomerow{977}{114}{1}{Get Local Regulator}{a = local regulator value}{7,7,6}{}{}
\genomerow{978}{114}{2}{Local Anchor}{register a local jump-to destination}{6,3,4}{}{}
\genomerow{979}{114}{3}{Get Local Regulator}{a = local regulator value}{4,2,4}{}{}
\genomerow{980}{114}{4}{Local Anchor}{register a local jump-to destination}{6,3,6}{}{}
\genomerow{981}{114}{5}{Add To Own State}{state += a}{1,2,5}{TransientNopState[1]}{}
\genomerow{982}{114}{6}{Terminal}{a = value}{7,2,5}{}{-2.08977}
\genomerow{983}{114}{7}{Terminal}{a = value}{4,2,2}{}{-383.982}
\genomerow{984}{114}{8}{Local Anchor}{register a local jump-to destination}{7,1,4}{}{}
\genomerow{985}{114}{9}{Local Anchor}{register a local jump-to destination}{2,6,3}{}{}
\genomerow{986}{114}{10}{Terminal}{a = value}{4,3,2}{}{-385.01}
\genomerow{987}{114}{11}{Bitwise Shift}{a = b << c}{7,1,4}{}{}
\genomerow{988}{114}{12}{Local Anchor}{register a local jump-to destination}{5,5,6}{}{}
\genomerow{989}{114}{13}{Bitwise And}{a = b \& c}{3,7,6}{}{}
\genomerow{990}{114}{14}{Send Intra-Cell Message}{if a, send message another cardinal within cell}{1,0,4}{}{}
\genomerow{991}{114}{15}{Terminal}{a = value}{2,2,4}{}{3.47062}
\genomerow{992}{114}{16}{Local Anchor}{register a local jump-to destination}{5,3,6}{}{}
\genomerow{993}{114}{17}{Read Own State}{a = state}{3,7,1}{NeighborFragmented[0]}{}
\genomerow{994}{114}{18}{Nop-2}{perform no operation}{1,0,6}{}{}
\genomerow{995}{114}{19}{Terminal}{a = value}{1,7,6}{}{-0.348885}
\genomerow{996}{114}{20}{Bitwise Not}{a = \textasciitilde{}b}{1,2,6}{}{}
\genomerow{997}{114}{21}{Read Own State}{a = state}{3,7,1}{NeighborFragmented[0]}{}
\genomerow{998}{114}{22}{Nop-2}{perform no operation}{1,0,6}{}{}
\genomerow{999}{114}{23}{Terminal}{a = value}{1,7,6}{}{-0.348885}
\genomerow{1000}{114}{24}{Bitwise Not}{a = \textasciitilde{}b}{1,2,6}{}{}
\genomerow{1001}{114}{25}{Read Own State}{a = state}{4,5,5}{KinGroupMatch[1]}{}
\genomerow{1002}{114}{26}{Multiply Own State}{state *= a}{5,0,4}{ResourceReceiveResistance[0]}{}
\genomerow{1003}{114}{27}{Read Own State}{a = state}{3,7,1}{KinGroupIDView[0]}{}
\genomerow{1004}{114}{28}{Nop-2}{perform no operation}{5,0,6}{}{}
\genomerow{1005}{114}{29}{Nop-2}{perform no operation}{1,0,6}{}{}
\genomerow{1006}{114}{30}{Add To Own State}{state += a}{1,7,6}{ResourceSendLimit[0]}{}
\genomerow{1007}{115}{0}{Global Anchor}{register a global jump-to destination, maybe terminate}{0,1,2}{}{}
\genomerow{1008}{115}{1}{Bitwise Xor}{a = b \textasciicircum{} c}{0,6,0}{}{}
\genomerow{1009}{116}{0}{Global Anchor}{register a global jump-to destination, maybe terminate}{4,1,6}{}{}
\genomerow{1010}{116}{1}{Bitwise And}{a = b \& c}{6,4,3}{}{}
\genomerow{1011}{116}{2}{Read Neighbor State}{a = state}{1,6,2}{IsParentCellOf[0]}{}
\genomerow{1012}{116}{3}{Global Jump If Not}{if !a, goto global anchor; if !b, clear registers}{2,5,3}{}{}
\genomerow{1013}{116}{4}{Write Own State}{if b, state = a}{5,2,4}{TransientNopState[0]}{}
\genomerow{1014}{116}{5}{Read Own State}{a = state}{3,6,6}{NopState[3]}{}
\genomerow{1015}{116}{6}{Generate Random Bool}{if p, a = 1; else a = 0}{3,2,5}{}{}
\genomerow{1016}{116}{7}{Increment}{a = a + 1}{3,2,5}{}{}
\genomerow{1017}{116}{8}{Local Anchor}{register a local jump-to destination}{7,7,1}{}{}
\genomerow{1018}{116}{9}{Multiply}{a = b * c}{6,5,1}{}{}
\genomerow{1019}{116}{10}{Set Secondary Global Regulator}{set global regulator value to a}{4,3,3}{}{}
\genomerow{1020}{116}{11}{Draw Random Value}{a = weighted draw from random number generator}{4,3,3}{}{}
\genomerow{1021}{116}{12}{Adjust Local Regulator}{adjust local regulator value by a}{2,5,7}{}{}
\genomerow{1022}{116}{13}{Local Anchor}{register a local jump-to destination}{7,3,1}{}{}
\genomerow{1023}{116}{14}{Draw Random Value}{a = weighted draw from random number generator}{4,3,1}{}{}
\genomerow{1024}{116}{15}{Local Anchor}{register a local jump-to destination}{3,4,6}{}{}
\genomerow{1025}{116}{16}{Terminal}{a = value}{3,6,5}{}{1.31542}
\genomerow{1026}{116}{17}{Count Ones}{a = popcnt( b )}{3,1,7}{}{}
\genomerow{1027}{116}{18}{Read Own State}{a = state}{4,2,3}{StockpileDepleted[0]}{}
\genomerow{1028}{116}{19}{Read Own State}{a = state}{1,6,4}{IsParentCellOf[0]}{}
\genomerow{1029}{117}{0}{Global Anchor}{register a global jump-to destination, maybe terminate}{3,4,2}{}{}
\genomerow{1030}{117}{1}{Write Own State}{if b, state = a}{3,6,7}{ResourceReserveRequest[0]}{}
\genomerow{1031}{117}{2}{Count Ones}{a = popcnt( b )}{3,5,7}{}{}
\genomerow{1032}{117}{3}{Generate Random Bool}{if p, a = 1; else a = 0}{5,6,7}{}{}
\genomerow{1033}{117}{4}{Read Own State}{a = state}{4,2,3}{StockpileDepleted[0]}{}
\genomerow{1034}{117}{5}{Read Own State}{a = state}{0,6,4}{IsParentCellOf[0]}{}
\genomerow{1035}{118}{0}{Global Anchor}{register a global jump-to destination, maybe terminate}{3,4,2}{}{}
\genomerow{1036}{118}{1}{Write Own State}{if b, state = a}{3,6,7}{ResourceReserveRequest[0]}{}
\genomerow{1037}{118}{2}{Count Ones}{a = popcnt( b )}{3,5,7}{}{}
\genomerow{1038}{118}{3}{Generate Random Bool}{if p, a = 1; else a = 0}{5,6,7}{}{}
\genomerow{1039}{118}{4}{Terminal}{a = value}{0,2,4}{}{2.96227}
\genomerow{1040}{118}{5}{Read Neighbor State}{a = state}{3,6,5}{PhylogeneticRootMatch[0]}{}
\genomerow{1041}{118}{6}{Bitwise Xor}{a = b \textasciicircum{} c}{3,5,7}{}{}
\genomerow{1042}{118}{7}{Not}{a = !a}{5,6,7}{}{}
\genomerow{1043}{118}{8}{Terminal}{a = value}{0,2,4}{}{2.96227}
\genomerow{1044}{118}{9}{Decay Global Regulator}{decay global regulator value by register a}{6,4,3}{}{}
\genomerow{1045}{119}{0}{Global Anchor}{register a global jump-to destination, maybe terminate}{0,0,5}{}{}
\genomerow{1046}{119}{1}{Send Inter-Cell Message}{if a, send message to neighbor cell}{7,6,0}{}{}
\genomerow{1047}{119}{2}{Terminal}{a = value}{0,2,2}{}{0.0234738}
\genomerow{1048}{119}{3}{Send Intra-Cell Message}{if a, send message another cardinal within cell}{4,2,7}{}{}
\genomerow{1049}{119}{4}{Terminal}{a = value}{2,6,4}{}{-5892.24}
\genomerow{1050}{119}{5}{Bitwise Shift}{a = b << c}{0,7,3}{}{}
\genomerow{1051}{119}{6}{Decay Local Regulator}{decay local regulator value by register a}{5,7,6}{}{}
\genomerow{1052}{119}{7}{Terminal}{a = value}{7,5,4}{}{1.99918}
\genomerow{1053}{119}{8}{Terminal}{a = value}{6,2,0}{}{-0.0674628}
\genomerow{1054}{119}{9}{Adjust Global Regulator}{adjust global regulator value by a}{0,4,3}{}{}
\genomerow{1055}{119}{10}{Terminal}{a = value}{5,2,1}{}{6.17348}
\genomerow{1056}{119}{11}{Local Anchor}{register a local jump-to destination}{3,7,6}{}{}
\genomerow{1057}{119}{12}{Terminal}{a = value}{7,5,4}{}{39.5265}
\genomerow{1058}{119}{13}{Multiply Own State}{state *= a}{3,3,4}{TransientNopState[2]}{}
\genomerow{1059}{119}{14}{Send Intra-Cell Message}{if a, send message another cardinal within cell}{1,1,4}{}{}
\genomerow{1060}{119}{15}{Send Intra-Cell Message}{if a, send message another cardinal within cell}{2,1,6}{}{}
\genomerow{1061}{119}{16}{Adjust Global Regulator}{adjust global regulator value by a}{0,7,5}{}{}
\genomerow{1062}{119}{17}{Local Anchor}{register a local jump-to destination}{6,1,6}{}{}
\genomerow{1063}{119}{18}{Send Intra-Cell Message}{if a, send message another cardinal within cell}{1,1,4}{}{}
\genomerow{1064}{119}{19}{Send Intra-Cell Message}{if a, send message another cardinal within cell}{6,1,6}{}{}
\genomerow{1065}{119}{20}{Local Anchor}{register a local jump-to destination}{0,0,6}{}{}
\genomerow{1066}{119}{21}{Bitwise Xor}{a = b \textasciicircum{} c}{2,2,0}{}{}
\genomerow{1067}{119}{22}{Less Than}{a = ( b < c )}{0,4,4}{}{}
\genomerow{1068}{119}{23}{Send Intra-Cell Message}{if a, send message another cardinal within cell}{6,1,6}{}{}
\genomerow{1069}{119}{24}{Set Global Regulator}{set global regulator value to a}{6,1,7}{}{}
\genomerow{1070}{119}{25}{Adjust Local Regulator}{adjust local regulator value by a}{5,1,1}{}{}
\genomerow{1071}{119}{26}{Decay Local Regulator}{decay local regulator value by register a}{2,1,0}{}{}
\genomerow{1072}{119}{27}{Send Inter-Cell Message}{if a, send message to neighbor cell}{5,0,5}{}{}
\genomerow{1073}{119}{28}{Get Local Regulator}{a = local regulator value}{7,5,7}{}{}
\genomerow{1074}{119}{29}{Local Anchor}{register a local jump-to destination}{7,2,3}{}{}
\genomerow{1075}{119}{30}{Read Own State}{a = state}{2,5,0}{IncomingIntraMessageCounter[0]}{}
\genomerow{1076}{119}{31}{Local Anchor}{register a local jump-to destination}{4,5,2}{}{}
\genomerow{1077}{119}{32}{Read Own State}{a = state}{4,4,1}{IsParentCellOf[0]}{}
\genomerow{1078}{119}{33}{Read Own State}{a = state}{4,4,1}{IsParentCellOf[0]}{}
\genomerow{1079}{119}{34}{Logical Or}{a = b \textbar{}\textbar{} c}{6,3,4}{}{}
\genomerow{1080}{119}{35}{Nop-2}{perform no operation}{0,4,0}{}{}
\genomerow{1081}{119}{36}{Read Own State}{a = state}{4,1,1}{ResourceSendRequest[0]}{}
\genomerow{1082}{119}{37}{Local Anchor}{register a local jump-to destination}{4,5,2}{}{}
\genomerow{1083}{119}{38}{Draw Random Value}{a = weighted draw from random number generator}{3,0,3}{}{}
\genomerow{1084}{119}{39}{Read Own State}{a = state}{4,4,1}{NeighborFragmented[0]}{}
\genomerow{1085}{119}{40}{Logical Or}{a = b \textbar{}\textbar{} c}{6,3,4}{}{}
\genomerow{1086}{119}{41}{Nop-2}{perform no operation}{0,4,0}{}{}
\genomerow{1087}{119}{42}{Read Own State}{a = state}{4,1,0}{ResourceSendRequest[0]}{}
\genomerow{1088}{119}{43}{Send Intra-Cell Message}{if a, send message another cardinal within cell}{6,6,6}{}{}
\genomerow{1089}{119}{44}{Bitwise And}{a = b \& c}{5,1,7}{}{}
\genomerow{1090}{119}{45}{Read Own State}{a = state}{4,1,0}{ResourceSendRequest[0]}{}
\genomerow{1091}{119}{46}{Send Intra-Cell Message}{if a, send message another cardinal within cell}{6,6,6}{}{}
\genomerow{1092}{119}{47}{Bitwise And}{a = b \& c}{5,1,7}{}{}
\genomerow{1093}{119}{48}{Less Than}{a = ( b < c )}{0,4,4}{}{}
\genomerow{1094}{119}{49}{Count Ones}{a = popcnt( b )}{1,7,6}{}{}
\genomerow{1095}{119}{50}{Set Local Regulator}{set local regulator value to a}{6,6,2}{}{}
\genomerow{1096}{119}{51}{Not Equal}{a = (b != c)}{6,7,6}{}{}
\genomerow{1097}{119}{52}{Bitwise Xor}{a = b \textasciicircum{} c}{5,1,7}{}{}
\genomerow{1098}{119}{53}{Less Than}{a = ( b < c )}{1,4,4}{}{}
\genomerow{1099}{119}{54}{Local Anchor}{register a local jump-to destination}{3,7,6}{}{}
\genomerow{1100}{120}{0}{Global Anchor}{register a global jump-to destination, maybe terminate}{0,5,3}{}{}
\genomerow{1101}{120}{1}{Read Own State}{a = state}{7,4,0}{IncomingInterMessageCounter[0]}{}
\genomerow{1102}{120}{2}{Read Own State}{a = state}{6,3,0}{NumKnownQuorumBits[0]}{}
\genomerow{1103}{121}{0}{Global Anchor}{register a global jump-to destination, maybe terminate}{0,5,3}{}{}
\genomerow{1104}{121}{1}{Read Own State}{a = state}{7,4,0}{NumKnownQuorumBits[0]}{}
\genomerow{1105}{121}{2}{Read Own State}{a = state}{6,3,0}{IsAlive[0]}{}
\genomerow{1106}{121}{3}{Local Anchor}{register a local jump-to destination}{2,4,3}{}{}
\genomerow{1107}{121}{4}{Read Own State}{a = state}{7,4,0}{NumKnownQuorumBits[0]}{}
\genomerow{1108}{121}{5}{Read Own State}{a = state}{7,4,0}{NumKnownQuorumBits[0]}{}
\genomerow{1109}{121}{6}{Read Own State}{a = state}{6,3,0}{IsAlive[0]}{}
\genomerow{1110}{122}{0}{Global Anchor}{register a global jump-to destination, maybe terminate}{6,4,3}{}{}
\genomerow{1111}{122}{1}{Set Global Regulator}{set global regulator value to a}{1,5,5}{}{}
\genomerow{1112}{122}{2}{Set Global Regulator}{set global regulator value to a}{1,5,4}{}{}
\genomerow{1113}{122}{3}{Local Anchor}{register a local jump-to destination}{5,1,4}{}{}
\genomerow{1114}{122}{4}{Logical Or}{a = b \textbar{}\textbar{} c}{3,2,3}{}{}
\genomerow{1115}{122}{5}{Bitwise Xor}{a = b \textasciicircum{} c}{2,6,2}{}{}
\genomerow{1116}{122}{6}{Terminal}{a = value}{4,4,0}{}{-9.53691}
\genomerow{1117}{122}{7}{Read Own State}{a = state}{7,7,7}{SpawnCount[0]}{}
\genomerow{1118}{122}{8}{Get Local Regulator}{a = local regulator value}{7,3,2}{}{}
\genomerow{1119}{122}{9}{Nop-2}{perform no operation}{0,2,5}{}{}
\genomerow{1120}{122}{10}{Local Anchor}{register a local jump-to destination}{5,1,0}{}{}
\genomerow{1121}{122}{11}{Adjust Local Regulator}{adjust local regulator value by a}{3,0,3}{}{}
\genomerow{1122}{122}{12}{Read Own State}{a = state}{6,7,1}{RicherThanNeighbor[0]}{}
\genomerow{1123}{123}{0}{Global Anchor}{register a global jump-to destination, maybe terminate}{6,2,4}{}{}
\genomerow{1124}{123}{1}{Not Equal}{a = (b != c)}{6,1,5}{}{}
\genomerow{1125}{123}{2}{Read Neighbor State}{a = state}{1,7,3}{KinGroupMatch[0]}{}
\genomerow{1126}{123}{3}{Multiply Own State}{state *= a}{6,6,2}{NopState[0]}{}
\genomerow{1127}{123}{4}{Local Anchor}{register a local jump-to destination}{3,1,0}{}{}
\genomerow{1128}{123}{5}{Local Anchor}{register a local jump-to destination}{0,4,0}{}{}
\genomerow{1129}{123}{6}{Add To Own State}{state += a}{6,6,2}{RepLevRequest[0]}{}
\genomerow{1130}{123}{7}{Local Anchor}{register a local jump-to destination}{0,4,0}{}{}
\genomerow{1131}{123}{8}{Equal}{a = ( b == c )}{7,5,0}{}{}
\genomerow{1132}{123}{9}{Global Jump If}{if a, goto tag match global anchor; if b, clear registers}{2,6,7}{}{}
\genomerow{1133}{123}{10}{Local Anchor}{register a local jump-to destination}{0,4,0}{}{}
\genomerow{1134}{123}{11}{Add To Own State}{state += a}{4,6,2}{RepLevRequest[0]}{}
\genomerow{1135}{123}{12}{Local Anchor}{register a local jump-to destination}{0,4,0}{}{}
\genomerow{1136}{123}{13}{Equal}{a = ( b == c )}{7,5,0}{}{}
\genomerow{1137}{123}{14}{Global Jump If}{if a, goto tag match global anchor; if b, clear registers}{2,6,7}{}{}
\genomerow{1138}{123}{15}{Set Secondary Global Regulator}{set global regulator value to a}{4,0,6}{}{}
\genomerow{1139}{123}{16}{Terminal}{a = value}{0,4,4}{}{1376.11}
\genomerow{1140}{123}{17}{Less Than}{a = ( b < c )}{4,3,7}{}{}
\genomerow{1141}{123}{18}{Terminal}{a = value}{2,7,1}{}{0.849186}
\genomerow{1142}{123}{19}{Add To Own State}{state += a}{6,6,2}{RepLevRequest[0]}{}
\genomerow{1143}{123}{20}{Local Anchor}{register a local jump-to destination}{0,4,0}{}{}
\genomerow{1144}{123}{21}{Equal}{a = ( b == c )}{7,5,0}{}{}
\genomerow{1145}{123}{22}{Write Own State}{if b, state = a}{0,1,0}{SpawnRequest[0]}{}
\genomerow{1146}{123}{23}{Local Jump If Not}{if !a, goto tag match local anchor}{3,5,0}{}{}
\genomerow{1147}{123}{24}{Count Ones}{a = popcnt( b )}{6,1,4}{}{}
\genomerow{1148}{123}{25}{Modulo}{a = b \% c}{4,7,7}{}{}
\genomerow{1149}{123}{26}{Send Intra-Cell Message}{if a, send message another cardinal within cell}{6,6,0}{}{}
\genomerow{1150}{123}{27}{Terminal}{a = value}{5,1,5}{}{0.439249}
\genomerow{1151}{123}{28}{Global Jump If}{if a, goto tag match global anchor; if b, clear registers}{2,6,7}{}{}
\genomerow{1152}{123}{29}{Set Secondary Global Regulator}{set global regulator value to a}{4,0,6}{}{}
\genomerow{1153}{123}{30}{Nop-1}{perform no operation}{5,0,6}{}{}
\genomerow{1154}{123}{31}{Terminal}{a = value}{6,3,7}{}{-10.3477}
\genomerow{1155}{123}{32}{Get Global Regulator}{a = global regulator value}{5,6,6}{}{}
\genomerow{1156}{123}{33}{Local Anchor}{register a local jump-to destination}{5,5,0}{}{}
\genomerow{1157}{123}{34}{Terminal}{a = value}{2,6,3}{}{-2.06402}
\genomerow{1158}{123}{35}{Local Anchor}{register a local jump-to destination}{6,4,2}{}{}
\genomerow{1159}{123}{36}{Terminal}{a = value}{6,1,5}{}{14.1507}
\genomerow{1160}{123}{37}{Terminal}{a = value}{7,1,7}{}{7.25388}
\genomerow{1161}{123}{38}{Write Own State}{if b, state = a}{0,3,7}{RepLevRequest[0]}{}
\genomerow{1162}{123}{39}{Adjust Secondary Global Regulator}{adjust global regulator value by a}{4,5,5}{}{}
\genomerow{1163}{123}{40}{Logical And}{a = b \&\& c}{0,5,1}{}{}
\genomerow{1164}{123}{41}{Local Anchor}{register a local jump-to destination}{6,4,7}{}{}
\genomerow{1165}{123}{42}{Local Anchor}{register a local jump-to destination}{0,3,6}{}{}
\genomerow{1166}{124}{0}{Global Anchor}{register a global jump-to destination, maybe terminate}{4,2,7}{}{}
\genomerow{1167}{124}{1}{Send Intra-Cell Message}{if a, send message another cardinal within cell}{5,2,7}{}{}
\genomerow{1168}{124}{2}{Read Own State}{a = state}{2,6,3}{NeighborOptimumQuorumExceeded[1]}{}
\genomerow{1169}{124}{3}{Broadcast Intra-Cell Message}{if a, send message to all other cardinals within cell}{7,6,2}{}{}
\genomerow{1170}{124}{4}{Less Than}{a = ( b < c )}{6,4,6}{}{}
\genomerow{1171}{124}{5}{Local Anchor}{register a local jump-to destination}{4,2,1}{}{}
\genomerow{1172}{124}{6}{Read Own State}{a = state}{2,6,3}{NeighborOptimumQuorumExceeded[1]}{}
\genomerow{1173}{124}{7}{Broadcast Intra-Cell Message}{if a, send message to all other cardinals within cell}{7,6,2}{}{}
\genomerow{1174}{124}{8}{Less Than}{a = ( b < c )}{6,4,6}{}{}
\genomerow{1175}{124}{9}{Local Anchor}{register a local jump-to destination}{4,3,1}{}{}
\genomerow{1176}{125}{0}{Global Anchor}{register a global jump-to destination, maybe terminate}{4,2,6}{}{}
\genomerow{1177}{126}{0}{Global Anchor}{register a global jump-to destination, maybe terminate}{2,1,1}{}{}
\genomerow{1178}{126}{1}{Not Equal}{a = (b != c)}{2,4,3}{}{}
\genomerow{1179}{126}{2}{Read Neighbor State}{a = state}{0,6,3}{IsChildGroupOf[1]}{}
\genomerow{1180}{126}{3}{Read Own State}{a = state}{0,0,0}{IsParentGroupOf[1]}{}
\genomerow{1181}{126}{4}{Local Anchor}{register a local jump-to destination}{4,2,1}{}{}
\genomerow{1182}{127}{0}{Global Anchor}{register a global jump-to destination, maybe terminate}{2,1,1}{}{}
\genomerow{1183}{127}{1}{Not Equal}{a = (b != c)}{2,4,3}{}{}
\genomerow{1184}{127}{2}{Read Neighbor State}{a = state}{0,6,3}{KinGroupIDView[1]}{}
\genomerow{1185}{127}{3}{Read Own State}{a = state}{0,0,0}{IsParentGroupOf[1]}{}
\genomerow{1186}{127}{4}{Local Anchor}{register a local jump-to destination}{4,2,1}{}{}
\genomerow{1187}{128}{0}{Global Anchor}{register a global jump-to destination, maybe terminate}{2,1,1}{}{}
\genomerow{1188}{128}{1}{Not Equal}{a = (b != c)}{2,4,3}{}{}
\genomerow{1189}{128}{2}{Read Neighbor State}{a = state}{0,6,3}{IsNewborn[0]}{}
\genomerow{1190}{128}{3}{Read Own State}{a = state}{0,0,0}{IsChildGroupOf[1]}{}
\genomerow{1191}{128}{4}{Global Jump If}{if a, goto tag match global anchor; if b, clear registers}{2,4,2}{}{}
\genomerow{1192}{128}{5}{Send Intra-Cell Message}{if a, send message another cardinal within cell}{3,0,3}{}{}
\genomerow{1193}{128}{6}{Terminal}{a = value}{5,4,6}{}{516.236}
\genomerow{1194}{128}{7}{Divide}{a = b / c}{2,7,4}{}{}
\genomerow{1195}{128}{8}{Less Than}{a = ( b < c )}{0,0,5}{}{}
\genomerow{1196}{128}{9}{Increment}{a = a + 1}{5,5,0}{}{}
\genomerow{1197}{128}{10}{Adjust Local Regulator}{adjust local regulator value by a}{2,2,1}{}{}
\genomerow{1198}{128}{11}{Fork If}{if a, submit fork request (processed when core terminates)}{1,5,1}{}{}
\genomerow{1199}{128}{12}{Terminal}{a = value}{4,2,7}{}{0.325813}
\genomerow{1200}{128}{13}{Local Jump If}{if a, goto tag match local anchor}{4,4,6}{}{}
\genomerow{1201}{128}{14}{Local Jump If Not}{if !a, goto tag match local anchor}{7,5,0}{}{}
\genomerow{1202}{128}{15}{Terminate If}{if a, terminate core}{4,6,2}{}{}
\genomerow{1203}{128}{16}{Terminal}{a = value}{5,5,5}{}{4.61325}
\genomerow{1204}{128}{17}{Local Anchor}{register a local jump-to destination}{1,0,7}{}{}
\genomerow{1205}{128}{18}{Generate Random Bool}{if p, a = 1; else a = 0}{1,7,4}{}{}
\genomerow{1206}{129}{0}{Global Anchor}{register a global jump-to destination, maybe terminate}{2,5,0}{}{}
\genomerow{1207}{129}{1}{Bitwise And}{a = b \& c}{1,0,3}{}{}
\genomerow{1208}{129}{2}{Get Local Regulator}{a = local regulator value}{2,1,1}{}{}
\genomerow{1209}{129}{3}{Terminal}{a = value}{2,5,6}{}{-1.72006}
\genomerow{1210}{129}{4}{Negate}{a = -a}{0,5,0}{}{}
\genomerow{1211}{129}{5}{Terminal}{a = value}{1,0,6}{}{3.35032}
\genomerow{1212}{129}{6}{Modulo}{a = b \% c}{1,0,7}{}{}
\genomerow{1213}{129}{7}{Terminal}{a = value}{4,2,7}{}{28.9132}
\genomerow{1214}{129}{8}{Local Anchor}{register a local jump-to destination}{4,4,6}{}{}
\genomerow{1215}{129}{9}{Local Anchor}{register a local jump-to destination}{6,6,0}{}{}
\genomerow{1216}{129}{10}{Nop-0}{perform no operation}{1,1,1}{}{}
\genomerow{1217}{129}{11}{Set Local Regulator}{set local regulator value to a}{7,0,3}{}{}
\genomerow{1218}{129}{12}{Write Own State}{if b, state = a}{3,7,4}{ResourceSendLimit[0]}{}
\genomerow{1219}{129}{13}{Broadcast Intra-Cell Message}{if a, send message to all other cardinals within cell}{5,3,3}{}{}
\genomerow{1220}{129}{14}{Terminal}{a = value}{5,3,6}{}{131.774}
\genomerow{1221}{129}{15}{Generate Random Bool}{if p, a = 1; else a = 0}{7,2,4}{}{}
\genomerow{1222}{129}{16}{Local Anchor}{register a local jump-to destination}{2,2,4}{}{}
\genomerow{1223}{129}{17}{Bitwise Xor}{a = b \textasciicircum{} c}{4,5,6}{}{}
\genomerow{1224}{130}{0}{Global Anchor}{register a global jump-to destination, maybe terminate}{5,6,6}{}{}
\genomerow{1225}{130}{1}{Terminal}{a = value}{2,0,3}{}{-2.22184}
\genomerow{1226}{130}{2}{Terminate If}{if a, terminate core}{4,1,4}{}{}
\genomerow{1227}{130}{3}{Bitwise Xor}{a = b \textasciicircum{} c}{7,0,4}{}{}
\genomerow{1228}{130}{4}{Terminal}{a = value}{7,1,6}{}{-14.4179}
\genomerow{1229}{130}{5}{Terminal}{a = value}{1,4,6}{}{-0.344612}
\genomerow{1230}{130}{6}{Terminal}{a = value}{2,7,5}{}{-5.36944}
\genomerow{1231}{130}{7}{Set Local Regulator}{set local regulator value to a}{3,5,4}{}{}
\genomerow{1232}{130}{8}{Add To Own State}{state += a}{4,3,6}{HeirRequest[0]}{}
\genomerow{1233}{130}{9}{Terminal}{a = value}{2,6,0}{}{-299.121}
\genomerow{1234}{130}{10}{Add To Own State}{state += a}{2,3,7}{ResourceSendLimit[0]}{}
\genomerow{1235}{130}{11}{Local Jump If}{if a, goto tag match local anchor}{4,6,0}{}{}
\genomerow{1236}{130}{12}{Less Than}{a = ( b < c )}{4,6,6}{}{}
\genomerow{1237}{130}{13}{Terminal}{a = value}{2,0,3}{}{-2.22184}
\genomerow{1238}{130}{14}{Bitwise Xor}{a = b \textasciicircum{} c}{0,2,4}{}{}
\genomerow{1239}{130}{15}{Global Jump If}{if a, goto tag match global anchor; if b, clear registers}{6,0,4}{}{}
\genomerow{1240}{130}{16}{Local Anchor}{register a local jump-to destination}{0,2,4}{}{}
\genomerow{1241}{130}{17}{Global Jump If}{if a, goto tag match global anchor; if b, clear registers}{6,0,4}{}{}
\genomerow{1242}{130}{18}{Terminal}{a = value}{4,6,1}{}{2.22617}
\genomerow{1243}{130}{19}{Terminal}{a = value}{5,4,4}{}{-38.1068}
\genomerow{1244}{130}{20}{Send Intra-Cell Message}{if a, send message another cardinal within cell}{6,3,6}{}{}
\genomerow{1245}{130}{21}{Terminate If}{if a, terminate core}{4,1,4}{}{}
\genomerow{1246}{130}{22}{Negate}{a = -a}{4,0,2}{}{}
\genomerow{1247}{130}{23}{Terminal}{a = value}{3,7,3}{}{-0.536215}
\genomerow{1248}{130}{24}{Adjust Secondary Global Regulator}{adjust global regulator value by a}{2,5,2}{}{}
\genomerow{1249}{130}{25}{Terminal}{a = value}{5,4,6}{}{-0.322304}
\genomerow{1250}{130}{26}{Not Equal}{a = (b != c)}{2,7,1}{}{}
\genomerow{1251}{130}{27}{Terminal}{a = value}{3,2,7}{}{1.30452}
\genomerow{1252}{130}{28}{Bitwise Shift}{a = b << c}{7,6,1}{}{}
\genomerow{1253}{130}{29}{Terminal}{a = value}{6,6,5}{}{9.74099}
\genomerow{1254}{131}{0}{Global Anchor}{register a global jump-to destination, maybe terminate}{2,7,7}{}{}
\genomerow{1255}{132}{0}{Global Anchor}{register a global jump-to destination, maybe terminate}{3,1,0}{}{}
\genomerow{1256}{132}{1}{Nop-0}{perform no operation}{5,5,3}{}{}
\genomerow{1257}{132}{2}{Terminal}{a = value}{3,5,2}{}{-11.152}
\genomerow{1258}{132}{3}{Terminal}{a = value}{4,2,7}{}{-38.1454}
\genomerow{1259}{132}{4}{Divide}{a = b / c}{6,4,1}{}{}
\genomerow{1260}{132}{5}{Local Jump If Not}{if !a, goto tag match local anchor}{1,0,2}{}{}
\genomerow{1261}{133}{0}{Global Anchor}{register a global jump-to destination, maybe terminate}{3,5,5}{}{}
\genomerow{1262}{133}{1}{Nop-0}{perform no operation}{2,5,2}{}{}
\genomerow{1263}{133}{2}{Decay Global Regulator}{decay global regulator value by register a}{6,7,0}{}{}
\genomerow{1264}{133}{3}{Local Anchor}{register a local jump-to destination}{5,6,2}{}{}
\genomerow{1265}{134}{0}{Global Anchor}{register a global jump-to destination, maybe terminate}{4,7,3}{}{}
\genomerow{1266}{134}{1}{Bitwise Not}{a = \textasciitilde{}b}{6,7,5}{}{}
\genomerow{1267}{134}{2}{Write Own State}{if b, state = a}{6,6,3}{HeirRequest[0]}{}
\genomerow{1268}{134}{3}{Terminal}{a = value}{0,1,4}{}{0.485842}
\genomerow{1269}{134}{4}{Read Neighbor State}{a = state}{0,5,5}{NeighborApoptosis[0]}{}
\genomerow{1270}{134}{5}{Read Neighbor State}{a = state}{2,7,0}{MostRecentCauseOfDeath[0]}{}
\genomerow{1271}{134}{6}{Read Neighbor State}{a = state}{0,5,1}{NopState[1]}{}
\genomerow{1272}{134}{7}{Read Own State}{a = state}{7,1,0}{ResourceSendLimit[0]}{}
\genomerow{1273}{134}{8}{Local Anchor}{register a local jump-to destination}{0,3,7}{}{}
\genomerow{1274}{134}{9}{Equal}{a = ( b == c )}{0,0,5}{}{}
\genomerow{1275}{134}{10}{Modulo}{a = b \% c}{4,7,6}{}{}
\genomerow{1276}{134}{11}{Local Anchor}{register a local jump-to destination}{0,3,4}{}{}
\genomerow{1277}{134}{12}{Multiply Own State}{state *= a}{0,0,4}{RepLevRequest[0]}{}
\genomerow{1278}{134}{13}{Add To Own State}{state += a}{3,7,5}{TransientNopState[1]}{}
\genomerow{1279}{135}{0}{Global Anchor}{register a global jump-to destination, maybe terminate}{7,6,3}{}{}
\genomerow{1280}{135}{1}{Read Neighbor State}{a = state}{1,1,7}{IsChildCellOf[0]}{}
\genomerow{1281}{135}{2}{Bitwise And}{a = b \& c}{6,5,0}{}{}
\genomerow{1282}{135}{3}{Terminal}{a = value}{4,4,0}{}{-0.907272}
\genomerow{1283}{136}{0}{Global Anchor}{register a global jump-to destination, maybe terminate}{6,4,1}{}{}
\genomerow{1284}{136}{1}{Get Local Regulator}{a = local regulator value}{6,7,0}{}{}
\genomerow{1285}{136}{2}{Get Local Regulator}{a = local regulator value}{6,6,0}{}{}
\genomerow{1286}{136}{3}{Terminal}{a = value}{4,5,5}{}{0.0279268}
\genomerow{1287}{136}{4}{Terminal}{a = value}{4,5,5}{}{0.0281404}
\genomerow{1288}{136}{5}{Bitwise And}{a = b \& c}{6,5,0}{}{}
\genomerow{1289}{136}{6}{Read Own State}{a = state}{1,0,4}{KinGroupAge[0]}{}
\genomerow{1290}{136}{7}{Bitwise And}{a = b \& c}{4,5,0}{}{}
\genomerow{1291}{137}{0}{Global Anchor}{register a global jump-to destination, maybe terminate}{7,4,2}{}{}
\genomerow{1292}{137}{1}{Get Local Regulator}{a = local regulator value}{6,5,4}{}{}
\genomerow{1293}{137}{2}{Get Local Regulator}{a = local regulator value}{6,6,0}{}{}
\genomerow{1294}{137}{3}{Terminal}{a = value}{4,5,5}{}{0.027408}
\genomerow{1295}{137}{4}{Bitwise And}{a = b \& c}{6,5,0}{}{}
\genomerow{1296}{137}{5}{Read Own State}{a = state}{1,0,4}{TransientNopState[2]}{}
\genomerow{1297}{137}{6}{Bitwise And}{a = b \& c}{4,5,0}{}{}
\genomerow{1298}{138}{0}{Global Anchor}{register a global jump-to destination, maybe terminate}{7,4,2}{}{}
\genomerow{1299}{138}{1}{Terminal}{a = value}{6,5,0}{}{0.560613}
\genomerow{1300}{138}{2}{Terminal}{a = value}{7,2,1}{}{-5.58893}
\genomerow{1301}{138}{3}{Bitwise Xor}{a = b \textasciicircum{} c}{2,1,5}{}{}
\genomerow{1302}{138}{4}{Less Than}{a = ( b < c )}{6,4,2}{}{}
\genomerow{1303}{138}{5}{Get Local Regulator}{a = local regulator value}{1,2,3}{}{}
\genomerow{1304}{138}{6}{Terminal}{a = value}{3,6,6}{}{-6.39391}
\genomerow{1305}{138}{7}{Set Local Regulator}{set local regulator value to a}{6,2,0}{}{}
\genomerow{1306}{138}{8}{Not Equal}{a = (b != c)}{4,3,7}{}{}
\genomerow{1307}{139}{0}{Global Anchor}{register a global jump-to destination, maybe terminate}{4,3,2}{}{}
\genomerow{1308}{139}{1}{Terminal}{a = value}{3,6,7}{}{1.00371}
\genomerow{1309}{139}{2}{Not Equal}{a = (b != c)}{5,3,1}{}{}
\genomerow{1310}{139}{3}{Terminal}{a = value}{4,2,7}{}{0.537759}
\genomerow{1311}{139}{4}{Get Global Regulator}{a = global regulator value}{0,2,5}{}{}
\genomerow{1312}{139}{5}{Generate Random Bool}{if p, a = 1; else a = 0}{6,5,3}{}{}
\genomerow{1313}{139}{6}{Get Global Regulator}{a = global regulator value}{7,5,7}{}{}
\genomerow{1314}{139}{7}{Divide}{a = b / c}{4,7,4}{}{}
\genomerow{1315}{139}{8}{Local Anchor}{register a local jump-to destination}{0,4,0}{}{}
\genomerow{1316}{139}{9}{Less Than}{a = ( b < c )}{4,2,2}{}{}
\genomerow{1317}{140}{0}{Global Anchor}{register a global jump-to destination, maybe terminate}{5,3,1}{}{}
\genomerow{1318}{140}{1}{Multiply}{a = b * c}{1,4,0}{}{}
\genomerow{1319}{140}{2}{Read Neighbor State}{a = state}{4,2,2}{KinGroupWillExpire[1]}{}
\genomerow{1320}{140}{3}{Local Anchor}{register a local jump-to destination}{5,6,5}{}{}
\genomerow{1321}{140}{4}{Divide}{a = b / c}{6,5,6}{}{}
\genomerow{1322}{141}{0}{Global Anchor}{register a global jump-to destination, maybe terminate}{0,1,2}{}{}
\genomerow{1323}{141}{1}{Logical And}{a = b \&\& c}{1,7,4}{}{}
\genomerow{1324}{141}{2}{Terminal}{a = value}{4,7,0}{}{181.55}
\genomerow{1325}{141}{3}{Local Anchor}{register a local jump-to destination}{3,6,1}{}{}
\genomerow{1326}{141}{4}{Global Jump If}{if a, goto tag match global anchor; if b, clear registers}{6,0,2}{}{}
\genomerow{1327}{141}{5}{Subtract}{a = b - c}{3,4,0}{}{}
\genomerow{1328}{141}{6}{Terminal}{a = value}{7,7,7}{}{-25191}
\genomerow{1329}{141}{7}{Read Own State}{a = state}{4,6,6}{SpawnedFrom[0]}{}
\genomerow{1330}{141}{8}{Read Own State}{a = state}{1,7,0}{KinGroupIDAncestorView[1]}{}
\genomerow{1331}{141}{9}{Multiply Own State}{state *= a}{5,4,1}{ResourceReserveRequest[0]}{}
\genomerow{1332}{141}{10}{Read Neighbor State}{a = state}{5,1,1}{IsParentCellOf[0]}{}
\genomerow{1333}{141}{11}{Bitwise And}{a = b \& c}{1,4,1}{}{}
\genomerow{1334}{141}{12}{Nop-2}{perform no operation}{1,3,7}{}{}
\genomerow{1335}{141}{13}{Terminal}{a = value}{5,4,2}{}{9.55006}
\genomerow{1336}{141}{14}{Get Global Regulator}{a = global regulator value}{6,4,7}{}{}
\genomerow{1337}{141}{15}{Send Inter-Cell Message}{if a, send message to neighbor cell}{1,2,2}{}{}
\genomerow{1338}{141}{16}{Send Intra-Cell Message}{if a, send message another cardinal within cell}{0,0,3}{}{}
\genomerow{1339}{141}{17}{Fork If}{if a, submit fork request (processed when core terminates)}{6,3,4}{}{}
\genomerow{1340}{141}{18}{Terminal}{a = value}{3,4,0}{}{13.261}
\genomerow{1341}{141}{19}{Bitwise Not}{a = \textasciitilde{}b}{5,5,2}{}{}
\genomerow{1342}{141}{20}{Get Global Regulator}{a = global regulator value}{3,5,7}{}{}
\genomerow{1343}{141}{21}{Broadcast Intra-Cell Message}{if a, send message to all other cardinals within cell}{3,5,6}{}{}
\genomerow{1344}{141}{22}{Terminal}{a = value}{2,2,2}{}{0.918583}
\genomerow{1345}{141}{23}{Decay Local Regulator}{decay local regulator value by register a}{5,2,0}{}{}
\genomerow{1346}{141}{24}{Read Own State}{a = state}{6,5,1}{KinGroupIDView[1]}{}
\genomerow{1347}{141}{25}{Set Local Regulator}{set local regulator value to a}{7,0,5}{}{}
\genomerow{1348}{141}{26}{Fork If}{if a, submit fork request (processed when core terminates)}{3,1,2}{}{}
\genomerow{1349}{141}{27}{Set Global Regulator}{set global regulator value to a}{2,3,5}{}{}
\genomerow{1350}{141}{28}{Logical And}{a = b \&\& c}{6,4,4}{}{}
\genomerow{1351}{141}{29}{Logical And}{a = b \&\& c}{6,4,4}{}{}
\genomerow{1352}{142}{0}{Global Anchor}{register a global jump-to destination, maybe terminate}{4,4,7}{}{}
\genomerow{1353}{142}{1}{Read Neighbor State}{a = state}{7,6,0}{KinGroupIDAncestorView[1]}{}
\genomerow{1354}{142}{2}{Read Own State}{a = state}{5,1,6}{OptimumQuorumExceeded[1]}{}
\genomerow{1355}{142}{3}{Local Anchor}{register a local jump-to destination}{1,0,6}{}{}
\genomerow{1356}{143}{0}{Global Anchor}{register a global jump-to destination, maybe terminate}{6,0,6}{}{}
\genomerow{1357}{143}{1}{Read Own State}{a = state}{5,6,1}{IsChildGroupOf[0]}{}
\genomerow{1358}{143}{2}{Read Own State}{a = state}{7,0,5}{IncomingInterMessageCounter[0]}{}
\genomerow{1359}{143}{3}{Terminal}{a = value}{3,2,3}{}{2.81134}
\genomerow{1360}{143}{4}{Less Than}{a = ( b < c )}{5,1,2}{}{}
\genomerow{1361}{143}{5}{Get Global Regulator}{a = global regulator value}{1,1,5}{}{}
\genomerow{1362}{144}{0}{Global Anchor}{register a global jump-to destination, maybe terminate}{7,4,2}{}{}
\genomerow{1363}{144}{1}{Local Anchor}{register a local jump-to destination}{1,2,1}{}{}
\genomerow{1364}{144}{2}{Local Anchor}{register a local jump-to destination}{6,1,3}{}{}
\genomerow{1365}{144}{3}{Bitwise Shift}{a = b << c}{6,3,3}{}{}
\genomerow{1366}{144}{4}{Decay Global Regulator}{decay global regulator value by register a}{6,7,6}{}{}
\genomerow{1367}{144}{5}{Broadcast Intra-Cell Message}{if a, send message to all other cardinals within cell}{4,1,0}{}{}
\genomerow{1368}{144}{6}{Not}{a = !a}{3,1,2}{}{}
\genomerow{1369}{144}{7}{Not Equal}{a = (b != c)}{5,1,2}{}{}
\genomerow{1370}{144}{8}{Terminal}{a = value}{4,3,3}{}{1.31859}
\genomerow{1371}{144}{9}{Not}{a = !a}{1,7,1}{}{}
\genomerow{1372}{144}{10}{Terminal}{a = value}{0,2,2}{}{2.08505}
\genomerow{1373}{144}{11}{Generate Random Bool}{if p, a = 1; else a = 0}{2,6,4}{}{}
\genomerow{1374}{144}{12}{Read Neighbor State}{a = state}{0,1,1}{NopState[0]}{}
\genomerow{1375}{144}{13}{Terminal}{a = value}{1,2,2}{}{0.160325}
\genomerow{1376}{145}{0}{Global Anchor}{register a global jump-to destination, maybe terminate}{6,0,6}{}{}
\genomerow{1377}{145}{1}{Terminal}{a = value}{5,6,3}{}{-9398.09}
\genomerow{1378}{145}{2}{Terminal}{a = value}{7,0,5}{}{-2.70517}
\genomerow{1379}{145}{3}{Decay Global Regulator}{decay global regulator value by register a}{6,7,6}{}{}
\genomerow{1380}{145}{4}{Send Inter-Cell Message}{if a, send message to neighbor cell}{4,1,0}{}{}
\genomerow{1381}{145}{5}{Terminal}{a = value}{6,1,2}{}{0.656863}
\genomerow{1382}{145}{6}{Draw Random Value}{a = weighted draw from random number generator}{5,3,0}{}{}
\genomerow{1383}{145}{7}{Terminal}{a = value}{1,3,2}{}{-4.3052}
\genomerow{1384}{145}{8}{Terminal}{a = value}{5,7,0}{}{0.320069}
\genomerow{1385}{145}{9}{Multiply}{a = b * c}{7,4,4}{}{}
\genomerow{1386}{145}{10}{Terminal}{a = value}{0,2,7}{}{-0.923173}
\genomerow{1387}{145}{11}{Terminal}{a = value}{3,7,7}{}{35.5229}
\genomerow{1388}{145}{12}{Local Anchor}{register a local jump-to destination}{2,0,3}{}{}
\genomerow{1389}{145}{13}{Terminal}{a = value}{3,7,3}{}{-0.682453}
\genomerow{1390}{145}{14}{Send Inter-Cell Message}{if a, send message to neighbor cell}{4,1,0}{}{}
\genomerow{1391}{145}{15}{Terminal}{a = value}{6,1,2}{}{0.656863}
\genomerow{1392}{145}{16}{Draw Random Value}{a = weighted draw from random number generator}{5,3,0}{}{}
\genomerow{1393}{145}{17}{Decay Secondary Global Regulator}{decay global regulator value by register a}{1,1,2}{}{}
\genomerow{1394}{145}{18}{Terminal}{a = value}{3,7,3}{}{-3.14618}
\genomerow{1395}{145}{19}{Terminal}{a = value}{2,2,5}{}{0.974192}
\genomerow{1396}{145}{20}{Terminal}{a = value}{2,2,5}{}{-0.0258076}
\genomerow{1397}{146}{0}{Global Anchor}{register a global jump-to destination, maybe terminate}{0,2,3}{}{}
\genomerow{1398}{146}{1}{Terminal}{a = value}{4,1,5}{}{72.1054}
\genomerow{1399}{146}{2}{Logical Or}{a = b \textbar{}\textbar{} c}{4,6,6}{}{}
\genomerow{1400}{146}{3}{Equal}{a = ( b == c )}{4,2,6}{}{}
\genomerow{1401}{146}{4}{Terminal}{a = value}{1,5,1}{}{-0.619654}
\genomerow{1402}{146}{5}{Terminal}{a = value}{4,1,5}{}{75.7399}
\genomerow{1403}{146}{6}{Logical Or}{a = b \textbar{}\textbar{} c}{4,6,6}{}{}
\genomerow{1404}{146}{7}{Equal}{a = ( b == c )}{4,2,6}{}{}
\genomerow{1405}{146}{8}{Terminal}{a = value}{1,5,1}{}{-0.619654}
\genomerow{1406}{146}{9}{Decay Secondary Global Regulator}{decay global regulator value by register a}{7,4,6}{}{}
\genomerow{1407}{146}{10}{Local Anchor}{register a local jump-to destination}{0,7,7}{}{}
\genomerow{1408}{146}{11}{Terminal}{a = value}{6,6,0}{}{-9.41647}
\genomerow{1409}{146}{12}{Add}{a = b + c}{4,3,4}{}{}
\genomerow{1410}{146}{13}{Read Own State}{a = state}{5,0,1}{ResourceReserveRequest[0]}{}
\genomerow{1411}{146}{14}{Get Local Regulator}{a = local regulator value}{3,3,3}{}{}
\genomerow{1412}{146}{15}{Subtract}{a = b - c}{1,0,0}{}{}
\genomerow{1413}{146}{16}{Not Equal}{a = (b != c)}{1,1,5}{}{}
\genomerow{1414}{147}{0}{Global Anchor}{register a global jump-to destination, maybe terminate}{1,4,7}{}{}
\genomerow{1415}{147}{1}{Set Secondary Global Regulator}{set global regulator value to a}{6,7,0}{}{}
\genomerow{1416}{147}{2}{Nop-0}{perform no operation}{4,5,5}{}{}
\genomerow{1417}{147}{3}{Adjust Global Regulator}{adjust global regulator value by a}{0,1,7}{}{}
\genomerow{1418}{147}{4}{Modulo}{a = b \% c}{6,3,3}{}{}
\genomerow{1419}{147}{5}{Terminal}{a = value}{3,6,6}{}{-2.83514}
\genomerow{1420}{147}{6}{Bitwise Xor}{a = b \textasciicircum{} c}{3,6,1}{}{}
\genomerow{1421}{147}{7}{Add}{a = b + c}{5,0,6}{}{}
\genomerow{1422}{148}{0}{Global Anchor}{register a global jump-to destination, maybe terminate}{6,1,6}{}{}
\genomerow{1423}{149}{0}{Global Anchor}{register a global jump-to destination, maybe terminate}{2,4,5}{}{}
\genomerow{1424}{149}{1}{Local Anchor}{register a local jump-to destination}{7,1,1}{}{}
\genomerow{1425}{149}{2}{Local Anchor}{register a local jump-to destination}{5,5,0}{}{}
\genomerow{1426}{149}{3}{Local Anchor}{register a local jump-to destination}{5,3,4}{}{}
\genomerow{1427}{150}{0}{Global Anchor}{register a global jump-to destination, maybe terminate}{5,0,4}{}{}
\genomerow{1428}{150}{1}{Terminal}{a = value}{5,5,4}{}{-0.402138}
\genomerow{1429}{150}{2}{Local Anchor}{register a local jump-to destination}{3,5,1}{}{}
\genomerow{1430}{150}{3}{Multiply Own State}{state *= a}{7,2,3}{ResourceReceiveResistance[0]}{}
\genomerow{1431}{150}{4}{Read Own State}{a = state}{1,4,1}{SpawnCount[0]}{}
\genomerow{1432}{150}{5}{Decay Local Regulator}{decay local regulator value by register a}{2,2,7}{}{}
\genomerow{1433}{150}{6}{Subtract}{a = b - c}{1,1,5}{}{}
\genomerow{1434}{150}{7}{Terminal}{a = value}{6,3,6}{}{-1.10395e+06}
\genomerow{1435}{151}{0}{Global Anchor}{register a global jump-to destination, maybe terminate}{2,5,2}{}{}
\genomerow{1436}{151}{1}{Terminal}{a = value}{3,0,0}{}{0.41467}
\genomerow{1437}{151}{2}{Local Anchor}{register a local jump-to destination}{5,7,1}{}{}
\genomerow{1438}{151}{3}{Read Own State}{a = state}{7,2,3}{NopState[1]}{}
\genomerow{1439}{151}{4}{Logical And}{a = b \&\& c}{5,7,3}{}{}
\genomerow{1440}{151}{5}{Count Ones}{a = popcnt( b )}{0,6,6}{}{}
\genomerow{1441}{151}{6}{Local Anchor}{register a local jump-to destination}{4,2,1}{}{}
\genomerow{1442}{151}{7}{Terminal}{a = value}{7,6,3}{}{-14.6687}
\genomerow{1443}{151}{8}{Local Anchor}{register a local jump-to destination}{5,7,1}{}{}
\genomerow{1444}{151}{9}{Read Own State}{a = state}{7,2,3}{NopState[1]}{}
\genomerow{1445}{151}{10}{Terminate If}{if a, terminate core}{5,7,3}{}{}
\genomerow{1446}{151}{11}{Count Ones}{a = popcnt( b )}{0,6,6}{}{}
\genomerow{1447}{151}{12}{Local Anchor}{register a local jump-to destination}{4,2,1}{}{}
\genomerow{1448}{151}{13}{Terminal}{a = value}{7,6,3}{}{-14.6687}
\genomerow{1449}{151}{14}{Terminal}{a = value}{6,5,4}{}{-4.23025}
\genomerow{1450}{151}{15}{Terminal}{a = value}{7,6,3}{}{-14.6687}
\genomerow{1451}{151}{16}{Terminal}{a = value}{6,5,4}{}{-4.23025}
\genomerow{1452}{151}{17}{Decrement}{a = a - 1}{3,1,4}{}{}
\genomerow{1453}{151}{18}{Logical Or}{a = b \textbar{}\textbar{} c}{5,5,1}{}{}
\genomerow{1454}{151}{19}{Read Neighbor State}{a = state}{6,0,4}{NopState[1]}{}
\genomerow{1455}{151}{20}{Decay Global Regulator}{decay global regulator value by register a}{1,6,1}{}{}
\genomerow{1456}{151}{21}{Read Neighbor State}{a = state}{1,0,1}{NumKnownQuorumBits[1]}{}
\genomerow{1457}{152}{0}{Global Anchor}{register a global jump-to destination, maybe terminate}{1,0,1}{}{}
\genomerow{1458}{152}{1}{Set Local Regulator}{set local regulator value to a}{7,2,2}{}{}
\genomerow{1459}{152}{2}{Decrement}{a = a - 1}{6,2,1}{}{}
\genomerow{1460}{152}{3}{Decay Global Regulator}{decay global regulator value by register a}{1,6,1}{}{}
\genomerow{1461}{152}{4}{Read Neighbor State}{a = state}{1,0,1}{NumKnownQuorumBits[1]}{}
\genomerow{1462}{152}{5}{Read Neighbor State}{a = state}{1,0,1}{SpawnCount[0]}{}
\genomerow{1463}{152}{6}{Set Local Regulator}{set local regulator value to a}{7,2,2}{}{}
\genomerow{1464}{152}{7}{Decrement}{a = a - 1}{6,2,1}{}{}
\genomerow{1465}{152}{8}{Read Neighbor State}{a = state}{1,5,0}{HeirRequest[0]}{}
\genomerow{1466}{152}{9}{Local Anchor}{register a local jump-to destination}{0,2,5}{}{}
\genomerow{1467}{152}{10}{Read Own State}{a = state}{0,6,5}{NeighborOptimumQuorumExceeded[1]}{}
\genomerow{1468}{152}{11}{Read Own State}{a = state}{0,6,3}{NeighborApoptosis[0]}{}
\genomerow{1469}{152}{12}{Set Local Regulator}{set local regulator value to a}{3,5,4}{}{}
\genomerow{1470}{153}{0}{Global Anchor}{register a global jump-to destination, maybe terminate}{7,5,7}{}{}
\genomerow{1471}{153}{1}{Local Anchor}{register a local jump-to destination}{3,6,0}{}{}
\genomerow{1472}{153}{2}{Terminal}{a = value}{7,1,1}{}{-3.49294}
\genomerow{1473}{153}{3}{Bitwise Xor}{a = b \textasciicircum{} c}{5,3,0}{}{}
\genomerow{1474}{153}{4}{Fork If}{if a, submit fork request (processed when core terminates)}{1,5,3}{}{}
\genomerow{1475}{153}{5}{Local Anchor}{register a local jump-to destination}{0,2,5}{}{}
\genomerow{1476}{153}{6}{Terminal}{a = value}{1,6,5}{}{4.26208}
\genomerow{1477}{153}{7}{Set Local Regulator}{set local regulator value to a}{3,5,4}{}{}
\genomerow{1478}{154}{0}{Global Anchor}{register a global jump-to destination, maybe terminate}{7,5,7}{}{}
\genomerow{1479}{154}{1}{Local Anchor}{register a local jump-to destination}{3,6,0}{}{}
\genomerow{1480}{154}{2}{Terminal}{a = value}{7,1,1}{}{-3.49294}
\genomerow{1481}{154}{3}{Bitwise Xor}{a = b \textasciicircum{} c}{5,3,0}{}{}
\genomerow{1482}{154}{4}{Local Jump If}{if a, goto tag match local anchor}{4,5,3}{}{}
\genomerow{1483}{154}{5}{Multiply Own State}{state *= a}{5,1,1}{RepLevRequest[0]}{}
\genomerow{1484}{154}{6}{Decay Global Regulator}{decay global regulator value by register a}{6,6,1}{}{}
\genomerow{1485}{154}{7}{Terminal}{a = value}{0,5,5}{}{-5706.19}
\genomerow{1486}{154}{8}{Less Than}{a = ( b < c )}{6,2,3}{}{}
\genomerow{1487}{154}{9}{Terminal}{a = value}{5,7,6}{}{7606.48}
\genomerow{1488}{154}{10}{Set Local Regulator}{set local regulator value to a}{1,2,3}{}{}
\genomerow{1489}{154}{11}{Fork If}{if a, submit fork request (processed when core terminates)}{1,5,3}{}{}
\genomerow{1490}{154}{12}{Local Anchor}{register a local jump-to destination}{0,2,5}{}{}
\genomerow{1491}{154}{13}{Terminal}{a = value}{1,6,5}{}{4.26208}
\genomerow{1492}{154}{14}{Fork If}{if a, submit fork request (processed when core terminates)}{1,5,3}{}{}
\genomerow{1493}{154}{15}{Add To Own State}{state += a}{2,2,5}{ResourceSendRequest[0]}{}
\genomerow{1494}{154}{16}{Divide}{a = b / c}{3,6,2}{}{}
\genomerow{1495}{154}{17}{Get Local Regulator}{a = local regulator value}{1,1,2}{}{}
\genomerow{1496}{154}{18}{Terminate If}{if a, terminate core}{0,6,6}{}{}
\genomerow{1497}{154}{19}{Local Anchor}{register a local jump-to destination}{2,5,1}{}{}
\genomerow{1498}{154}{20}{Read Neighbor State}{a = state}{4,0,0}{ResourceSendRequest[0]}{}
\genomerow{1499}{154}{21}{Decay Secondary Global Regulator}{decay global regulator value by register a}{1,7,5}{}{}
\genomerow{1500}{155}{0}{Global Anchor}{register a global jump-to destination, maybe terminate}{3,5,3}{}{}
\genomerow{1501}{155}{1}{Broadcast Intra-Cell Message}{if a, send message to all other cardinals within cell}{7,7,4}{}{}
\genomerow{1502}{155}{2}{Modulo}{a = b \% c}{2,1,6}{}{}
\genomerow{1503}{155}{3}{Terminal}{a = value}{4,6,7}{}{0.527316}
\genomerow{1504}{155}{4}{Terminal}{a = value}{4,1,5}{}{-0.773296}
\genomerow{1505}{155}{5}{Fork If}{if a, submit fork request (processed when core terminates)}{5,1,7}{}{}
\genomerow{1506}{155}{6}{Less Than}{a = ( b < c )}{7,1,2}{}{}
\genomerow{1507}{155}{7}{Write Own State}{if b, state = a}{0,2,2}{ResourceSendLimit[0]}{}
\genomerow{1508}{155}{8}{Terminal}{a = value}{5,4,5}{}{0.8735}
\genomerow{1509}{155}{9}{Bitwise Xor}{a = b \textasciicircum{} c}{0,7,3}{}{}
\genomerow{1510}{155}{10}{Nop-1}{perform no operation}{6,0,4}{}{}
\genomerow{1511}{155}{11}{Add To Own State}{state += a}{2,3,3}{TransientNopState[3]}{}
\genomerow{1512}{155}{12}{Not}{a = !a}{0,1,3}{}{}
\genomerow{1513}{155}{13}{Decay Secondary Global Regulator}{decay global regulator value by register a}{7,0,2}{}{}
\genomerow{1514}{156}{0}{Global Anchor}{register a global jump-to destination, maybe terminate}{2,4,7}{}{}
\genomerow{1515}{156}{1}{Terminal}{a = value}{1,3,3}{}{1.69959}
\genomerow{1516}{156}{2}{Send Inter-Cell Message}{if a, send message to neighbor cell}{5,3,7}{}{}
\genomerow{1517}{156}{3}{Read Own State}{a = state}{5,1,5}{RicherThanNeighbor[0]}{}
\genomerow{1518}{156}{4}{Local Jump If Not}{if !a, goto tag match local anchor}{5,0,3}{}{}
\genomerow{1519}{156}{5}{Get Global Regulator}{a = global regulator value}{0,2,1}{}{}
\genomerow{1520}{156}{6}{Draw Random Value}{a = weighted draw from random number generator}{6,5,5}{}{}
\genomerow{1521}{156}{7}{Terminal}{a = value}{4,4,5}{}{-0.154192}
\genomerow{1522}{156}{8}{Terminal}{a = value}{4,4,5}{}{-0.154192}
\genomerow{1523}{156}{9}{Bitwise Not}{a = \textasciitilde{}b}{4,6,4}{}{}
\genomerow{1524}{156}{10}{Modulo}{a = b \% c}{2,3,6}{}{}
\genomerow{1525}{156}{11}{Get Global Regulator}{a = global regulator value}{1,5,4}{}{}
\genomerow{1526}{156}{12}{Terminal}{a = value}{4,6,7}{}{-0.559602}
\genomerow{1527}{156}{13}{Terminal}{a = value}{4,4,5}{}{-0.154192}
\genomerow{1528}{156}{14}{Bitwise Not}{a = \textasciitilde{}b}{4,6,4}{}{}
\genomerow{1529}{156}{15}{Modulo}{a = b \% c}{2,3,6}{}{}
\genomerow{1530}{156}{16}{Get Global Regulator}{a = global regulator value}{1,5,4}{}{}
\genomerow{1531}{156}{17}{Terminal}{a = value}{4,6,7}{}{-0.56009}
\genomerow{1532}{156}{18}{RandomFill}{fill a with uniformly-distributed random bits}{3,3,7}{}{}
\genomerow{1533}{156}{19}{Nop-0}{perform no operation}{7,0,3}{}{}
\genomerow{1534}{156}{20}{Count Ones}{a = popcnt( b )}{4,2,2}{}{}
\genomerow{1535}{156}{21}{Logical Or}{a = b \textbar{}\textbar{} c}{7,1,0}{}{}
\genomerow{1536}{156}{22}{Terminal}{a = value}{4,2,1}{}{-15.3926}
\genomerow{1537}{156}{23}{Terminate If}{if a, terminate core}{4,4,6}{}{}
\genomerow{1538}{156}{24}{Terminal}{a = value}{2,5,5}{}{0.912099}
\genomerow{1539}{156}{25}{Less Than}{a = ( b < c )}{7,5,1}{}{}
\genomerow{1540}{156}{26}{Local Anchor}{register a local jump-to destination}{1,3,2}{}{}
\genomerow{1541}{156}{27}{Get Global Regulator}{a = global regulator value}{6,4,0}{}{}
\genomerow{1542}{156}{28}{Send Inter-Cell Message}{if a, send message to neighbor cell}{7,4,2}{}{}
\genomerow{1543}{156}{29}{Equal}{a = ( b == c )}{2,6,3}{}{}
\genomerow{1544}{156}{30}{Terminate If}{if a, terminate core}{4,2,5}{}{}
\genomerow{1545}{156}{31}{Divide}{a = b / c}{5,6,5}{}{}
\genomerow{1546}{156}{32}{RandomFill}{fill a with uniformly-distributed random bits}{5,1,3}{}{}
\genomerow{1547}{156}{33}{Decrement}{a = a - 1}{3,5,2}{}{}
\genomerow{1548}{156}{34}{Less Than}{a = ( b < c )}{6,0,2}{}{}
\genomerow{1549}{156}{35}{Local Anchor}{register a local jump-to destination}{4,1,6}{}{}
\genomerow{1550}{156}{36}{Local Anchor}{register a local jump-to destination}{3,0,4}{}{}
\genomerow{1551}{156}{37}{Decrement}{a = a - 1}{3,1,6}{}{}
\genomerow{1552}{156}{38}{Modulo}{a = b \% c}{7,0,7}{}{}
\genomerow{1553}{157}{0}{Global Anchor}{register a global jump-to destination, maybe terminate}{2,2,6}{}{}
\genomerow{1554}{157}{1}{Local Anchor}{register a local jump-to destination}{7,4,4}{}{}
\genomerow{1555}{158}{0}{Global Anchor}{register a global jump-to destination, maybe terminate}{0,0,3}{}{}
\genomerow{1556}{158}{1}{Add To Own State}{state += a}{0,7,2}{NopState[1]}{}
\genomerow{1557}{158}{2}{Add To Own State}{state += a}{4,0,3}{ResourceSendLimit[0]}{}
\genomerow{1558}{158}{3}{RandomFill}{fill a with uniformly-distributed random bits}{6,0,6}{}{}
\genomerow{1559}{158}{4}{Read Own State}{a = state}{6,1,6}{IsChildCellOf[0]}{}
\genomerow{1560}{158}{5}{Send Intra-Cell Message}{if a, send message another cardinal within cell}{4,7,1}{}{}
\genomerow{1561}{159}{0}{Global Anchor}{register a global jump-to destination, maybe terminate}{1,2,0}{}{}
\genomerow{1562}{159}{1}{Terminal}{a = value}{5,0,0}{}{-0.548841}
\genomerow{1563}{159}{2}{Terminal}{a = value}{1,5,1}{}{-33.4397}
\genomerow{1564}{159}{3}{Adjust Local Regulator}{adjust local regulator value by a}{4,3,2}{}{}
\genomerow{1565}{159}{4}{Broadcast Intra-Cell Message}{if a, send message to all other cardinals within cell}{3,7,6}{}{}
\genomerow{1566}{160}{0}{Global Anchor}{register a global jump-to destination, maybe terminate}{5,5,5}{}{}
\genomerow{1567}{160}{1}{Add To Own State}{state += a}{6,0,6}{ResourceSendRequest[0]}{}
\genomerow{1568}{160}{2}{Read Neighbor State}{a = state}{7,7,3}{NumKnownQuorumBits[1]}{}
\genomerow{1569}{160}{3}{Less Than}{a = ( b < c )}{4,0,4}{}{}
\genomerow{1570}{160}{4}{Terminal}{a = value}{7,4,2}{}{-1802.75}
\genomerow{1571}{160}{5}{Send Intra-Cell Message}{if a, send message another cardinal within cell}{6,7,1}{}{}
\genomerow{1572}{161}{0}{Global Anchor}{register a global jump-to destination, maybe terminate}{7,4,0}{}{}
\genomerow{1573}{161}{1}{Add To Own State}{state += a}{6,0,6}{ResourceSendRequest[0]}{}
\genomerow{1574}{161}{2}{Read Neighbor State}{a = state}{7,7,3}{NumKnownQuorumBits[1]}{}
\genomerow{1575}{161}{3}{Less Than}{a = ( b < c )}{4,0,4}{}{}
\genomerow{1576}{161}{4}{Terminal}{a = value}{7,4,2}{}{-1802.75}
\genomerow{1577}{161}{5}{Send Intra-Cell Message}{if a, send message another cardinal within cell}{6,7,1}{}{}
\genomerow{1578}{162}{0}{Global Anchor}{register a global jump-to destination, maybe terminate}{7,4,1}{}{}
\genomerow{1579}{162}{1}{Read Neighbor State}{a = state}{6,1,0}{IncomingInterMessageCounter[0]}{}
\genomerow{1580}{162}{2}{Write Own State}{if b, state = a}{1,2,3}{RepLevRequest[0]}{}
\genomerow{1581}{163}{0}{Global Anchor}{register a global jump-to destination, maybe terminate}{7,4,0}{}{}
\genomerow{1582}{163}{1}{Read Neighbor State}{a = state}{6,1,0}{IncomingInterMessageCounter[0]}{}
\genomerow{1583}{163}{2}{Write Own State}{if b, state = a}{1,2,3}{RepLevRequest[0]}{}
\genomerow{1584}{163}{3}{Decay Secondary Global Regulator}{decay global regulator value by register a}{4,7,1}{}{}
\genomerow{1585}{164}{0}{Global Anchor}{register a global jump-to destination, maybe terminate}{1,5,5}{}{}
\genomerow{1586}{164}{1}{Send Intra-Cell Message}{if a, send message another cardinal within cell}{7,7,3}{}{}
\genomerow{1587}{164}{2}{Less Than}{a = ( b < c )}{4,0,0}{}{}
\genomerow{1588}{164}{3}{Terminal}{a = value}{7,4,4}{}{-328.623}
\genomerow{1589}{164}{4}{Less Than}{a = ( b < c )}{4,0,0}{}{}
\genomerow{1590}{164}{5}{Terminal}{a = value}{7,4,2}{}{-356.699}
\genomerow{1591}{164}{6}{Terminal}{a = value}{1,5,5}{}{-1299.58}
\genomerow{1592}{164}{7}{Logical And}{a = b \&\& c}{6,1,2}{}{}
\genomerow{1593}{164}{8}{Set Local Regulator}{set local regulator value to a}{7,7,3}{}{}
\genomerow{1594}{164}{9}{Decay Secondary Global Regulator}{decay global regulator value by register a}{4,3,5}{}{}
\genomerow{1595}{164}{10}{Write Own State}{if b, state = a}{6,6,4}{NopState[1]}{}
\genomerow{1596}{164}{11}{Multiply}{a = b * c}{6,3,5}{}{}
\genomerow{1597}{164}{12}{Read Own State}{a = state}{7,7,5}{KinGroupIDView[0]}{}
\genomerow{1598}{164}{13}{Bitwise Or}{a = b \textbar{} c}{7,5,0}{}{}
\genomerow{1599}{164}{14}{Local Jump If Not}{if !a, goto tag match local anchor}{4,0,5}{}{}
\genomerow{1600}{165}{0}{Global Anchor}{register a global jump-to destination, maybe terminate}{6,1,2}{}{}
\genomerow{1601}{165}{1}{Decrement}{a = a - 1}{5,7,2}{}{}
\genomerow{1602}{165}{2}{Terminal}{a = value}{5,7,2}{}{3.05429}
\genomerow{1603}{165}{3}{Terminal}{a = value}{7,7,2}{}{2.65213}
\genomerow{1604}{166}{0}{Global Anchor}{register a global jump-to destination, maybe terminate}{7,0,0}{}{}
\genomerow{1605}{166}{1}{Set Local Regulator}{set local regulator value to a}{7,7,3}{}{}
\genomerow{1606}{166}{2}{Terminal}{a = value}{3,7,2}{}{2.65213}
\genomerow{1607}{167}{0}{Global Anchor}{register a global jump-to destination, maybe terminate}{7,0,0}{}{}
\genomerow{1608}{167}{1}{Read Own State}{a = state}{6,1,2}{MostRecentCauseOfDeath[0]}{}
\genomerow{1609}{167}{2}{Decrement}{a = a - 1}{5,7,2}{}{}
\genomerow{1610}{167}{3}{Terminal}{a = value}{7,7,2}{}{2.67385}
\genomerow{1611}{167}{4}{Global Jump If Not}{if !a, goto global anchor; if !b, clear registers}{7,5,1}{}{}
\genomerow{1612}{167}{5}{Send Intra-Cell Message}{if a, send message another cardinal within cell}{3,6,6}{}{}
\genomerow{1613}{167}{6}{Global Jump If Not}{if !a, goto global anchor; if !b, clear registers}{7,5,1}{}{}
\genomerow{1614}{167}{7}{Send Intra-Cell Message}{if a, send message another cardinal within cell}{3,6,6}{}{}
\genomerow{1615}{167}{8}{Read Own State}{a = state}{6,1,2}{NeighborFragmented[0]}{}
\genomerow{1616}{167}{9}{Terminal}{a = value}{5,7,2}{}{0.449612}
\genomerow{1617}{167}{10}{Terminal}{a = value}{7,7,2}{}{2.67385}
\genomerow{1618}{168}{0}{Global Anchor}{register a global jump-to destination, maybe terminate}{4,2,4}{}{}
\genomerow{1619}{168}{1}{Terminal}{a = value}{6,1,6}{}{0.265916}
\genomerow{1620}{169}{0}{Global Anchor}{register a global jump-to destination, maybe terminate}{4,2,0}{}{}
\genomerow{1621}{170}{0}{Global Anchor}{register a global jump-to destination, maybe terminate}{4,2,4}{}{}
\genomerow{1622}{170}{1}{Terminal}{a = value}{6,1,6}{}{53.1825}
\genomerow{1623}{171}{0}{Global Anchor}{register a global jump-to destination, maybe terminate}{4,2,0}{}{}
\genomerow{1624}{172}{0}{Global Anchor}{register a global jump-to destination, maybe terminate}{4,2,4}{}{}
\genomerow{1625}{172}{1}{Decay Secondary Global Regulator}{decay global regulator value by register a}{6,1,6}{}{}
\genomerow{1626}{173}{0}{Global Anchor}{register a global jump-to destination, maybe terminate}{4,2,0}{}{}
\genomerow{1627}{174}{0}{Global Anchor}{register a global jump-to destination, maybe terminate}{7,0,0}{}{}
\genomerow{1628}{175}{0}{Global Anchor}{register a global jump-to destination, maybe terminate}{6,0,0}{}{}
\genomerow{1629}{175}{1}{Read Own State}{a = state}{6,1,2}{NeighborFragmented[0]}{}
\genomerow{1630}{175}{2}{Decrement}{a = a - 1}{1,7,2}{}{}
\genomerow{1631}{176}{0}{Global Anchor}{register a global jump-to destination, maybe terminate}{5,6,0}{}{}
\genomerow{1632}{176}{1}{Bitwise Shift}{a = b << c}{4,2,5}{}{}
\genomerow{1633}{176}{2}{Read Own State}{a = state}{0,4,3}{ReceivedResourceFrom[0]}{}
\genomerow{1634}{176}{3}{Read Own State}{a = state}{3,6,2}{StockpileDepleted[0]}{}
\genomerow{1635}{176}{4}{Terminal}{a = value}{1,6,2}{}{-1.02362}
\genomerow{1636}{176}{5}{Read Neighbor State}{a = state}{4,0,2}{IsChildGroupOf[1]}{}
\genomerow{1637}{176}{6}{Equal}{a = ( b == c )}{5,5,5}{}{}
\genomerow{1638}{176}{7}{Terminal}{a = value}{3,6,5}{}{-1.70521}
\genomerow{1639}{176}{8}{Get Secondary Global Regulator}{a = global regulator value}{5,6,5}{}{}
\genomerow{1640}{176}{9}{Terminal}{a = value}{2,6,1}{}{-0.950089}
\genomerow{1641}{176}{10}{Local Anchor}{register a local jump-to destination}{5,6,0}{}{}
\genomerow{1642}{176}{11}{Bitwise Shift}{a = b << c}{4,2,1}{}{}
\genomerow{1643}{176}{12}{Read Own State}{a = state}{0,0,3}{NopState[2]}{}
\genomerow{1644}{176}{13}{Bitwise Shift}{a = b << c}{4,2,5}{}{}
\genomerow{1645}{176}{14}{Read Own State}{a = state}{0,4,3}{ReceivedResourceFrom[0]}{}
\genomerow{1646}{176}{15}{Read Own State}{a = state}{3,6,2}{StockpileDepleted[0]}{}
\genomerow{1647}{176}{16}{Get Secondary Global Regulator}{a = global regulator value}{5,6,5}{}{}
\genomerow{1648}{176}{17}{Add To Own State}{state += a}{6,0,1}{NopState[0]}{}
\genomerow{1649}{176}{18}{Add To Own State}{state += a}{6,0,0}{NopState[0]}{}
\genomerow{1650}{176}{19}{Bitwise Shift}{a = b << c}{6,4,5}{}{}
\genomerow{1651}{176}{20}{Read Own State}{a = state}{4,6,3}{StockpileFecund[0]}{}
\genomerow{1652}{176}{21}{Add To Own State}{state += a}{6,0,0}{NopState[2]}{}
\genomerow{1653}{176}{22}{Decrement}{a = a - 1}{5,4,6}{}{}
\genomerow{1654}{176}{23}{Terminal}{a = value}{3,6,2}{}{3.04407}
\genomerow{1655}{176}{24}{Bitwise Shift}{a = b << c}{5,5,7}{}{}
\genomerow{1656}{176}{25}{Decrement}{a = a - 1}{5,6,6}{}{}
\genomerow{1657}{176}{26}{Terminal}{a = value}{7,3,7}{}{0.945789}
\genomerow{1658}{176}{27}{Send Intra-Cell Message}{if a, send message another cardinal within cell}{7,5,5}{}{}
\genomerow{1659}{176}{28}{Terminal}{a = value}{7,4,1}{}{0.825638}
\genomerow{1660}{176}{29}{Local Anchor}{register a local jump-to destination}{1,3,2}{}{}
\genomerow{1661}{176}{30}{Read Own State}{a = state}{4,4,5}{IsParentGroupOf[0]}{}
\genomerow{1662}{176}{31}{Multiply Own State}{state *= a}{5,3,7}{SpawnRequest[0]}{}
\genomerow{1663}{176}{32}{Terminal}{a = value}{2,0,7}{}{3.35585}
\genomerow{1664}{176}{33}{Not Equal}{a = (b != c)}{6,1,2}{}{}
\genomerow{1665}{176}{34}{Local Anchor}{register a local jump-to destination}{5,0,7}{}{}
\genomerow{1666}{176}{35}{Send Intra-Cell Message}{if a, send message another cardinal within cell}{7,0,5}{}{}
\genomerow{1667}{176}{36}{Decay Secondary Global Regulator}{decay global regulator value by register a}{2,3,7}{}{}
\genomerow{1668}{176}{37}{Draw Random Value}{a = weighted draw from random number generator}{2,3,7}{}{}
\genomerow{1669}{176}{38}{Decrement}{a = a - 1}{3,2,3}{}{}
\genomerow{1670}{176}{39}{Terminal}{a = value}{2,1,4}{}{9.54632}
\genomerow{1671}{176}{40}{Bitwise And}{a = b \& c}{0,6,0}{}{}
\genomerow{1672}{177}{0}{Global Anchor}{register a global jump-to destination, maybe terminate}{6,1,4}{}{}
\genomerow{1673}{177}{1}{Bitwise And}{a = b \& c}{4,0,4}{}{}
\genomerow{1674}{177}{2}{Terminal}{a = value}{6,1,6}{}{-0.858233}
\genomerow{1675}{177}{3}{Decay Secondary Global Regulator}{decay global regulator value by register a}{2,1,4}{}{}
\genomerow{1676}{177}{4}{Bitwise And}{a = b \& c}{0,6,4}{}{}
\genomerow{1677}{177}{5}{Draw Random Value}{a = weighted draw from random number generator}{6,1,5}{}{}
\genomerow{1678}{177}{6}{Bitwise And}{a = b \& c}{4,2,4}{}{}
\genomerow{1679}{177}{7}{Decay Secondary Global Regulator}{decay global regulator value by register a}{1,4,7}{}{}
\genomerow{1680}{177}{8}{Bitwise And}{a = b \& c}{4,0,4}{}{}
\genomerow{1681}{177}{9}{Terminal}{a = value}{6,1,6}{}{-0.858233}
\genomerow{1682}{177}{10}{Decay Secondary Global Regulator}{decay global regulator value by register a}{2,1,4}{}{}
\genomerow{1683}{177}{11}{Bitwise And}{a = b \& c}{0,6,6}{}{}
\genomerow{1684}{177}{12}{Set Secondary Global Regulator}{set global regulator value to a}{6,1,5}{}{}
\genomerow{1685}{177}{13}{Bitwise And}{a = b \& c}{4,2,4}{}{}
\genomerow{1686}{177}{14}{Terminal}{a = value}{1,4,7}{}{0.290945}
\genomerow{1687}{177}{15}{Read Own State}{a = state}{5,5,4}{NeighborFragmented[0]}{}
\genomerow{1688}{177}{16}{Decay Secondary Global Regulator}{decay global regulator value by register a}{6,1,6}{}{}
\genomerow{1689}{177}{17}{Terminal}{a = value}{5,1,0}{}{-2.20725}
\genomerow{1690}{177}{18}{Read Own State}{a = state}{0,2,1}{NumKnownQuorumBits[1]}{}
\genomerow{1691}{177}{19}{Terminal}{a = value}{7,7,6}{}{-0.241979}
\genomerow{1692}{178}{0}{Global Anchor}{register a global jump-to destination, maybe terminate}{1,3,0}{}{}
\genomerow{1693}{178}{1}{Terminal}{a = value}{1,2,0}{}{0.408177}
\genomerow{1694}{178}{2}{Less Than}{a = ( b < c )}{2,6,2}{}{}
\genomerow{1695}{178}{3}{Local Anchor}{register a local jump-to destination}{4,7,4}{}{}
\genomerow{1696}{178}{4}{Send Inter-Cell Message}{if a, send message to neighbor cell}{5,4,3}{}{}
\genomerow{1697}{178}{5}{Decay Local Regulator}{decay local regulator value by register a}{3,7,0}{}{}
\genomerow{1698}{178}{6}{Terminal}{a = value}{7,0,5}{}{0.882176}
\genomerow{1699}{178}{7}{Nop-2}{perform no operation}{1,5,4}{}{}
\genomerow{1700}{178}{8}{Local Jump If}{if a, goto tag match local anchor}{2,6,2}{}{}
\genomerow{1701}{178}{9}{Terminal}{a = value}{1,2,0}{}{0.470562}
\genomerow{1702}{179}{0}{Global Anchor}{register a global jump-to destination, maybe terminate}{2,6,2}{}{}
\genomerow{1703}{179}{1}{Local Anchor}{register a local jump-to destination}{4,5,0}{}{}
\genomerow{1704}{179}{2}{Read Neighbor State}{a = state}{5,4,3}{KinGroupIDAncestorView[1]}{}
\genomerow{1705}{179}{3}{Local Jump If}{if a, goto tag match local anchor}{2,6,6}{}{}
\genomerow{1706}{179}{4}{Local Anchor}{register a local jump-to destination}{0,6,7}{}{}
\genomerow{1707}{179}{5}{Modulo}{a = b \% c}{3,0,6}{}{}
\genomerow{1708}{179}{6}{Terminal}{a = value}{7,7,7}{}{-0.241002}
\genomerow{1709}{179}{7}{Terminal}{a = value}{6,5,7}{}{0.134555}
\genomerow{1710}{179}{8}{Add}{a = b + c}{5,4,6}{}{}
\genomerow{1711}{179}{9}{Logical Or}{a = b \textbar{}\textbar{} c}{3,4,6}{}{}
\genomerow{1712}{179}{10}{Write Own State}{if b, state = a}{4,7,4}{TransientNopState[0]}{}
\genomerow{1713}{179}{11}{Logical Or}{a = b \textbar{}\textbar{} c}{4,0,7}{}{}
\genomerow{1714}{179}{12}{Terminal}{a = value}{3,1,0}{}{0.852688}
\genomerow{1715}{179}{13}{Terminal}{a = value}{3,3,0}{}{0.851711}
\genomerow{1716}{179}{14}{Logical Or}{a = b \textbar{}\textbar{} c}{4,0,7}{}{}
\genomerow{1717}{179}{15}{Terminal}{a = value}{3,1,0}{}{0.852932}
\genomerow{1718}{179}{16}{Terminal}{a = value}{3,3,0}{}{0.851711}
\genomerow{1719}{179}{17}{Terminal}{a = value}{5,3,2}{}{0.934728}
\genomerow{1720}{179}{18}{Decay Global Regulator}{decay global regulator value by register a}{5,7,4}{}{}
\genomerow{1721}{180}{0}{Global Anchor}{register a global jump-to destination, maybe terminate}{3,3,4}{}{}
\genomerow{1722}{180}{1}{Local Anchor}{register a local jump-to destination}{5,2,6}{}{}
\genomerow{1723}{180}{2}{Send Intra-Cell Message}{if a, send message another cardinal within cell}{2,1,6}{}{}
\genomerow{1724}{180}{3}{Less Than}{a = ( b < c )}{2,4,3}{}{}
\genomerow{1725}{180}{4}{Generate Random Bool}{if p, a = 1; else a = 0}{4,5,4}{}{}
\genomerow{1726}{180}{5}{Decay Secondary Global Regulator}{decay global regulator value by register a}{7,5,7}{}{}
\genomerow{1727}{180}{6}{Local Anchor}{register a local jump-to destination}{2,4,1}{}{}
\genomerow{1728}{180}{7}{Decay Secondary Global Regulator}{decay global regulator value by register a}{1,3,4}{}{}
\genomerow{1729}{180}{8}{Local Anchor}{register a local jump-to destination}{6,4,7}{}{}
\genomerow{1730}{180}{9}{Generate Random Bool}{if p, a = 1; else a = 0}{6,2,4}{}{}
\genomerow{1731}{180}{10}{Bitwise Shift}{a = b << c}{6,5,7}{}{}
\genomerow{1732}{180}{11}{Draw Random Value}{a = weighted draw from random number generator}{2,5,3}{}{}
\genomerow{1733}{180}{12}{Modulo}{a = b \% c}{6,2,0}{}{}
\genomerow{1734}{180}{13}{Terminal}{a = value}{3,3,2}{}{-30.7864}
\genomerow{1735}{180}{14}{Set Global Regulator}{set global regulator value to a}{7,6,3}{}{}
\genomerow{1736}{180}{15}{Logical And}{a = b \&\& c}{7,2,5}{}{}
\genomerow{1737}{180}{16}{Generate Random Bool}{if p, a = 1; else a = 0}{6,6,5}{}{}
\genomerow{1738}{180}{17}{Get Global Regulator}{a = global regulator value}{3,1,5}{}{}
\genomerow{1739}{180}{18}{Logical And}{a = b \&\& c}{7,2,4}{}{}
\genomerow{1740}{180}{19}{Terminal}{a = value}{6,4,5}{}{-0.581094}
\genomerow{1741}{180}{20}{Local Anchor}{register a local jump-to destination}{3,1,1}{}{}
\genomerow{1742}{181}{0}{Global Anchor}{register a global jump-to destination, maybe terminate}{7,1,1}{}{}
\genomerow{1743}{182}{0}{Global Anchor}{register a global jump-to destination, maybe terminate}{7,7,3}{}{}
\genomerow{1744}{182}{1}{Bitwise Not}{a = \textasciitilde{}b}{3,0,0}{}{}
\genomerow{1745}{182}{2}{Terminal}{a = value}{1,7,4}{}{28.5109}
\genomerow{1746}{182}{3}{Logical Or}{a = b \textbar{}\textbar{} c}{0,2,1}{}{}
\genomerow{1747}{182}{4}{Read Own State}{a = state}{6,7,5}{IsChildGroupOf[0]}{}
\genomerow{1748}{182}{5}{Read Neighbor State}{a = state}{4,7,5}{IncomingInterMessageCounter[0]}{}
\genomerow{1749}{182}{6}{Generate Random Bool}{if p, a = 1; else a = 0}{4,5,4}{}{}
\genomerow{1750}{182}{7}{Terminal}{a = value}{1,6,5}{}{0.934051}
\genomerow{1751}{182}{8}{Local Anchor}{register a local jump-to destination}{1,5,3}{}{}
\genomerow{1752}{182}{9}{Set Local Regulator}{set local regulator value to a}{2,0,2}{}{}
\genomerow{1753}{183}{0}{Global Anchor}{register a global jump-to destination, maybe terminate}{5,7,5}{}{}
\genomerow{1754}{183}{1}{Terminal}{a = value}{4,5,4}{}{-0.791573}
\genomerow{1755}{183}{2}{Terminal}{a = value}{1,6,4}{}{0.933925}
\genomerow{1756}{183}{3}{Generate Random Bool}{if p, a = 1; else a = 0}{4,5,4}{}{}
\genomerow{1757}{184}{0}{Global Anchor}{register a global jump-to destination, maybe terminate}{5,7,5}{}{}
\genomerow{1758}{184}{1}{Terminal}{a = value}{4,5,4}{}{-0.291573}
\genomerow{1759}{184}{2}{Terminal}{a = value}{1,6,4}{}{0.933925}
\genomerow{1760}{184}{3}{Local Anchor}{register a local jump-to destination}{1,5,3}{}{}
\genomerow{1761}{184}{4}{Set Local Regulator}{set local regulator value to a}{2,2,2}{}{}
\genomerow{1762}{184}{5}{Send Intra-Cell Message}{if a, send message another cardinal within cell}{0,7,6}{}{}
\genomerow{1763}{184}{6}{Send Intra-Cell Message}{if a, send message another cardinal within cell}{2,7,7}{}{}
\genomerow{1764}{184}{7}{Get Global Regulator}{a = global regulator value}{7,1,1}{}{}
\genomerow{1765}{184}{8}{Fork If}{if a, submit fork request (processed when core terminates)}{2,4,0}{}{}
\genomerow{1766}{184}{9}{Get Local Regulator}{a = local regulator value}{5,2,4}{}{}
\genomerow{1767}{184}{10}{Local Anchor}{register a local jump-to destination}{5,1,0}{}{}
\genomerow{1768}{184}{11}{Local Anchor}{register a local jump-to destination}{7,1,0}{}{}
\genomerow{1769}{184}{12}{Multiply Own State}{state *= a}{6,7,6}{ResourceSendRequest[0]}{}
\genomerow{1770}{184}{13}{Terminal}{a = value}{0,0,3}{}{-1.22769}
\genomerow{1771}{184}{14}{Write Own State}{if b, state = a}{0,6,5}{ResourceReserveRequest[0]}{}
\genomerow{1772}{184}{15}{Fork If}{if a, submit fork request (processed when core terminates)}{2,4,2}{}{}
\genomerow{1773}{184}{16}{Terminal}{a = value}{5,2,4}{}{0.101411}
\genomerow{1774}{185}{0}{Global Anchor}{register a global jump-to destination, maybe terminate}{5,1,0}{}{}
\genomerow{1775}{185}{1}{Local Anchor}{register a local jump-to destination}{7,1,0}{}{}
\genomerow{1776}{185}{2}{Logical Or}{a = b \textbar{}\textbar{} c}{6,4,0}{}{}
\genomerow{1777}{185}{3}{Read Own State}{a = state}{1,7,2}{TransientNopState[2]}{}
\genomerow{1778}{185}{4}{Read Neighbor State}{a = state}{1,6,2}{ParentFragmented[0]}{}
\genomerow{1779}{185}{5}{Get Global Regulator}{a = global regulator value}{7,6,6}{}{}
\genomerow{1780}{185}{6}{Write Own State}{if b, state = a}{1,2,4}{TransientNopState[1]}{}
\genomerow{1781}{185}{7}{Local Anchor}{register a local jump-to destination}{3,5,4}{}{}
\genomerow{1782}{185}{8}{Read Own State}{a = state}{5,0,7}{MostRecentCauseOfDeath[0]}{}
\genomerow{1783}{185}{9}{Terminal}{a = value}{2,4,3}{}{-1.22767}
\genomerow{1784}{185}{10}{Write Own State}{if b, state = a}{0,4,5}{NopState[3]}{}
\genomerow{1785}{185}{11}{Fork If}{if a, submit fork request (processed when core terminates)}{2,4,2}{}{}
\genomerow{1786}{185}{12}{Get Local Regulator}{a = local regulator value}{5,2,4}{}{}
\genomerow{1787}{186}{0}{Global Anchor}{register a global jump-to destination, maybe terminate}{5,1,2}{}{}
\genomerow{1788}{186}{1}{Not}{a = !a}{2,0,4}{}{}
\genomerow{1789}{187}{0}{Global Anchor}{register a global jump-to destination, maybe terminate}{7,7,2}{}{}
\genomerow{1790}{187}{1}{Get Global Regulator}{a = global regulator value}{7,7,0}{}{}
\genomerow{1791}{187}{2}{Get Global Regulator}{a = global regulator value}{5,3,1}{}{}
\genomerow{1792}{187}{3}{Decay Global Regulator}{decay global regulator value by register a}{3,0,0}{}{}
\genomerow{1793}{187}{4}{Local Jump If}{if a, goto tag match local anchor}{6,1,5}{}{}
\genomerow{1794}{187}{5}{Add To Own State}{state += a}{6,5,3}{TransientNopState[0]}{}
\genomerow{1795}{187}{6}{Local Jump If}{if a, goto tag match local anchor}{6,1,5}{}{}
\genomerow{1796}{187}{7}{Read Own State}{a = state}{6,5,3}{IsChildCellOf[0]}{}
\genomerow{1797}{188}{0}{Global Anchor}{register a global jump-to destination, maybe terminate}{5,3,1}{}{}
\genomerow{1798}{188}{1}{Send Inter-Cell Message}{if a, send message to neighbor cell}{2,6,4}{}{}
\genomerow{1799}{188}{2}{Generate Random Bool}{if p, a = 1; else a = 0}{1,3,6}{}{}
\genomerow{1800}{188}{3}{Read Own State}{a = state}{3,1,3}{ResourceSendLimit[0]}{}
\genomerow{1801}{188}{4}{Read Own State}{a = state}{5,3,7}{KinGroupIDView[1]}{}
\genomerow{1802}{188}{5}{Nop-2}{perform no operation}{6,5,5}{}{}
\genomerow{1803}{188}{6}{Generate Random Bool}{if p, a = 1; else a = 0}{6,7,7}{}{}
\genomerow{1804}{188}{7}{Get Local Regulator}{a = local regulator value}{7,3,3}{}{}
\genomerow{1805}{188}{8}{Terminal}{a = value}{4,1,0}{}{-0.391305}
\genomerow{1806}{188}{9}{Terminate If}{if a, terminate core}{1,4,1}{}{}
\genomerow{1807}{188}{10}{Logical And}{a = b \&\& c}{6,7,6}{}{}
\genomerow{1808}{188}{11}{Local Jump If}{if a, goto tag match local anchor}{3,2,4}{}{}
\genomerow{1809}{188}{12}{Terminal}{a = value}{5,7,7}{}{-11.8496}
\genomerow{1810}{188}{13}{Logical And}{a = b \&\& c}{6,7,2}{}{}
\genomerow{1811}{188}{14}{Local Anchor}{register a local jump-to destination}{3,3,5}{}{}
\genomerow{1812}{188}{15}{Less Than}{a = ( b < c )}{5,6,3}{}{}
\genomerow{1813}{189}{0}{Global Anchor}{register a global jump-to destination, maybe terminate}{6,7,3}{}{}
\genomerow{1814}{189}{1}{Bitwise Not}{a = \textasciitilde{}b}{5,7,7}{}{}
\genomerow{1815}{189}{2}{Read Own State}{a = state}{6,7,4}{NopState[3]}{}
\genomerow{1816}{189}{3}{Local Anchor}{register a local jump-to destination}{3,2,0}{}{}
\genomerow{1817}{189}{4}{Get Local Regulator}{a = local regulator value}{5,7,6}{}{}
\genomerow{1818}{189}{5}{Read Own State}{a = state}{2,7,4}{SpawnedFrom[0]}{}
\genomerow{1819}{189}{6}{Local Anchor}{register a local jump-to destination}{3,3,5}{}{}
\genomerow{1820}{190}{0}{Global Anchor}{register a global jump-to destination, maybe terminate}{3,6,0}{}{}
\genomerow{1821}{190}{1}{Local Anchor}{register a local jump-to destination}{0,5,3}{}{}
\genomerow{1822}{190}{2}{Terminate If}{if a, terminate core}{7,3,6}{}{}
\genomerow{1823}{191}{0}{Global Anchor}{register a global jump-to destination, maybe terminate}{5,0,0}{}{}
\genomerow{1824}{191}{1}{Send Intra-Cell Message}{if a, send message another cardinal within cell}{6,0,3}{}{}
\genomerow{1825}{191}{2}{Terminate If}{if a, terminate core}{2,6,5}{}{}
\genomerow{1826}{191}{3}{Divide}{a = b / c}{2,1,5}{}{}
\genomerow{1827}{191}{4}{Terminal}{a = value}{3,6,0}{}{0.346903}
\genomerow{1828}{191}{5}{Divide}{a = b / c}{2,3,5}{}{}
\genomerow{1829}{191}{6}{Terminal}{a = value}{3,6,0}{}{0.346903}
\genomerow{1830}{192}{0}{Global Anchor}{register a global jump-to destination, maybe terminate}{3,2,0}{}{}
\genomerow{1831}{192}{1}{Logical Or}{a = b \textbar{}\textbar{} c}{2,5,3}{}{}
\genomerow{1832}{193}{0}{Global Anchor}{register a global jump-to destination, maybe terminate}{0,6,4}{}{}
\genomerow{1833}{193}{1}{Local Anchor}{register a local jump-to destination}{1,7,5}{}{}
\genomerow{1834}{193}{2}{Decay Global Regulator}{decay global regulator value by register a}{7,4,3}{}{}
\genomerow{1835}{193}{3}{Not}{a = !a}{4,5,3}{}{}
\genomerow{1836}{193}{4}{Local Anchor}{register a local jump-to destination}{0,4,4}{}{}
\genomerow{1837}{193}{5}{Decay Secondary Global Regulator}{decay global regulator value by register a}{3,2,0}{}{}
\genomerow{1838}{194}{0}{Global Anchor}{register a global jump-to destination, maybe terminate}{2,4,2}{}{}
\genomerow{1839}{194}{1}{Local Anchor}{register a local jump-to destination}{4,2,1}{}{}
\genomerow{1840}{194}{2}{Multiply Own State}{state *= a}{4,5,5}{ResourceReceiveResistance[0]}{}
\genomerow{1841}{195}{0}{Global Anchor}{register a global jump-to destination, maybe terminate}{3,5,1}{}{}
\genomerow{1842}{195}{1}{Terminal}{a = value}{6,6,6}{}{0.0104148}
\genomerow{1843}{195}{2}{Terminal}{a = value}{6,7,1}{}{0.0479729}
\genomerow{1844}{195}{3}{Terminal}{a = value}{0,5,3}{}{368.507}
\genomerow{1845}{195}{4}{Terminal}{a = value}{3,5,4}{}{-168145}
\genomerow{1846}{195}{5}{Modulo}{a = b \% c}{2,5,3}{}{}
\genomerow{1847}{196}{0}{Global Anchor}{register a global jump-to destination, maybe terminate}{0,6,4}{}{}
\genomerow{1848}{196}{1}{Local Anchor}{register a local jump-to destination}{1,7,5}{}{}
\genomerow{1849}{196}{2}{Read Own State}{a = state}{7,3,1}{KinGroupMatch[1]}{}
\genomerow{1850}{196}{3}{Send Intra-Cell Message}{if a, send message another cardinal within cell}{7,0,3}{}{}
\genomerow{1851}{196}{4}{Read Neighbor State}{a = state}{3,1,4}{KinGroupIDView[1]}{}
\genomerow{1852}{196}{5}{Local Anchor}{register a local jump-to destination}{5,2,3}{}{}
\genomerow{1853}{196}{6}{Global Jump If Not}{if !a, goto global anchor; if !b, clear registers}{3,7,4}{}{}
\genomerow{1854}{197}{0}{Global Anchor}{register a global jump-to destination, maybe terminate}{0,6,4}{}{}
\genomerow{1855}{197}{1}{Local Anchor}{register a local jump-to destination}{1,7,5}{}{}
\genomerow{1856}{198}{0}{Global Anchor}{register a global jump-to destination, maybe terminate}{6,3,7}{}{}
\genomerow{1857}{199}{0}{Global Anchor}{register a global jump-to destination, maybe terminate}{0,6,4}{}{}
\genomerow{1858}{199}{1}{Local Anchor}{register a local jump-to destination}{3,2,5}{}{}
\genomerow{1859}{199}{2}{Local Anchor}{register a local jump-to destination}{0,5,0}{}{}
\genomerow{1860}{199}{3}{Local Anchor}{register a local jump-to destination}{0,3,7}{}{}
\genomerow{1861}{199}{4}{Read Own State}{a = state}{7,3,1}{KinGroupMatch[1]}{}
\genomerow{1862}{199}{5}{Read Own State}{a = state}{7,0,3}{KinGroupAge[0]}{}
\genomerow{1863}{199}{6}{Read Neighbor State}{a = state}{3,1,0}{ResourceSendRequest[0]}{}
\genomerow{1864}{199}{7}{Local Anchor}{register a local jump-to destination}{5,2,1}{}{}
\genomerow{1865}{199}{8}{Decay Global Regulator}{decay global regulator value by register a}{3,7,4}{}{}
\genomerow{1866}{199}{9}{Count Ones}{a = popcnt( b )}{3,2,4}{}{}
\genomerow{1867}{199}{10}{Bitwise Or}{a = b \textbar{} c}{2,1,3}{}{}
\genomerow{1868}{199}{11}{Terminal}{a = value}{2,6,2}{}{8.97414}
\genomerow{1869}{199}{12}{Terminal}{a = value}{3,7,1}{}{7.11963}
\genomerow{1870}{200}{0}{Global Anchor}{register a global jump-to destination, maybe terminate}{5,2,1}{}{}
\genomerow{1871}{200}{1}{Decay Global Regulator}{decay global regulator value by register a}{3,7,4}{}{}
\genomerow{1872}{200}{2}{Bitwise Shift}{a = b << c}{3,2,4}{}{}
\genomerow{1873}{200}{3}{Bitwise Xor}{a = b \textasciicircum{} c}{6,3,3}{}{}
\genomerow{1874}{201}{0}{Global Anchor}{register a global jump-to destination, maybe terminate}{2,6,2}{}{}
\genomerow{1875}{201}{1}{Terminal}{a = value}{1,7,1}{}{2.19202}
\genomerow{1876}{201}{2}{Bitwise Not}{a = \textasciitilde{}b}{4,2,3}{}{}
\genomerow{1877}{201}{3}{Terminal}{a = value}{3,7,1}{}{7.12089}
\genomerow{1878}{201}{4}{Read Neighbor State}{a = state}{5,2,1}{KinGroupIDAncestorView[0]}{}
\genomerow{1879}{201}{5}{Decay Global Regulator}{decay global regulator value by register a}{3,7,4}{}{}
\genomerow{1880}{201}{6}{Bitwise Shift}{a = b << c}{3,2,4}{}{}
\genomerow{1881}{201}{7}{Bitwise Or}{a = b \textbar{} c}{6,3,3}{}{}
\genomerow{1882}{202}{0}{Global Anchor}{register a global jump-to destination, maybe terminate}{2,6,2}{}{}
\genomerow{1883}{202}{1}{Terminal}{a = value}{1,7,1}{}{7.11963}
\genomerow{1884}{202}{2}{Bitwise Not}{a = \textasciitilde{}b}{5,2,3}{}{}
\genomerow{1885}{202}{3}{Get Local Regulator}{a = local regulator value}{3,5,3}{}{}
\genomerow{1886}{203}{0}{Global Anchor}{register a global jump-to destination, maybe terminate}{5,2,1}{}{}
\genomerow{1887}{203}{1}{Terminal}{a = value}{2,7,4}{}{0.487232}
\genomerow{1888}{203}{2}{Count Ones}{a = popcnt( b )}{3,0,0}{}{}
\genomerow{1889}{203}{3}{Terminal}{a = value}{2,7,4}{}{0.486988}
\genomerow{1890}{203}{4}{Subtract}{a = b - c}{3,0,4}{}{}
\genomerow{1891}{203}{5}{Bitwise Or}{a = b \textbar{} c}{2,1,3}{}{}
\genomerow{1892}{203}{6}{Local Anchor}{register a local jump-to destination}{3,6,1}{}{}
\genomerow{1893}{203}{7}{Send Intra-Cell Message}{if a, send message another cardinal within cell}{2,6,1}{}{}
\genomerow{1894}{203}{8}{Local Anchor}{register a local jump-to destination}{3,7,1}{}{}
\genomerow{1895}{203}{9}{Bitwise Not}{a = \textasciitilde{}b}{4,2,3}{}{}
\genomerow{1896}{203}{10}{Terminal}{a = value}{3,1,3}{}{8.58029}
\genomerow{1897}{203}{11}{Bitwise Or}{a = b \textbar{} c}{7,3,4}{}{}
\genomerow{1898}{203}{12}{Bitwise Xor}{a = b \textasciicircum{} c}{0,7,6}{}{}
\genomerow{1899}{203}{13}{Get Global Regulator}{a = global regulator value}{5,5,1}{}{}
\genomerow{1900}{204}{0}{Global Anchor}{register a global jump-to destination, maybe terminate}{5,2,7}{}{}
\genomerow{1901}{204}{1}{Add}{a = b + c}{6,1,6}{}{}
\genomerow{1902}{204}{2}{Nop-2}{perform no operation}{5,7,7}{}{}
\genomerow{1903}{204}{3}{Terminal}{a = value}{5,6,3}{}{-0.453247}
\genomerow{1904}{204}{4}{Bitwise Not}{a = \textasciitilde{}b}{6,3,1}{}{}
\genomerow{1905}{204}{5}{Get Global Regulator}{a = global regulator value}{2,1,5}{}{}
\genomerow{1906}{204}{6}{Modulo}{a = b \% c}{0,7,1}{}{}
\genomerow{1907}{205}{0}{Global Anchor}{register a global jump-to destination, maybe terminate}{5,2,7}{}{}
\genomerow{1908}{205}{1}{Add}{a = b + c}{6,1,6}{}{}
\genomerow{1909}{205}{2}{Nop-2}{perform no operation}{5,7,7}{}{}
\genomerow{1910}{205}{3}{Terminal}{a = value}{5,6,3}{}{-0.453247}
\genomerow{1911}{205}{4}{Bitwise Not}{a = \textasciitilde{}b}{6,3,1}{}{}
\genomerow{1912}{205}{5}{Get Global Regulator}{a = global regulator value}{2,1,5}{}{}
\genomerow{1913}{205}{6}{Modulo}{a = b \% c}{0,7,1}{}{}
\genomerow{1914}{205}{7}{Local Anchor}{register a local jump-to destination}{5,5,3}{}{}
\genomerow{1915}{205}{8}{Local Anchor}{register a local jump-to destination}{5,3,1}{}{}
\genomerow{1916}{205}{9}{Bitwise Xor}{a = b \textasciicircum{} c}{2,7,6}{}{}
\genomerow{1917}{205}{10}{Logical And}{a = b \&\& c}{7,3,1}{}{}
\genomerow{1918}{205}{11}{Bitwise Or}{a = b \textbar{} c}{5,5,3}{}{}
\genomerow{1919}{205}{12}{Logical Or}{a = b \textbar{}\textbar{} c}{5,7,1}{}{}
\genomerow{1920}{205}{13}{Bitwise Xor}{a = b \textasciicircum{} c}{2,7,6}{}{}
\genomerow{1921}{205}{14}{Get Global Regulator}{a = global regulator value}{5,5,1}{}{}
\genomerow{1922}{205}{15}{Get Global Regulator}{a = global regulator value}{5,5,1}{}{}
\genomerow{1923}{205}{16}{Logical Or}{a = b \textbar{}\textbar{} c}{0,0,3}{}{}
\genomerow{1924}{206}{0}{Global Anchor}{register a global jump-to destination, maybe terminate}{0,6,4}{}{}
\genomerow{1925}{207}{0}{Global Anchor}{register a global jump-to destination, maybe terminate}{0,6,4}{}{}
\genomerow{1926}{207}{1}{Bitwise Xor}{a = b \textasciicircum{} c}{4,4,2}{}{}
\genomerow{1927}{207}{2}{Read Neighbor State}{a = state}{2,6,5}{NeighborKinGroupWillExpire[1]}{}
\genomerow{1928}{208}{0}{Global Anchor}{register a global jump-to destination, maybe terminate}{6,7,7}{}{}
\genomerow{1929}{208}{1}{Draw Random Value}{a = weighted draw from random number generator}{4,5,5}{}{}
\genomerow{1930}{208}{2}{Bitwise Not}{a = \textasciitilde{}b}{6,6,4}{}{}
\genomerow{1931}{208}{3}{Terminal}{a = value}{3,0,1}{}{-1.23276}
\genomerow{1932}{208}{4}{Decay Global Regulator}{decay global regulator value by register a}{3,4,1}{}{}
\genomerow{1933}{208}{5}{Logical Or}{a = b \textbar{}\textbar{} c}{6,5,4}{}{}
\genomerow{1934}{209}{0}{Global Anchor}{register a global jump-to destination, maybe terminate}{3,1,4}{}{}
\genomerow{1935}{209}{1}{Adjust Global Regulator}{adjust global regulator value by a}{1,6,3}{}{}
\genomerow{1936}{209}{2}{Local Anchor}{register a local jump-to destination}{5,1,3}{}{}
\genomerow{1937}{209}{3}{Draw Random Value}{a = weighted draw from random number generator}{2,4,1}{}{}
\genomerow{1938}{209}{4}{Get Global Regulator}{a = global regulator value}{3,7,3}{}{}
\genomerow{1939}{209}{5}{Terminate If}{if a, terminate core}{2,3,3}{}{}
\genomerow{1940}{209}{6}{Read Neighbor State}{a = state}{3,7,2}{KinGroupAge[1]}{}
\genomerow{1941}{210}{0}{Global Anchor}{register a global jump-to destination, maybe terminate}{2,2,4}{}{}
\genomerow{1942}{210}{1}{Read Own State}{a = state}{1,4,6}{RepLevRequest[0]}{}
\genomerow{1943}{210}{2}{Local Anchor}{register a local jump-to destination}{6,3,6}{}{}
\genomerow{1944}{210}{3}{Bitwise And}{a = b \& c}{1,6,2}{}{}
\genomerow{1945}{210}{4}{Not Equal}{a = (b != c)}{7,4,0}{}{}
\genomerow{1946}{210}{5}{Local Anchor}{register a local jump-to destination}{4,4,0}{}{}
\genomerow{1947}{210}{6}{Add To Own State}{state += a}{7,5,5}{ResourceReserveRequest[0]}{}
\genomerow{1948}{210}{7}{Subtract}{a = b - c}{3,6,0}{}{}
\genomerow{1949}{210}{8}{Multiply Own State}{state *= a}{7,5,5}{ResourceReserveRequest[0]}{}
\genomerow{1950}{210}{9}{Subtract}{a = b - c}{3,6,0}{}{}
\genomerow{1951}{210}{10}{Not Equal}{a = (b != c)}{1,0,1}{}{}
\genomerow{1952}{211}{0}{Global Anchor}{register a global jump-to destination, maybe terminate}{0,7,5}{}{}
\genomerow{1953}{211}{1}{Decay Secondary Global Regulator}{decay global regulator value by register a}{2,6,5}{}{}
\genomerow{1954}{212}{0}{Global Anchor}{register a global jump-to destination, maybe terminate}{0,7,5}{}{}
\genomerow{1955}{212}{1}{Terminal}{a = value}{4,7,6}{}{-1.68841}
\genomerow{1956}{212}{2}{Less Than}{a = ( b < c )}{1,4,3}{}{}
\genomerow{1957}{212}{3}{Read Own State}{a = state}{1,0,1}{IncomingInterMessageCounter[0]}{}
\genomerow{1958}{213}{0}{Global Anchor}{register a global jump-to destination, maybe terminate}{4,3,4}{}{}
\genomerow{1959}{213}{1}{Terminal}{a = value}{4,1,3}{}{-0.885759}
\genomerow{1960}{214}{0}{Global Anchor}{register a global jump-to destination, maybe terminate}{4,1,3}{}{}
\genomerow{1961}{214}{1}{Terminal}{a = value}{6,4,6}{}{-0.30155}
\genomerow{1962}{214}{2}{Terminal}{a = value}{7,1,7}{}{11.5699}
\genomerow{1963}{214}{3}{Read Own State}{a = state}{6,6,5}{IsParentCellOf[0]}{}
\genomerow{1964}{214}{4}{Less Than}{a = ( b < c )}{2,6,5}{}{}
\genomerow{1965}{214}{5}{Read Own State}{a = state}{7,7,1}{RepLevRequest[1]}{}
\genomerow{1966}{214}{6}{Greater Than}{a = ( b > c )}{2,6,5}{}{}
\genomerow{1967}{214}{7}{Send Inter-Cell Message}{if a, send message to neighbor cell}{5,1,3}{}{}
\genomerow{1968}{214}{8}{Read Own State}{a = state}{2,6,2}{NeighborApoptosis[0]}{}
\genomerow{1969}{214}{9}{Local Anchor}{register a local jump-to destination}{2,6,5}{}{}
\genomerow{1970}{214}{10}{Logical And}{a = b \&\& c}{1,2,4}{}{}
\genomerow{1971}{214}{11}{Terminal}{a = value}{3,2,1}{}{-12.6989}
\genomerow{1972}{214}{12}{Write Own State}{if b, state = a}{6,1,1}{NopState[1]}{}
\genomerow{1973}{214}{13}{Adjust Local Regulator}{adjust local regulator value by a}{3,4,5}{}{}
\genomerow{1974}{214}{14}{Local Jump If}{if a, goto tag match local anchor}{7,3,0}{}{}
\genomerow{1975}{214}{15}{Terminal}{a = value}{5,4,0}{}{-0.790204}
\genomerow{1976}{214}{16}{Terminal}{a = value}{6,3,1}{}{-0.0163396}
\genomerow{1977}{214}{17}{Not Equal}{a = (b != c)}{5,3,0}{}{}
\genomerow{1978}{214}{18}{Generate Random Bool}{if p, a = 1; else a = 0}{5,3,0}{}{}
\genomerow{1979}{214}{19}{Terminal}{a = value}{4,6,7}{}{292.184}
\genomerow{1980}{214}{20}{Nop-2}{perform no operation}{4,2,1}{}{}
\genomerow{1981}{215}{0}{Global Anchor}{register a global jump-to destination, maybe terminate}{2,1,0}{}{}
\genomerow{1982}{216}{0}{Global Anchor}{register a global jump-to destination, maybe terminate}{4,2,4}{}{}
\genomerow{1983}{217}{0}{Global Anchor}{register a global jump-to destination, maybe terminate}{4,4,7}{}{}
\genomerow{1984}{217}{1}{Terminal}{a = value}{6,3,4}{}{2746.07}
\genomerow{1985}{217}{2}{Add To Own State}{state += a}{5,2,7}{RepLevRequest[0]}{}
\genomerow{1986}{217}{3}{Terminal}{a = value}{2,7,0}{}{-76.7227}
\genomerow{1987}{217}{4}{Terminal}{a = value}{2,2,7}{}{0.221822}
\genomerow{1988}{218}{0}{Global Anchor}{register a global jump-to destination, maybe terminate}{4,4,7}{}{}
\genomerow{1989}{218}{1}{Terminal}{a = value}{4,3,6}{}{2746.07}
\genomerow{1990}{218}{2}{Add To Own State}{state += a}{5,2,7}{RepLevRequest[0]}{}
\genomerow{1991}{218}{3}{Terminal}{a = value}{2,7,0}{}{-76.7227}
\genomerow{1992}{218}{4}{Terminal}{a = value}{2,2,7}{}{-0.778178}
\genomerow{1993}{218}{5}{Terminal}{a = value}{1,6,7}{}{-3.13266}
\genomerow{1994}{218}{6}{Terminal}{a = value}{1,4,1}{}{0.819635}
\genomerow{1995}{219}{0}{Global Anchor}{register a global jump-to destination, maybe terminate}{5,3,7}{}{}
\genomerow{1996}{219}{1}{Write Own State}{if b, state = a}{4,3,6}{ResourceReceiveResistance[0]}{}
\genomerow{1997}{219}{2}{Greater Than}{a = ( b > c )}{2,6,5}{}{}
\genomerow{1998}{219}{3}{Write Own State}{if b, state = a}{6,1,1}{NopState[1]}{}
\genomerow{1999}{219}{4}{Terminal}{a = value}{6,6,3}{}{-31.6452}
\genomerow{2000}{219}{5}{Local Anchor}{register a local jump-to destination}{6,2,2}{}{}
\genomerow{2001}{219}{6}{Terminal}{a = value}{0,3,7}{}{-1.57359}
\genomerow{2002}{219}{7}{Generate Random Bool}{if p, a = 1; else a = 0}{4,6,6}{}{}
\genomerow{2003}{219}{8}{Negate}{a = -a}{2,7,5}{}{}
\genomerow{2004}{219}{9}{Local Jump If}{if a, goto tag match local anchor}{3,5,2}{}{}
\genomerow{2005}{219}{10}{Not Equal}{a = (b != c)}{6,6,2}{}{}
\genomerow{2006}{219}{11}{Local Anchor}{register a local jump-to destination}{0,0,6}{}{}
\genomerow{2007}{219}{12}{Local Jump If}{if a, goto tag match local anchor}{3,5,3}{}{}
\genomerow{2008}{219}{13}{Read Neighbor State}{a = state}{3,5,0}{NeighborIsNewborn[0]}{}
\genomerow{2009}{219}{14}{Get Local Regulator}{a = local regulator value}{3,1,0}{}{}
\genomerow{2010}{219}{15}{Send Inter-Cell Message}{if a, send message to neighbor cell}{6,6,2}{}{}
\genomerow{2011}{219}{16}{Terminal}{a = value}{0,0,6}{}{-0.631959}
\genomerow{2012}{219}{17}{Local Anchor}{register a local jump-to destination}{4,7,5}{}{}
\genomerow{2013}{219}{18}{RandomFill}{fill a with uniformly-distributed random bits}{3,5,1}{}{}
\genomerow{2014}{219}{19}{Local Anchor}{register a local jump-to destination}{3,2,7}{}{}
\genomerow{2015}{219}{20}{Terminal}{a = value}{3,3,2}{}{-0.809985}
\genomerow{2016}{219}{21}{Bitwise And}{a = b \& c}{5,5,1}{}{}
\genomerow{2017}{219}{22}{Read Own State}{a = state}{4,3,1}{IncomingInterMessageCounter[0]}{}
\genomerow{2018}{219}{23}{Local Anchor}{register a local jump-to destination}{4,7,5}{}{}
\genomerow{2019}{219}{24}{RandomFill}{fill a with uniformly-distributed random bits}{3,5,0}{}{}
\genomerow{2020}{219}{25}{Local Anchor}{register a local jump-to destination}{3,2,7}{}{}
\genomerow{2021}{219}{26}{Terminal}{a = value}{1,3,2}{}{-0.809985}
\genomerow{2022}{219}{27}{Bitwise And}{a = b \& c}{5,5,3}{}{}
\genomerow{2023}{219}{28}{Read Own State}{a = state}{4,3,1}{IncomingInterMessageCounter[0]}{}
\genomerow{2024}{220}{0}{Global Anchor}{register a global jump-to destination, maybe terminate}{4,7,7}{}{}
\genomerow{2025}{220}{1}{Nop-2}{perform no operation}{0,0,1}{}{}
\genomerow{2026}{220}{2}{Add}{a = b + c}{3,7,4}{}{}
\genomerow{2027}{221}{0}{Global Anchor}{register a global jump-to destination, maybe terminate}{4,2,3}{}{}
\genomerow{2028}{221}{1}{Write Own State}{if b, state = a}{3,0,4}{RepLevRequest[1]}{}
\genomerow{2029}{221}{2}{Read Own State}{a = state}{0,2,7}{IsNewborn[0]}{}
\genomerow{2030}{222}{0}{Global Anchor}{register a global jump-to destination, maybe terminate}{0,2,4}{}{}
\genomerow{2031}{222}{1}{Read Own State}{a = state}{1,1,4}{SpawnArrest[0]}{}
\genomerow{2032}{222}{2}{Local Anchor}{register a local jump-to destination}{6,0,7}{}{}
\genomerow{2033}{223}{0}{Global Anchor}{register a global jump-to destination, maybe terminate}{4,2,3}{}{}
\genomerow{2034}{223}{1}{Write Own State}{if b, state = a}{3,0,4}{TransientNopState[3]}{}
\genomerow{2035}{223}{2}{Read Own State}{a = state}{4,2,7}{NopState[3]}{}
\genomerow{2036}{223}{3}{Read Own State}{a = state}{0,2,4}{TransientNopState[3]}{}
\genomerow{2037}{223}{4}{Read Own State}{a = state}{1,1,4}{IsParentGroupOf[1]}{}
\genomerow{2038}{223}{5}{Local Anchor}{register a local jump-to destination}{6,0,7}{}{}
\genomerow{2039}{224}{0}{Global Anchor}{register a global jump-to destination, maybe terminate}{4,2,3}{}{}
\genomerow{2040}{224}{1}{Write Own State}{if b, state = a}{3,0,4}{ResourceSendRequest[0]}{}
\genomerow{2041}{224}{2}{Local Anchor}{register a local jump-to destination}{1,2,4}{}{}
\genomerow{2042}{224}{3}{Write Own State}{if b, state = a}{3,0,4}{TransientNopState[3]}{}
\genomerow{2043}{224}{4}{Read Own State}{a = state}{4,2,7}{NopState[3]}{}
\genomerow{2044}{224}{5}{Read Own State}{a = state}{0,2,4}{TransientNopState[3]}{}
\genomerow{2045}{224}{6}{Read Own State}{a = state}{1,1,4}{IsParentGroupOf[1]}{}
\genomerow{2046}{224}{7}{Local Anchor}{register a local jump-to destination}{6,0,7}{}{}
\genomerow{2047}{225}{0}{Global Anchor}{register a global jump-to destination, maybe terminate}{4,2,3}{}{}
\genomerow{2048}{225}{1}{Write Own State}{if b, state = a}{3,0,0}{ResourceSendRequest[0]}{}
\genomerow{2049}{225}{2}{Local Anchor}{register a local jump-to destination}{1,2,4}{}{}
\genomerow{2050}{225}{3}{Write Own State}{if b, state = a}{3,0,4}{ResourceSendRequest[0]}{}
\genomerow{2051}{225}{4}{Local Anchor}{register a local jump-to destination}{1,2,4}{}{}
\genomerow{2052}{225}{5}{Multiply}{a = b * c}{2,2,0}{}{}
\genomerow{2053}{225}{6}{Local Jump If}{if a, goto tag match local anchor}{7,7,4}{}{}
\genomerow{2054}{226}{0}{Global Anchor}{register a global jump-to destination, maybe terminate}{6,1,6}{}{}
\genomerow{2055}{226}{1}{Decay Secondary Global Regulator}{decay global regulator value by register a}{5,7,5}{}{}
\genomerow{2056}{226}{2}{Adjust Local Regulator}{adjust local regulator value by a}{3,0,2}{}{}
\genomerow{2057}{227}{0}{Global Anchor}{register a global jump-to destination, maybe terminate}{3,0,6}{}{}
\genomerow{2058}{227}{1}{Bitwise Xor}{a = b \textasciicircum{} c}{4,6,5}{}{}
\genomerow{2059}{227}{2}{Terminal}{a = value}{6,6,2}{}{1831.26}
\genomerow{2060}{228}{0}{Global Anchor}{register a global jump-to destination, maybe terminate}{6,2,2}{}{}
\genomerow{2061}{229}{0}{Global Anchor}{register a global jump-to destination, maybe terminate}{1,3,2}{}{}
\genomerow{2062}{229}{1}{Read Own State}{a = state}{3,0,5}{KinGroupWillExpire[0]}{}
\genomerow{2063}{229}{2}{Terminal}{a = value}{3,1,4}{}{181.581}
\genomerow{2064}{229}{3}{Read Own State}{a = state}{7,2,2}{IsChildCellOf[0]}{}
\genomerow{2065}{229}{4}{Terminal}{a = value}{0,2,3}{}{-23.846}
\genomerow{2066}{230}{0}{Global Anchor}{register a global jump-to destination, maybe terminate}{6,2,2}{}{}
\genomerow{2067}{230}{1}{Read Neighbor State}{a = state}{1,3,2}{SpawnRequest[0]}{}
\genomerow{2068}{230}{2}{Read Own State}{a = state}{3,0,5}{KinGroupWillExpire[0]}{}
\genomerow{2069}{230}{3}{Terminal}{a = value}{3,1,4}{}{181.581}
\genomerow{2070}{230}{4}{Read Own State}{a = state}{7,2,2}{IsChildCellOf[0]}{}
\genomerow{2071}{230}{5}{Terminal}{a = value}{0,2,3}{}{-23.846}
\genomerow{2072}{230}{6}{Less Than}{a = ( b < c )}{5,5,6}{}{}
\genomerow{2073}{230}{7}{Decrement}{a = a - 1}{7,5,5}{}{}
\genomerow{2074}{230}{8}{Get Global Regulator}{a = global regulator value}{5,1,1}{}{}
\genomerow{2075}{230}{9}{Bitwise Not}{a = \textasciitilde{}b}{5,5,1}{}{}
\genomerow{2076}{230}{10}{Send Inter-Cell Message}{if a, send message to neighbor cell}{4,6,5}{}{}
\genomerow{2077}{230}{11}{Get Local Regulator}{a = local regulator value}{2,0,6}{}{}
\genomerow{2078}{230}{12}{Terminal}{a = value}{6,7,6}{}{-0.517813}
\genomerow{2079}{230}{13}{Decrement}{a = a - 1}{6,5,5}{}{}
\genomerow{2080}{230}{14}{Get Global Regulator}{a = global regulator value}{5,1,1}{}{}
\genomerow{2081}{230}{15}{Bitwise Not}{a = \textasciitilde{}b}{5,5,1}{}{}
\genomerow{2082}{230}{16}{Send Inter-Cell Message}{if a, send message to neighbor cell}{4,6,5}{}{}
\genomerow{2083}{230}{17}{Get Local Regulator}{a = local regulator value}{2,0,2}{}{}
\genomerow{2084}{230}{18}{Terminal}{a = value}{3,7,6}{}{-0.517813}
\genomerow{2085}{230}{19}{Divide}{a = b / c}{5,6,6}{}{}
\genomerow{2086}{230}{20}{Local Anchor}{register a local jump-to destination}{0,0,2}{}{}
\genomerow{2087}{230}{21}{Send Intra-Cell Message}{if a, send message another cardinal within cell}{0,6,5}{}{}
\genomerow{2088}{230}{22}{Terminal}{a = value}{4,1,1}{}{0.308834}
\genomerow{2089}{230}{23}{Bitwise Not}{a = \textasciitilde{}b}{0,6,5}{}{}
\genomerow{2090}{230}{24}{Get Local Regulator}{a = local regulator value}{2,0,2}{}{}
\genomerow{2091}{230}{25}{Terminal}{a = value}{7,7,6}{}{-0.517813}
\genomerow{2092}{230}{26}{Nop-1}{perform no operation}{5,6,6}{}{}
\genomerow{2093}{230}{27}{Local Anchor}{register a local jump-to destination}{0,0,2}{}{}
\genomerow{2094}{230}{28}{Write Own State}{if b, state = a}{0,6,5}{TransientNopState[1]}{}
\genomerow{2095}{230}{29}{Terminal}{a = value}{4,1,1}{}{0.308834}
\genomerow{2096}{230}{30}{Bitwise Not}{a = \textasciitilde{}b}{0,6,7}{}{}
\genomerow{2097}{230}{31}{Get Local Regulator}{a = local regulator value}{2,0,2}{}{}
\genomerow{2098}{230}{32}{Terminal}{a = value}{7,7,6}{}{-0.517813}
\genomerow{2099}{230}{33}{Divide}{a = b / c}{5,6,6}{}{}
\genomerow{2100}{230}{34}{Local Anchor}{register a local jump-to destination}{0,0,2}{}{}
\genomerow{2101}{230}{35}{Send Intra-Cell Message}{if a, send message another cardinal within cell}{0,6,5}{}{}
\genomerow{2102}{230}{36}{Terminal}{a = value}{4,1,1}{}{0.310787}
\genomerow{2103}{230}{37}{Bitwise Not}{a = \textasciitilde{}b}{0,6,7}{}{}
\genomerow{2104}{230}{38}{Send Intra-Cell Message}{if a, send message another cardinal within cell}{2,6,3}{}{}
\genomerow{2105}{230}{39}{Local Anchor}{register a local jump-to destination}{3,1,5}{}{}
\genomerow{2106}{230}{40}{Terminal}{a = value}{6,4,1}{}{388.118}
\genomerow{2107}{230}{41}{Draw Random Value}{a = weighted draw from random number generator}{6,4,0}{}{}
\genomerow{2108}{230}{42}{Write Own State}{if b, state = a}{3,1,2}{TransientNopState[0]}{}
\genomerow{2109}{230}{43}{Local Anchor}{register a local jump-to destination}{2,0,2}{}{}
\genomerow{2110}{230}{44}{Terminal}{a = value}{5,7,6}{}{-0.517814}
\genomerow{2111}{231}{0}{Global Anchor}{register a global jump-to destination, maybe terminate}{6,7,6}{}{}
\genomerow{2112}{232}{0}{Global Anchor}{register a global jump-to destination, maybe terminate}{6,0,2}{}{}
\genomerow{2113}{232}{1}{Decay Secondary Global Regulator}{decay global regulator value by register a}{7,4,0}{}{}
\genomerow{2114}{232}{2}{Get Global Regulator}{a = global regulator value}{5,6,4}{}{}
\genomerow{2115}{232}{3}{Local Anchor}{register a local jump-to destination}{5,1,3}{}{}
\genomerow{2116}{232}{4}{Bitwise Shift}{a = b << c}{5,5,1}{}{}
\genomerow{2117}{233}{0}{Global Anchor}{register a global jump-to destination, maybe terminate}{0,7,1}{}{}
\genomerow{2118}{234}{0}{Global Anchor}{register a global jump-to destination, maybe terminate}{2,0,2}{}{}
\genomerow{2119}{235}{0}{Global Anchor}{register a global jump-to destination, maybe terminate}{0,7,2}{}{}
\genomerow{2120}{235}{1}{Set Local Regulator}{set local regulator value to a}{1,1,1}{}{}
\genomerow{2121}{235}{2}{Nop-1}{perform no operation}{4,7,7}{}{}
\genomerow{2122}{235}{3}{Local Jump If Not}{if !a, goto tag match local anchor}{1,7,2}{}{}
\genomerow{2123}{235}{4}{Get Secondary Global Regulator}{a = global regulator value}{4,6,5}{}{}
\genomerow{2124}{235}{5}{Read Own State}{a = state}{4,1,6}{IncomingIntraMessageCounter[0]}{}
\genomerow{2125}{235}{6}{Subtract}{a = b - c}{7,2,4}{}{}
\genomerow{2126}{235}{7}{Terminal}{a = value}{2,7,6}{}{0.756388}
\genomerow{2127}{235}{8}{Send Inter-Cell Message}{if a, send message to neighbor cell}{5,4,2}{}{}
\genomerow{2128}{236}{0}{Global Anchor}{register a global jump-to destination, maybe terminate}{0,0,6}{}{}
\genomerow{2129}{236}{1}{Get Local Regulator}{a = local regulator value}{1,0,2}{}{}
\genomerow{2130}{236}{2}{Local Anchor}{register a local jump-to destination}{3,5,7}{}{}
\genomerow{2131}{236}{3}{Multiply Own State}{state *= a}{7,6,5}{RepLevRequest[1]}{}
\genomerow{2132}{236}{4}{Terminal}{a = value}{7,0,5}{}{0.912495}
\genomerow{2133}{236}{5}{Local Jump If}{if a, goto tag match local anchor}{5,4,1}{}{}
\genomerow{2134}{236}{6}{Multiply Own State}{state *= a}{7,6,4}{ResourceReserveRequest[0]}{}
\genomerow{2135}{236}{7}{Read Own State}{a = state}{5,0,5}{MostRecentCauseOfDeath[0]}{}
\genomerow{2136}{236}{8}{Count Ones}{a = popcnt( b )}{5,4,1}{}{}
\genomerow{2137}{236}{9}{Read Neighbor State}{a = state}{7,6,4}{CellAge[0]}{}
\genomerow{2138}{236}{10}{Terminal}{a = value}{3,2,4}{}{0.912502}
\genomerow{2139}{236}{11}{Local Anchor}{register a local jump-to destination}{6,1,6}{}{}
\genomerow{2140}{236}{12}{Terminal}{a = value}{1,3,1}{}{-5.8736}
\genomerow{2141}{236}{13}{Terminal}{a = value}{3,4,2}{}{0.656454}
\genomerow{2142}{236}{14}{Local Jump If Not}{if !a, goto tag match local anchor}{6,6,3}{}{}
\genomerow{2143}{236}{15}{Local Anchor}{register a local jump-to destination}{2,6,7}{}{}
\genomerow{2144}{236}{16}{Local Anchor}{register a local jump-to destination}{6,2,1}{}{}
\genomerow{2145}{237}{0}{Global Anchor}{register a global jump-to destination, maybe terminate}{2,7,3}{}{}
\genomerow{2146}{238}{0}{Global Anchor}{register a global jump-to destination, maybe terminate}{2,2,0}{}{}
\genomerow{2147}{239}{0}{Global Anchor}{register a global jump-to destination, maybe terminate}{0,7,3}{}{}
\genomerow{2148}{239}{1}{Terminal}{a = value}{0,1,0}{}{-1.65225}
\genomerow{2149}{239}{2}{Local Anchor}{register a local jump-to destination}{0,4,0}{}{}
\genomerow{2150}{239}{3}{Terminal}{a = value}{7,6,1}{}{96.0258}
\genomerow{2151}{239}{4}{Equal}{a = ( b == c )}{1,0,5}{}{}
\genomerow{2152}{239}{5}{Read Own State}{a = state}{5,0,1}{CellAge[0]}{}
\genomerow{2153}{239}{6}{Bitwise And}{a = b \& c}{1,5,3}{}{}
\genomerow{2154}{239}{7}{Terminate If}{if a, terminate core}{3,3,4}{}{}
\genomerow{2155}{239}{8}{Read Own State}{a = state}{7,4,6}{IsAlive[0]}{}
\genomerow{2156}{239}{9}{Multiply Own State}{state *= a}{1,6,1}{ApoptosisRequest[0]}{}
\genomerow{2157}{239}{10}{Terminate If}{if a, terminate core}{3,3,4}{}{}
\genomerow{2158}{239}{11}{Broadcast Intra-Cell Message}{if a, send message to all other cardinals within cell}{1,5,3}{}{}
\genomerow{2159}{239}{12}{Local Anchor}{register a local jump-to destination}{3,3,3}{}{}
\genomerow{2160}{239}{13}{Bitwise And}{a = b \& c}{1,5,3}{}{}
\genomerow{2161}{239}{14}{Negate}{a = -a}{7,4,6}{}{}
\genomerow{2162}{239}{15}{Multiply Own State}{state *= a}{1,6,1}{TransientNopState[2]}{}
\genomerow{2163}{239}{16}{Read Own State}{a = state}{2,1,1}{NopState[3]}{}
\genomerow{2164}{239}{17}{Set Local Regulator}{set local regulator value to a}{3,3,4}{}{}
\genomerow{2165}{239}{18}{Broadcast Intra-Cell Message}{if a, send message to all other cardinals within cell}{1,5,3}{}{}
\genomerow{2166}{239}{19}{Local Anchor}{register a local jump-to destination}{3,3,3}{}{}
\genomerow{2167}{239}{20}{Bitwise And}{a = b \& c}{1,5,3}{}{}
\genomerow{2168}{239}{21}{Write Own State}{if b, state = a}{7,4,6}{TransientNopState[3]}{}
\genomerow{2169}{239}{22}{Nop-2}{perform no operation}{6,4,6}{}{}
\genomerow{2170}{239}{23}{Decay Secondary Global Regulator}{decay global regulator value by register a}{1,3,6}{}{}
\genomerow{2171}{239}{24}{Local Anchor}{register a local jump-to destination}{3,2,0}{}{}
\genomerow{2172}{239}{25}{Less Than}{a = ( b < c )}{0,7,0}{}{}
\genomerow{2173}{239}{26}{Multiply Own State}{state *= a}{1,6,1}{ApoptosisRequest[0]}{}
\genomerow{2174}{239}{27}{Terminal}{a = value}{4,1,4}{}{0.213481}
\genomerow{2175}{239}{28}{Equal}{a = ( b == c )}{6,4,4}{}{}
\genomerow{2176}{239}{29}{Write Own State}{if b, state = a}{5,0,1}{TransientNopState[2]}{}
\genomerow{2177}{239}{30}{Send Inter-Cell Message}{if a, send message to neighbor cell}{0,5,5}{}{}
\genomerow{2178}{239}{31}{Divide}{a = b / c}{1,3,6}{}{}
\genomerow{2179}{239}{32}{Read Own State}{a = state}{3,1,1}{NopState[3]}{}
\genomerow{2180}{239}{33}{Adjust Global Regulator}{adjust global regulator value by a}{2,6,3}{}{}
\genomerow{2181}{239}{34}{Read Own State}{a = state}{3,1,1}{CellAge[0]}{}
\genomerow{2182}{239}{35}{Local Anchor}{register a local jump-to destination}{2,6,3}{}{}
\genomerow{2183}{239}{36}{Send Intra-Cell Message}{if a, send message another cardinal within cell}{6,4,0}{}{}
\genomerow{2184}{239}{37}{Get Local Regulator}{a = local regulator value}{5,1,6}{}{}
\genomerow{2185}{239}{38}{Terminal}{a = value}{2,4,5}{}{-13.9577}
\genomerow{2186}{240}{0}{Global Anchor}{register a global jump-to destination, maybe terminate}{6,1,4}{}{}
\genomerow{2187}{240}{1}{Set Global Regulator}{set global regulator value to a}{7,6,4}{}{}
\genomerow{2188}{240}{2}{Send Intra-Cell Message}{if a, send message another cardinal within cell}{1,3,0}{}{}
\genomerow{2189}{241}{0}{Global Anchor}{register a global jump-to destination, maybe terminate}{3,2,2}{}{}
\genomerow{2190}{241}{1}{Read Own State}{a = state}{6,5,3}{ParentFragmented[0]}{}
\genomerow{2191}{242}{0}{Global Anchor}{register a global jump-to destination, maybe terminate}{4,0,2}{}{}
\genomerow{2192}{242}{1}{Local Jump If Not}{if !a, goto tag match local anchor}{4,5,4}{}{}
\genomerow{2193}{242}{2}{Local Anchor}{register a local jump-to destination}{5,5,4}{}{}
\genomerow{2194}{242}{3}{Terminal}{a = value}{6,2,5}{}{-3.16524}
\genomerow{2195}{242}{4}{Terminal}{a = value}{1,4,6}{}{-0.948168}
\genomerow{2196}{242}{5}{Send Inter-Cell Message}{if a, send message to neighbor cell}{5,7,7}{}{}
\genomerow{2197}{242}{6}{Terminal}{a = value}{1,2,1}{}{2840.1}
\genomerow{2198}{243}{0}{Global Anchor}{register a global jump-to destination, maybe terminate}{5,3,4}{}{}
\genomerow{2199}{243}{1}{Local Anchor}{register a local jump-to destination}{7,5,0}{}{}
\genomerow{2200}{243}{2}{Terminal}{a = value}{3,0,6}{}{-0.732814}
\genomerow{2201}{243}{3}{Not Equal}{a = (b != c)}{4,5,4}{}{}
\genomerow{2202}{244}{0}{Global Anchor}{register a global jump-to destination, maybe terminate}{2,1,4}{}{}
\genomerow{2203}{244}{1}{Local Anchor}{register a local jump-to destination}{5,2,2}{}{}
\genomerow{2204}{244}{2}{Local Anchor}{register a local jump-to destination}{4,4,4}{}{}
\genomerow{2205}{244}{3}{Set Global Regulator}{set global regulator value to a}{4,6,7}{}{}
\genomerow{2206}{244}{4}{Local Anchor}{register a local jump-to destination}{0,2,6}{}{}
\genomerow{2207}{244}{5}{Write Own State}{if b, state = a}{5,1,1}{ResourceReserveRequest[0]}{}
\genomerow{2208}{244}{6}{Not Equal}{a = (b != c)}{5,5,4}{}{}
\genomerow{2209}{244}{7}{Terminal}{a = value}{5,1,7}{}{0.211613}
\genomerow{2210}{244}{8}{Read Own State}{a = state}{1,1,0}{KinGroupAge[0]}{}
\genomerow{2211}{244}{9}{Terminal}{a = value}{7,1,3}{}{-16.1332}
\genomerow{2212}{244}{10}{Multiply Own State}{state *= a}{3,5,7}{ResourceSendRequest[0]}{}
\genomerow{2213}{244}{11}{Terminal}{a = value}{4,7,0}{}{-0.930367}
\genomerow{2214}{244}{12}{Read Own State}{a = state}{5,4,1}{ResourceSendLimit[0]}{}
\genomerow{2215}{244}{13}{Read Own State}{a = state}{3,4,6}{NeighborKinGroupWillExpire[0]}{}
\genomerow{2216}{244}{14}{Terminal}{a = value}{3,0,1}{}{-1.06671}
\genomerow{2217}{245}{0}{Global Anchor}{register a global jump-to destination, maybe terminate}{2,1,4}{}{}
\genomerow{2218}{245}{1}{Local Anchor}{register a local jump-to destination}{5,2,2}{}{}
\genomerow{2219}{246}{0}{Global Anchor}{register a global jump-to destination, maybe terminate}{3,1,6}{}{}
\genomerow{2220}{246}{1}{Write Own State}{if b, state = a}{5,0,7}{NopState[2]}{}
\genomerow{2221}{246}{2}{Decay Secondary Global Regulator}{decay global regulator value by register a}{7,0,0}{}{}
\genomerow{2222}{246}{3}{Terminal}{a = value}{3,0,2}{}{0.36437}
\genomerow{2223}{246}{4}{Terminal}{a = value}{5,3,4}{}{-0.701068}
\genomerow{2224}{246}{5}{Read Neighbor State}{a = state}{4,5,4}{RicherThanNeighbor[0]}{}
\genomerow{2225}{247}{0}{Global Anchor}{register a global jump-to destination, maybe terminate}{2,3,0}{}{}
\genomerow{2226}{247}{1}{Bitwise Xor}{a = b \textasciicircum{} c}{7,2,0}{}{}
\genomerow{2227}{247}{2}{Logical And}{a = b \&\& c}{7,2,0}{}{}
\genomerow{2228}{247}{3}{Terminal}{a = value}{1,6,3}{}{-31.609}
\genomerow{2229}{247}{4}{Terminal}{a = value}{0,3,0}{}{59625.1}
\genomerow{2230}{247}{5}{Write Own State}{if b, state = a}{0,5,5}{RepLevRequest[1]}{}
\genomerow{2231}{247}{6}{Send Inter-Cell Message}{if a, send message to neighbor cell}{1,4,0}{}{}
\genomerow{2232}{247}{7}{Local Jump If}{if a, goto tag match local anchor}{2,1,4}{}{}
\genomerow{2233}{247}{8}{Modulo}{a = b \% c}{4,4,6}{}{}
\genomerow{2234}{247}{9}{Local Anchor}{register a local jump-to destination}{6,2,0}{}{}
\genomerow{2235}{248}{0}{Global Anchor}{register a global jump-to destination, maybe terminate}{5,6,7}{}{}
\genomerow{2236}{248}{1}{Send Intra-Cell Message}{if a, send message another cardinal within cell}{3,1,0}{}{}
\genomerow{2237}{248}{2}{Multiply Own State}{state *= a}{0,3,1}{NopState[0]}{}
\genomerow{2238}{248}{3}{Read Neighbor State}{a = state}{0,5,4}{NeighborOptimumQuorumExceeded[0]}{}
\genomerow{2239}{248}{4}{Local Jump If}{if a, goto tag match local anchor}{2,1,4}{}{}
\genomerow{2240}{249}{0}{Global Anchor}{register a global jump-to destination, maybe terminate}{4,4,6}{}{}
\genomerow{2241}{250}{0}{Global Anchor}{register a global jump-to destination, maybe terminate}{7,2,5}{}{}
\genomerow{2242}{250}{1}{Send Inter-Cell Message}{if a, send message to neighbor cell}{4,1,4}{}{}
\genomerow{2243}{250}{2}{Read Own State}{a = state}{0,3,3}{NeighborOptimumQuorumExceeded[1]}{}
\genomerow{2244}{250}{3}{Local Anchor}{register a local jump-to destination}{3,2,5}{}{}
\genomerow{2245}{250}{4}{Send Inter-Cell Message}{if a, send message to neighbor cell}{4,1,4}{}{}
\genomerow{2246}{250}{5}{Local Anchor}{register a local jump-to destination}{4,1,5}{}{}
\genomerow{2247}{250}{6}{Send Intra-Cell Message}{if a, send message another cardinal within cell}{0,3,3}{}{}
\genomerow{2248}{250}{7}{Send Intra-Cell Message}{if a, send message another cardinal within cell}{7,2,6}{}{}
\genomerow{2249}{250}{8}{Local Anchor}{register a local jump-to destination}{4,1,5}{}{}
\genomerow{2250}{250}{9}{Local Anchor}{register a local jump-to destination}{6,7,2}{}{}
\genomerow{2251}{250}{10}{Terminal}{a = value}{6,1,5}{}{-1.48755}
\genomerow{2252}{250}{11}{Terminal}{a = value}{1,1,4}{}{-0.998119}
\genomerow{2253}{250}{12}{Set Local Regulator}{set local regulator value to a}{0,3,3}{}{}
\genomerow{2254}{250}{13}{Multiply Own State}{state *= a}{7,2,4}{ResourceReceiveResistance[0]}{}
\genomerow{2255}{250}{14}{Local Anchor}{register a local jump-to destination}{7,1,1}{}{}
\genomerow{2256}{250}{15}{Send Intra-Cell Message}{if a, send message another cardinal within cell}{0,7,2}{}{}
\genomerow{2257}{250}{16}{Terminal}{a = value}{6,3,3}{}{0.896961}
\genomerow{2258}{250}{17}{Get Local Regulator}{a = local regulator value}{5,6,1}{}{}
\genomerow{2259}{250}{18}{Terminal}{a = value}{7,5,0}{}{-6.14391}
\genomerow{2260}{250}{19}{Adjust Local Regulator}{adjust local regulator value by a}{2,7,1}{}{}
\genomerow{2261}{250}{20}{Bitwise Xor}{a = b \textasciicircum{} c}{7,7,6}{}{}
\genomerow{2262}{250}{21}{Local Anchor}{register a local jump-to destination}{0,3,0}{}{}
\genomerow{2263}{250}{22}{Local Jump If}{if a, goto tag match local anchor}{7,0,2}{}{}
\genomerow{2264}{250}{23}{Global Jump If}{if a, goto tag match global anchor; if b, clear registers}{6,4,6}{}{}
\genomerow{2265}{250}{24}{Terminal}{a = value}{2,0,7}{}{0.413024}
\genomerow{2266}{250}{25}{Terminal}{a = value}{7,1,2}{}{-0.453116}
\genomerow{2267}{250}{26}{Terminal}{a = value}{2,7,2}{}{-0.514833}
\genomerow{2268}{251}{0}{Global Anchor}{register a global jump-to destination, maybe terminate}{6,3,3}{}{}
\genomerow{2269}{251}{1}{Local Anchor}{register a local jump-to destination}{2,4,1}{}{}
\genomerow{2270}{251}{2}{Bitwise And}{a = b \& c}{2,6,5}{}{}
\genomerow{2271}{251}{3}{Terminal}{a = value}{6,3,3}{}{0.896961}
\genomerow{2272}{251}{4}{Get Local Regulator}{a = local regulator value}{7,6,3}{}{}
\genomerow{2273}{251}{5}{Bitwise And}{a = b \& c}{2,6,5}{}{}
\genomerow{2274}{251}{6}{Terminal}{a = value}{6,2,3}{}{0.928211}
\genomerow{2275}{251}{7}{Local Anchor}{register a local jump-to destination}{0,4,0}{}{}
\genomerow{2276}{251}{8}{Read Own State}{a = state}{4,1,5}{OptimumQuorumExceeded[1]}{}
\genomerow{2277}{251}{9}{Read Own State}{a = state}{5,7,2}{KinGroupMatch[1]}{}
\genomerow{2278}{251}{10}{Less Than}{a = ( b < c )}{4,4,3}{}{}
\genomerow{2279}{251}{11}{Terminal}{a = value}{7,7,0}{}{0.0545166}
\genomerow{2280}{251}{12}{Send Inter-Cell Message}{if a, send message to neighbor cell}{3,7,5}{}{}
\genomerow{2281}{252}{0}{Global Anchor}{register a global jump-to destination, maybe terminate}{1,1,2}{}{}
\genomerow{2282}{252}{1}{Less Than}{a = ( b < c )}{0,4,3}{}{}
\genomerow{2283}{252}{2}{Local Anchor}{register a local jump-to destination}{2,6,0}{}{}
\genomerow{2284}{252}{3}{Terminal}{a = value}{0,1,2}{}{0.806044}
\genomerow{2285}{252}{4}{Local Anchor}{register a local jump-to destination}{2,3,7}{}{}
\genomerow{2286}{252}{5}{Local Anchor}{register a local jump-to destination}{3,6,5}{}{}
\genomerow{2287}{252}{6}{Fork If}{if a, submit fork request (processed when core terminates)}{5,1,4}{}{}
\genomerow{2288}{252}{7}{Greater Than}{a = ( b > c )}{4,5,3}{}{}
\genomerow{2289}{252}{8}{Terminal}{a = value}{7,7,0}{}{0.0540246}
\genomerow{2290}{252}{9}{Set Local Regulator}{set local regulator value to a}{0,6,3}{}{}
\genomerow{2291}{252}{10}{Terminal}{a = value}{6,4,2}{}{66.9332}
\genomerow{2292}{252}{11}{Get Local Regulator}{a = local regulator value}{7,6,1}{}{}
\genomerow{2293}{252}{12}{Send Inter-Cell Message}{if a, send message to neighbor cell}{2,5,0}{}{}
\genomerow{2294}{252}{13}{Terminal}{a = value}{2,4,0}{}{-2.56818}
\genomerow{2295}{252}{14}{Local Anchor}{register a local jump-to destination}{6,4,4}{}{}
\genomerow{2296}{252}{15}{Local Anchor}{register a local jump-to destination}{6,0,6}{}{}
\genomerow{2297}{252}{16}{Send Inter-Cell Message}{if a, send message to neighbor cell}{3,7,5}{}{}
\genomerow{2298}{252}{17}{Logical Or}{a = b \textbar{}\textbar{} c}{2,5,6}{}{}
\genomerow{2299}{252}{18}{Not Equal}{a = (b != c)}{1,1,5}{}{}
\genomerow{2300}{252}{19}{Terminal}{a = value}{1,5,2}{}{0.874001}
\genomerow{2301}{252}{20}{Local Anchor}{register a local jump-to destination}{2,2,3}{}{}
\genomerow{2302}{252}{21}{Nop-0}{perform no operation}{2,1,5}{}{}
\genomerow{2303}{252}{22}{Local Anchor}{register a local jump-to destination}{2,2,3}{}{}
\genomerow{2304}{252}{23}{Equal}{a = ( b == c )}{0,1,5}{}{}
\genomerow{2305}{252}{24}{Local Anchor}{register a local jump-to destination}{0,3,3}{}{}
\genomerow{2306}{252}{25}{Read Own State}{a = state}{4,1,7}{KinGroupWillExpire[0]}{}
\genomerow{2307}{253}{0}{Global Anchor}{register a global jump-to destination, maybe terminate}{1,3,6}{}{}
\genomerow{2308}{253}{1}{Less Than}{a = ( b < c )}{5,4,1}{}{}
\genomerow{2309}{253}{2}{Terminal}{a = value}{1,5,2}{}{0.874001}
\genomerow{2310}{253}{3}{Nop-0}{perform no operation}{2,1,5}{}{}
\genomerow{2311}{253}{4}{Local Anchor}{register a local jump-to destination}{2,2,3}{}{}
\genomerow{2312}{253}{5}{Equal}{a = ( b == c )}{0,1,5}{}{}
\genomerow{2313}{253}{6}{Local Anchor}{register a local jump-to destination}{0,3,3}{}{}
\genomerow{2314}{253}{7}{Read Own State}{a = state}{4,1,7}{KinGroupWillExpire[0]}{}
\genomerow{2315}{254}{0}{Global Anchor}{register a global jump-to destination, maybe terminate}{1,3,6}{}{}
\genomerow{2316}{254}{1}{Less Than}{a = ( b < c )}{5,4,1}{}{}
\genomerow{2317}{254}{2}{Terminal}{a = value}{1,5,2}{}{0.874001}
\genomerow{2318}{254}{3}{Bitwise And}{a = b \& c}{1,0,5}{}{}
\genomerow{2319}{254}{4}{Local Anchor}{register a local jump-to destination}{2,6,7}{}{}
\genomerow{2320}{254}{5}{Terminal}{a = value}{6,3,2}{}{0.709214}
\genomerow{2321}{254}{6}{Terminal}{a = value}{3,0,3}{}{1.27752}
\genomerow{2322}{254}{7}{Send Intra-Cell Message}{if a, send message another cardinal within cell}{4,1,7}{}{}
\genomerow{2323}{254}{8}{Bitwise And}{a = b \& c}{1,0,5}{}{}
\genomerow{2324}{254}{9}{Local Anchor}{register a local jump-to destination}{2,6,7}{}{}
\genomerow{2325}{254}{10}{Terminal}{a = value}{6,3,2}{}{0.709214}
\genomerow{2326}{254}{11}{Terminal}{a = value}{3,0,3}{}{1.27752}
\genomerow{2327}{254}{12}{Send Intra-Cell Message}{if a, send message another cardinal within cell}{4,1,7}{}{}
\genomerow{2328}{254}{13}{Adjust Global Regulator}{adjust global regulator value by a}{5,5,6}{}{}
\genomerow{2329}{254}{14}{Bitwise Not}{a = \textasciitilde{}b}{3,7,2}{}{}
\genomerow{2330}{254}{15}{Read Own State}{a = state}{4,1,7}{IsChildGroupOf[0]}{}
\genomerow{2331}{255}{0}{Global Anchor}{register a global jump-to destination, maybe terminate}{1,2,6}{}{}
\genomerow{2332}{255}{1}{Less Than}{a = ( b < c )}{5,4,1}{}{}
\genomerow{2333}{255}{2}{Terminal}{a = value}{1,5,3}{}{-0.141623}
\genomerow{2334}{255}{3}{Bitwise And}{a = b \& c}{1,2,5}{}{}
\genomerow{2335}{255}{4}{Read Own State}{a = state}{6,3,2}{KinGroupMatch[1]}{}
\genomerow{2336}{255}{5}{Read Own State}{a = state}{7,2,7}{StockpileDepleted[0]}{}
\genomerow{2337}{256}{0}{Global Anchor}{register a global jump-to destination, maybe terminate}{6,7,2}{}{}
\genomerow{2338}{256}{1}{Divide}{a = b / c}{7,0,3}{}{}
\genomerow{2339}{256}{2}{Get Global Regulator}{a = global regulator value}{5,7,0}{}{}
\genomerow{2340}{256}{3}{Terminal}{a = value}{7,0,3}{}{-2.16061}
\genomerow{2341}{256}{4}{Terminal}{a = value}{2,6,1}{}{1.93478}
\genomerow{2342}{256}{5}{Multiply Own State}{state *= a}{0,5,0}{TransientNopState[1]}{}
\genomerow{2343}{256}{6}{Nop-2}{perform no operation}{5,6,3}{}{}
\genomerow{2344}{256}{7}{RandomFill}{fill a with uniformly-distributed random bits}{1,1,4}{}{}
\genomerow{2345}{256}{8}{Terminal}{a = value}{0,0,2}{}{0.192858}
\genomerow{2346}{256}{9}{Get Secondary Global Regulator}{a = global regulator value}{6,0,0}{}{}
\genomerow{2347}{256}{10}{Decay Secondary Global Regulator}{decay global regulator value by register a}{6,0,0}{}{}
\genomerow{2348}{256}{11}{Terminal}{a = value}{7,2,3}{}{-0.716824}
\genomerow{2349}{256}{12}{Local Anchor}{register a local jump-to destination}{4,0,4}{}{}
\genomerow{2350}{256}{13}{Write Own State}{if b, state = a}{4,0,4}{TransientNopState[2]}{}
\genomerow{2351}{256}{14}{Get Global Regulator}{a = global regulator value}{3,1,6}{}{}
\genomerow{2352}{256}{15}{Decay Secondary Global Regulator}{decay global regulator value by register a}{6,0,0}{}{}
\genomerow{2353}{256}{16}{Terminal}{a = value}{7,2,3}{}{-0.716824}
\genomerow{2354}{257}{0}{Global Anchor}{register a global jump-to destination, maybe terminate}{4,0,4}{}{}
\genomerow{2355}{257}{1}{Send Intra-Cell Message}{if a, send message another cardinal within cell}{4,0,4}{}{}
\genomerow{2356}{257}{2}{Terminate If}{if a, terminate core}{6,4,2}{}{}
%
\end{longtable}
\end{scriptsize}
\end{landscape}
